% V : produire un rapport — Français Tle Tech — Cours signé OnBuch+
% Programme officiel MINESEC. Compiler avec : tectonic N17-v-produire-rapport.tex
\newcommand{\DOCMATIERE}{Français}
\newcommand{\DOCNIVEAU}{Tle Tech}
\newcommand{\DOCMODULE}{Français – classe de Terminale technique}
\newcommand{\DOCLECON}{N17}
\newcommand{\DOCTITRE}{V : produire un rapport}
% OnBuch+ — préambule « cours signés » ENRICHI (compilé avec tectonic / XeLaTeX)
% Système visuel « Pop » d'OnBuch+ : fond crème (jamais blanc), bordures encre
% épaisses, ombres « dures » décalées, titres Archivo Black en capitales.
% Version MATHÉMATIQUES : graphes (pgfplots/tikz), boîtes pédagogiques
% (définition, propriété, méthode, attention, le savais-tu, exercice résolu,
% exercice, TP, bilan), page de garde et signature OnBuch+ ancrée en bas.
%
% Macros attendues dans main.tex AVANT \input{preamble} :
%   \DOCMATIERE  \DOCNIVEAU  \DOCMODULE  \DOCLECON (numéro)  \DOCTITRE
\documentclass[11pt]{article}

\usepackage{fontspec}
\usepackage[french]{babel}
\usepackage[a4paper,margin=2cm,top=3.4cm,bottom=2.8cm,headheight=1.4cm,footskip=1.3cm]{geometry}
\usepackage{amsmath,amssymb,amsfonts}
\usepackage{mathtools} % \xmapsto, \coloneqq... souvent utilisés par les modèles
\usepackage{cancel} % simplifications barrées (bilans de forces)
% \abs{z} (module, valeur absolue) et \norm{u} (norme), utilisés par les modèles
\providecommand{\abs}{}\let\abs\relax\DeclarePairedDelimiter{\abs}{\lvert}{\rvert}
\providecommand{\norm}{}\let\norm\relax\DeclarePairedDelimiter{\norm}{\lVert}{\rVert}
\usepackage[table]{xcolor} % [table] : fournit aussi \rowcolors (alternance de lignes)
\usepackage{pagecolor}
\usepackage{tikz}
\usetikzlibrary{arrows.meta,positioning,calc,shapes.geometric,shapes.misc,shapes.arrows,decorations.pathmorphing,decorations.pathreplacing,decorations.markings,patterns,shadows,mindmap,trees,fit,backgrounds,matrix,intersections,angles,quotes}
\usepackage{pgfplots}
\usepackage[version=4]{mhchem} % équations nucléaires \ce{^{235}_{92}U}
\usepackage[european,siunitx]{circuitikz} % schémas électriques
\pgfplotsset{compat=1.18}
\usepgfplotslibrary{fillbetween,groupplots} % aires entre courbes (calcul intégral)
\usepackage{enumitem}
\usepackage{fancyhdr}
% [most] charge toutes les bibliothèques tcolorbox (listings, minted, raster,
% documentation...) et gonfle inutilement la mémoire TeX sur un document long
% (~300 boîtes) : on ne charge que ce qui est réellement utilisé.
\usepackage[skins,breakable]{tcolorbox}
\usepackage{graphicx}
\usepackage{colortbl}
\usepackage{booktabs}
\usepackage{tabularx}
\usepackage{multirow}
\usepackage{diagbox} % case d'en-tête diagonale des tableaux à double entrée
% Colonnes extensibles centrée / gauche / droite (C, L, R), souvent utilisées par les modèles
\newcolumntype{C}{>{\centering\arraybackslash}X}
\newcolumntype{L}{>{\raggedright\arraybackslash}X}
\newcolumntype{R}{>{\raggedleft\arraybackslash}X}
\usepackage{array}
\usepackage{multicol}
\usepackage{siunitx}
% parse-numbers=false : accepte aussi \SI{2,5 \times 10^{-3}}{...}
\sisetup{locale=FR,output-decimal-marker={,},group-minimum-digits=4,parse-numbers=false}
\DeclareSIUnit\ohm{\ensuremath{\symup{\Omega}}} % U+2126 absent de la police
\DeclareSIUnit\division{div} % réglages d'oscilloscope (ms/div, V/div)
\DeclareSIUnit\tr{tr} % tours (vitesse de rotation, ex. \SI{1500}{\tr\per\minute})
\DeclareSIUnit\molar{M} % concentration molaire (ex. \SI{3}{\molar}, \SI{1}{\micro\molar})
\DeclareSIUnit\an{an} % année (le modèle confond parfois avec \year, un compteur TeX) — ex. \SI{5}{\kilo\meter\per\an}
\DeclareSIUnit\FCFA{FCFA} % franc CFA (ex. \SI{500}{\FCFA\per\kilo\gram})
\DeclareSIUnit\degreeBrix{\textdegree Brix} % teneur en sucre d'un sirop/jus (réfractométrie)
\DeclareSIUnit\month{mois} % le modèle utilise parfois \month, un compteur TeX, comme unité de durée
\DeclareSIUnit\microliter{\micro\liter} % le modèle écrit parfois \microliter en un seul token
\DeclareSIUnit\million{millions} % dénombrement (ex. \SI{15}{\million\per\mL} spermatozoïdes)
\usepackage{qrcode}
% \degree : commande fréquemment utilisée par les modèles pour un angle ou une
% température en dehors de \SI{}{} (non fournie par les paquets chargés).
\providecommand{\degree}{^{\circ}}
% \qed (fin de démonstration), utilisé par les modèles sans amsthm
\providecommand{\qed}{\hfill$\square$}
% \barre : double barre des valeurs interdites dans un tableau de variations (tkz-tab)
\providecommand{\barre}{\ensuremath{\|}}
% environnement proof (amsthm non chargé) utilisé par les modèles
\ifdefined\proof\else\newenvironment{proof}[1][Démonstration]{\par\noindent\textit{#1.}\ }{\qed\par}\fi
\usepackage{lastpage}
\usepackage{hyperref}
\hypersetup{hidelinks}

% ============================================================
% Palette « Pop » OnBuch+ — mêmes hex que la classe P (plus_ui.dart)
% ============================================================
\definecolor{poporange}{HTML}{F59321}
\definecolor{popdark}{HTML}{E07A0C}
\definecolor{popcream}{HTML}{FFF6E8}
\definecolor{popink}{HTML}{1C1714}
\definecolor{poppurple}{HTML}{7A5AE0}
\definecolor{popgreen}{HTML}{2E9E6A}
\definecolor{poppink}{HTML}{E0457B}
\definecolor{popblue}{HTML}{2F7DE1}
\definecolor{popgold}{HTML}{C98A12}
\definecolor{popmuted}{HTML}{8A7F76}
\definecolor{popline}{HTML}{EADFD2}
\definecolor{popgray}{HTML}{8A7F76}
\colorlet{popgrey}{popgray}
% Alias tolérants (le modèle nomme parfois les couleurs différemment)
\colorlet{popwhite}{white}
\colorlet{popblack}{popink}
\colorlet{popred}{poppink}
\colorlet{popyellow}{popgold}
% Teintes claires (fonds de tableaux/encadrés) — le modèle nomme parfois
% les couleurs avec un suffixe L (ex. popblueL) sans les avoir définies.
\colorlet{poporangeL}{poporange!20}
\colorlet{popdarkL}{popdark!20}
\colorlet{poppurpleL}{poppurple!20}
\colorlet{popgreenL}{popgreen!20}
\colorlet{poppinkL}{poppink!20}
\colorlet{popblueL}{popblue!20}
\colorlet{popgoldL}{popgold!20}
\colorlet{popredL}{popred!20}
\colorlet{popgrayL}{popgray!20}
% Teintes claires pour fonds de tableaux / zones
\colorlet{popblueL}{popblue!10!white}
\colorlet{poporangeL}{poporange!14!white}
\colorlet{popgreenL}{popgreen!12!white}
\colorlet{poppurpleL}{poppurple!12!white}
\colorlet{poppinkL}{poppink!12!white}
\colorlet{popgoldL}{popgold!14!white}

\pagecolor{popcream}

% ============================================================
% Typographie
% ============================================================
\defaultfontfeatures{Ligatures=TeX}
\setmainfont{PlusJakartaSans-Medium.ttf}[Path=../../../fonts/,BoldFont=PlusJakartaSans-Bold.ttf]
% Maths en sans-serif (Fira Math) avec chiffres et lettres droites de Plus
% Jakarta, pour que les formules se fondent dans le texte.
\usepackage[math-style=ISO,bold-style=ISO]{unicode-math}
\setmathfont{FiraMath-Regular.otf}[Path=../../../fonts/]
\setmathfont{PlusJakartaSans-Medium.ttf}[Path=../../../fonts/,range={up/num,up/{Latin,latin},\mathup}]
\usepackage{icomma} % virgule décimale sans espace en mode math
% Caractères mathématiques Unicode absents des polices de texte (Plus Jakarta,
% Archivo Black) : rendus par leur équivalent mathématique, en texte comme en
% formule — sinon ils s'affichent en carré vide (ex. « Arithmétique dans ℤ »).
\usepackage{newunicodechar}
\newunicodechar{ℝ}{\ensuremath{\mathbb{R}}}
\newunicodechar{ℤ}{\ensuremath{\mathbb{Z}}}
\newunicodechar{ℂ}{\ensuremath{\mathbb{C}}}
\newunicodechar{ℕ}{\ensuremath{\mathbb{N}}}
\newunicodechar{ℚ}{\ensuremath{\mathbb{Q}}}
\newunicodechar{𝔻}{\ensuremath{\mathbb{D}}}
\newunicodechar{α}{\ensuremath{\alpha}}
\newunicodechar{β}{\ensuremath{\beta}}
\newunicodechar{γ}{\ensuremath{\gamma}}
\newunicodechar{δ}{\ensuremath{\delta}}
\newunicodechar{ε}{\ensuremath{\varepsilon}}
\newunicodechar{θ}{\ensuremath{\theta}}
\newunicodechar{λ}{\ensuremath{\lambda}}
\newunicodechar{μ}{\ensuremath{\mu}}
\newunicodechar{σ}{\ensuremath{\sigma}}
\newunicodechar{τ}{\ensuremath{\tau}}
\newunicodechar{φ}{\ensuremath{\varphi}}
\newunicodechar{ψ}{\ensuremath{\psi}}
\newunicodechar{ω}{\ensuremath{\omega}}
\newunicodechar{Σ}{\ensuremath{\Sigma}}
% majuscules produites par \MakeUppercase dans les titres (α-aminés -> Α)
\newunicodechar{Α}{\ensuremath{\alpha}}
\newunicodechar{Β}{\ensuremath{\beta}}
\newunicodechar{Γ}{\ensuremath{\Gamma}}
\newunicodechar{Δ}{\ensuremath{\Delta}}
\newunicodechar{Ε}{\ensuremath{\varepsilon}}
\newunicodechar{Θ}{\ensuremath{\Theta}}
\newunicodechar{Λ}{\ensuremath{\Lambda}}
\newunicodechar{Μ}{\ensuremath{\mu}}
\newunicodechar{Τ}{\ensuremath{\tau}}
\newunicodechar{Φ}{\ensuremath{\Phi}}
\newunicodechar{Ψ}{\ensuremath{\Psi}}
\newunicodechar{Ω}{\ensuremath{\Omega}}
\newunicodechar{↦}{\ensuremath{\mapsto}}
\newunicodechar{∈}{\ensuremath{\in}}
\newunicodechar{∉}{\ensuremath{\notin}}
\newunicodechar{∧}{\ensuremath{\wedge}}
\newunicodechar{⇔}{\ensuremath{\Leftrightarrow}}
\newunicodechar{⟺}{\ensuremath{\Longleftrightarrow}}
\newunicodechar{⇒}{\ensuremath{\Rightarrow}}
\newunicodechar{⟹}{\ensuremath{\Longrightarrow}}
\newunicodechar{∘}{\ensuremath{\circ}}
\newunicodechar{∪}{\ensuremath{\cup}}
\newunicodechar{∩}{\ensuremath{\cap}}
\newunicodechar{≡}{\ensuremath{\equiv}}
\newunicodechar{⊂}{\ensuremath{\subset}}
\newunicodechar{⊥}{\ensuremath{\perp}}
\newunicodechar{∃}{\ensuremath{\exists}}
\newunicodechar{∀}{\ensuremath{\forall}}
\newunicodechar{∛}{\ensuremath{\sqrt[3]{\,}}}
\newunicodechar{∜}{\ensuremath{\sqrt[4]{\,}}}
\newunicodechar{⃗}{\ensuremath{{}^{\to}}}
\newunicodechar{ⁿ}{\ifmmode^{n}\else\textsuperscript{\ensuremath{n}}\fi}
\newunicodechar{ˣ}{\ifmmode^{x}\else\textsuperscript{\ensuremath{x}}\fi}
\newunicodechar{ᵇ}{\ifmmode^{b}\else\textsuperscript{\ensuremath{b}}\fi}
\newunicodechar{ʳ}{\ifmmode^{r}\else\textsuperscript{\ensuremath{r}}\fi}
\newunicodechar{ᵘ}{\ifmmode^{u}\else\textsuperscript{\ensuremath{u}}\fi}
\newunicodechar{ʸ}{\ifmmode^{y}\else\textsuperscript{\ensuremath{y}}\fi}
\newunicodechar{ᵃ}{\ifmmode^{a}\else\textsuperscript{\ensuremath{a}}\fi}
\newunicodechar{ᵉ}{\ifmmode^{e}\else\textsuperscript{\ensuremath{e}}\fi}
\newunicodechar{ᵏ}{\ifmmode^{k}\else\textsuperscript{\ensuremath{k}}\fi}
\newunicodechar{ᵖ}{\ifmmode^{p}\else\textsuperscript{\ensuremath{p}}\fi}
\newunicodechar{ᴺ}{\ifmmode^{N}\else\textsuperscript{\ensuremath{N}}\fi}
\newunicodechar{ᶜ}{\ifmmode^{c}\else\textsuperscript{\ensuremath{c}}\fi}
\newunicodechar{ʰ}{\ifmmode^{h}\else\textsuperscript{\ensuremath{h}}\fi}
\newunicodechar{ᵠ}{\ifmmode^{q}\else\textsuperscript{\ensuremath{q}}\fi}
\newunicodechar{ₙ}{\ifmmode_{n}\else\textsubscript{\ensuremath{n}}\fi}
\newunicodechar{ₐ}{\ifmmode_{a}\else\textsubscript{\ensuremath{a}}\fi}
\newunicodechar{ᵢ}{\ifmmode_{i}\else\textsubscript{\ensuremath{i}}\fi}
\newunicodechar{ᵣ}{\ifmmode_{r}\else\textsubscript{\ensuremath{r}}\fi}
\newunicodechar{ₖ}{\ifmmode_{k}\else\textsubscript{\ensuremath{k}}\fi}
\newunicodechar{ᵦ}{\ifmmode_{\beta}\else\textsubscript{\ensuremath{\beta}}\fi}
\newunicodechar{ₓ}{\ifmmode_{x}\else\textsubscript{\ensuremath{x}}\fi}
\newfontfamily\popdisplay{ArchivoBlack-Regular.ttf}[Path=../../../fonts/]
\newfontfamily\poplogo{SpaceGrotesk-Bold.ttf}[Path=../../../fonts/]
\newfontfamily\popsemi{PlusJakartaSans-SemiBold.ttf}[Path=../../../fonts/]
\newfontfamily\popxbold{PlusJakartaSans-ExtraBold.ttf}[Path=../../../fonts/]
\newfontfamily\popxboldls{PlusJakartaSans-ExtraBold.ttf}[Path=../../../fonts/,LetterSpace=3.0]

\setlist[enumerate]{leftmargin=1.7em,itemsep=5pt,topsep=4pt}
\setlist[itemize]{leftmargin=1.7em,itemsep=4pt,topsep=4pt,label={\color{poporange}\textbullet}}
\setlength{\parindent}{0pt}
\setlength{\parskip}{5pt}
\renewcommand{\arraystretch}{1.25}

% pgfplots : style Pop par défaut
\pgfplotsset{
  every axis/.append style={axis line style={popink,line width=1pt},tick style={popink},
    grid style={popline,line width=0.5pt},label style={font=\small},tick label style={font=\footnotesize}},
  popcurve/.style={poporange,line width=1.8pt,no markers,smooth},
}
% Même style disponible dans un \draw TikZ simple (hors pgfplots)
\tikzset{popcurve/.style={poporange,line width=1.8pt}}
\tikzset{popgrid/.style={popline,very thin}}
\colorlet{popcurve}{poporange} % le nom du style est parfois utilisé comme couleur

% ============================================================
% Pill / squiggle
% ============================================================
\newcommand{\poppill}[3]{%
  \tikz[baseline=-0.55ex]\node[draw=popink,line width=1.2pt,rounded corners=2.6pt,%
    fill=#2,inner xsep=6.5pt,inner ysep=3pt]{%
    \popxboldls\fontsize{7}{7}\selectfont\color{#3}\MakeUppercase{#1}};%
}
\newcommand{\popsquiggle}[1]{%
  \tikz[baseline=-0.4ex]{
    \draw[#1,line width=2.6pt,line cap=round]
      plot [smooth] coordinates {(0,0.09) (0.34,0.01) (0.68,0.21) (1.02,0.01) (1.36,0.10)};
  }%
}
% Mot-clé surligné (marqueur orange sous le texte)
\newcommand{\cle}[1]{{\popxbold\color{popdark}#1}}

% ============================================================
% En-tête / pied
% ============================================================
\pagestyle{fancy}
\fancyhf{}
\renewcommand{\headrulewidth}{0pt}
\renewcommand{\footrulewidth}{0pt}
\lhead{\raisebox{-0.15ex}{\poplogo\fontsize{13}{13}\selectfont\color{popink}OnBuch\color{poporange}+}}
\rhead{\poppill{\DOCMATIERE\ \textbullet\ \DOCNIVEAU\ \textbullet\ Leçon \DOCLECON}{white}{popink}}
\lfoot{\popxbold\fontsize{7.6}{7.6}\selectfont\color{popmuted}\MakeUppercase{Cours signé OnBuch+}\ \textbullet\ {\popsemi\DOCTITRE}}
\rfoot{\tikz[baseline=-0.6ex]\node[rounded corners=8.5pt,fill=popink,minimum height=17pt,minimum width=17pt,inner xsep=5pt,inner sep=0]{\popxbold\fontsize{8}{8}\selectfont\color{popcream}\thepage\,/\,\pageref*{LastPage}};}

% ============================================================
% Titres
% ============================================================
\newcommand{\chaptitre}[2]{%
  \begin{tcolorbox}[enhanced,colback=poporange,colframe=popink,boxrule=2.6pt,arc=11pt,
      drop shadow=popink,left=15pt,right=15pt,top=12pt,bottom=11pt,before skip=3pt,after skip=18pt]
    \poppill{#2}{popink}{popcream}\par\vspace{9pt}
    {\raggedright\hyphenpenalty=10000\exhyphenpenalty=10000\popdisplay\fontsize{22}{24}\selectfont\color{popcream}\MakeUppercase{#1}\par}
  \end{tcolorbox}
}
% Section numérotée : pastille orange avec le numéro + titre Archivo
\newcounter{popsec}
\newcommand{\coursec}[1]{%
  \stepcounter{popsec}\setcounter{popsub}{0}%
  \par\vspace{16pt}\noindent
  \begin{minipage}{\linewidth}
  \tikz[baseline=-0.7ex]\node[draw=popink,line width=1.6pt,fill=poporange,rounded corners=4pt,minimum size=22pt,inner sep=0]{\popdisplay\fontsize{12}{12}\selectfont\color{popcream}\thepopsec};%
  \hspace{8pt}{\popdisplay\fontsize{15}{17}\selectfont\color{popink}\MakeUppercase{#1}}\par
  \vspace{4pt}\noindent{\color{poporange}\rule{36pt}{3.6pt}}
  \end{minipage}\par\nopagebreak\vspace{8pt}
}
\newcounter{popsub}
\newcommand{\courssub}[1]{%
  \stepcounter{popsub}%
  \par\vspace{9pt}\noindent{\popxbold\fontsize{11.5}{13}\selectfont\color{popink}{\color{poporange}\thepopsec.\thepopsub}\hspace{6pt}#1}\par\nopagebreak\vspace{4pt}
}

% ============================================================
% Boîtes pédagogiques — même recette : fond blanc, bordure épaisse colorée,
% ombre dure encre, badge plein. Titre optionnel [#1].
% ============================================================
\tcbset{popbase/.style={breakable,enhanced,colback=white,boxrule=2.3pt,arc=9pt,
  shadow={3pt}{-3pt}{0pt}{popink},left=14pt,right=14pt,top=11pt,bottom=11pt,before skip=11pt,after skip=15pt,
  fontupper=\popsemi\fontsize{10.6}{14.5}\selectfont\color{popink},
  fontlower=\popsemi\fontsize{10.6}{14.5}\selectfont\color{popink}}}
% Coche dessinée (le glyphe ✓ n'existe pas dans les polices du document)
\newcommand{\popcheck}[1]{\tikz[baseline=-0.5ex]\draw[#1,line width=1.6pt,line cap=round,line join=round](-0.13,0.0)--(-0.04,-0.09)--(0.13,0.1);}
\providecommand{\coche}{\popcheck{popgreen}}
\providecommand{\resultat}[1]{\par\smallskip\noindent\fcolorbox{popgreen}{popgreenL}{\parbox{\dimexpr\linewidth-2\fboxsep-2\fboxrule\relax}{\textbf{Résultat.} #1}}\par\smallskip}
\newcommand{\popboxhead}[3]{\poppill{#1}{#2}{white}\ifx\relax#3\relax\else\hspace{6pt}{\popxbold\fontsize{10.6}{12}\selectfont\color{popink}#3}\fi\par\vspace{7pt}}

\newtcolorbox{aretenir}[1][]{popbase,colframe=popgreen,before upper={\popboxhead{À retenir}{popgreen}{#1}}}
\newtcolorbox{exemplebox}[1][]{popbase,colframe=poporange,before upper={\popboxhead{Exemple}{poporange}{#1}}}
\newtcolorbox{definition}[1][]{popbase,colframe=popblue,before upper={\popboxhead{Définition}{popblue}{#1}}}
\newtcolorbox{propriete}[1][]{popbase,colframe=poppurple,before upper={\popboxhead{Propriété}{poppurple}{#1}}}
\newtcolorbox{methode}[1][]{popbase,colframe=popgold,before upper={\popboxhead{Méthode}{popgold}{#1}}}
\newtcolorbox{attention}[1][]{popbase,colframe=poppink,before upper={\popboxhead{Attention}{poppink}{#1}}}
\newtcolorbox{experience}[1][]{popbase,colframe=popblue,colback=popblueL,before upper={\popboxhead{Expérience}{popblue}{#1}}}
% « Le savais-tu ? » : carte violette pleine (contexte, culture, Cameroun)
\newtcolorbox{savaistu}[1][]{popbase,colback=poppurple,colframe=popink,
  fontupper=\popsemi\fontsize{10.4}{14.2}\selectfont\color{white},
  before upper={\poppill{Le savais-tu\,?}{white}{poppurple}\ifx\relax#1\relax\else\hspace{6pt}{\popxbold\color{white}#1}\fi\par\vspace{7pt}}}
% Exercice résolu : énoncé en haut, \tcblower puis la solution sur fond crème
\newcounter{popexo}\newcounter{popexr}
\newtcolorbox{exoresolu}[1][]{popbase,colframe=poporange,colbacklower=popcream,
  segmentation style={popink,line width=1.2pt,dashed},
  before upper={\stepcounter{popexr}\popboxhead{Exercice résolu \thepopexr}{poporange}{#1}},
  before lower={\poppill{Solution}{popink}{popcream}\par\vspace{6pt}}}
% Exercice (non résolu en place) — la correction est dans la partie Corrigés
\newtcolorbox{exercice}[1][]{popbase,colframe=popink,
  before upper={\stepcounter{popexo}\popboxhead{Exercice \thepopexo}{popink}{#1}}}
% Corrigé
\newtcolorbox{corrige}[1][]{popbase,colframe=popgreen,colback=popgreenL,
  before upper={\popboxhead{Corrigé}{popgreen}{#1}}}
% Objectifs / prérequis / bilan
\newtcolorbox{objectifs}{popbase,colframe=popink,colback=poporangeL,
  before upper={\popboxhead{Objectifs de la leçon}{popink}{}}}
\newtcolorbox{prerequis}{popbase,colframe=popink,colback=popblueL,
  before upper={\popboxhead{Prérequis}{popblue}{}}}
\newtcolorbox{bilan}{popbase,colframe=popink,colback=white,boxrule=2.8pt,
  before upper={\popboxhead{Fiche bilan}{poporange}{L'essentiel à connaître pour l'examen}}}
% Figure encadrée avec légende
\newtcolorbox{popfigure}[1][]{enhanced,breakable=false,colback=white,colframe=popline,boxrule=1.4pt,arc=8pt,
  left=10pt,right=10pt,top=10pt,bottom=8pt,before skip=10pt,after skip=12pt,halign=center,
  lower separated=false,halign lower=center,
  fontlower=\popsemi\fontsize{9}{11.5}\selectfont\color{popmuted},
  #1}
\newcommand{\legende}[1]{\par\vspace{6pt}{\popsemi\fontsize{9}{11.5}\selectfont\color{popmuted}#1}}

% ============================================================
% Signature OnBuch+
% ============================================================
\newcommand{\poplogoinline}[1]{{\poplogo\fontsize{#1}{#1}\selectfont\color{popink}OnBuch\color{poporange}+}}
\newcommand{\signatureonbuch}{%
  \begin{tcolorbox}[enhanced,colback=popink,colframe=popink,boxrule=2.2pt,arc=10pt,
      drop shadow=poporange,left=14pt,right=14pt,top=10pt,bottom=10pt,before skip=0pt,after skip=0pt]
    \tikz[baseline=-0.7ex]\node[circle,fill=popgreen,draw=popcream,line width=1.3pt,minimum size=22pt,inner sep=0]{\popcheck{white}};%
    \hspace{9pt}\begin{minipage}[c]{0.62\linewidth}
      {\popxbold\fontsize{12}{13}\selectfont\color{popcream}Signé\ \textbullet\ {\poplogo OnBuch\color{poporange}+}}\par\vspace{2pt}
      {\raggedright\popsemi\fontsize{8}{10.5}\selectfont\color{popline}\MakeUppercase{Équipe pédagogique OnBuch+}\\\MakeUppercase{Conforme au programme officiel MINESEC}\par}
    \end{minipage}\hfill
    {\poplogo\fontsize{22}{22}\selectfont\color{popcream}OnBuch\color{poporange}+}
  \end{tcolorbox}%
}

% ============================================================
% Page de garde
% #1 titre · #2 matière · #3 niveau · #4 module · #5 n° leçon · #6 durée ·
% #7 description / objectifs (texte ou itemize)
% ============================================================
\newcommand{\pagegarde}[7]{%
  \thispagestyle{empty}
  \newgeometry{a4paper,margin=1.7cm,top=1.6cm,bottom=1.4cm}%
  \begin{tikzpicture}[remember picture,overlay]
    \fill[poporange!18] ([xshift=-3.2cm,yshift=-3.6cm]current page.north east) circle (4.6cm);
    \fill[poppurple!14] ([xshift=2.4cm,yshift=7.6cm]current page.south west) circle (3.8cm);
  \end{tikzpicture}%
  {\poplogo\fontsize{30}{30}\selectfont\color{popink}OnBuch\color{poporange}+}\hfill
  \raisebox{0.6ex}{\poppill{Cours signé \textbullet\ Premium}{popink}{popcream}}\par
  \vspace{1pt}\popsquiggle{poppurple}\par
  \vspace{8pt}
  \begin{tcolorbox}[enhanced,colback=poporange,colframe=popink,boxrule=2.8pt,arc=13pt,
      drop shadow=popink,left=18pt,right=18pt,top=12pt,bottom=13pt,after skip=10pt]
    \poppill{#2 \textbullet\ #3}{popink}{popcream}\hspace{5pt}\poppill{#4}{white}{popink}\par\vspace{8pt}
    {\popdisplay\fontsize{14}{14}\selectfont\color{popink}LEÇON #5}\par\vspace{4pt}
    {\raggedright\hyphenpenalty=10000\exhyphenpenalty=10000\popdisplay\fontsize{25}{27}\selectfont\color{popcream}\MakeUppercase{#1}\par}
    \vspace{8pt}
    {\popxbold\fontsize{9.5}{9.5}\selectfont\color{popink}\MakeUppercase{Durée conseillée : #6}}
  \end{tcolorbox}
  \begin{tcolorbox}[enhanced,colback=white,colframe=popink,boxrule=2.2pt,arc=10pt,
      drop shadow=popink,left=16pt,right=16pt,top=10pt,bottom=10pt,after skip=8pt]
    \poppill{Dans cette leçon}{poppurple}{white}\par\vspace{5pt}
    \popsemi\fontsize{9.3}{12.6}\selectfont\color{popink}#7
  \end{tcolorbox}
  \noindent\begin{minipage}[t]{0.487\linewidth}
    \begin{tcolorbox}[enhanced,colback=popgreen,colframe=popink,boxrule=2.2pt,arc=10pt,
        drop shadow=popink,left=11pt,right=11pt,top=10pt,bottom=10pt,height=3.1cm,valign=center]
      \tcbox[colback=white,colframe=popink,boxrule=1.3pt,arc=5pt,boxsep=0pt,nobeforeafter,
        left=3pt,right=3pt,top=3pt,bottom=3pt,valign=center]{\qrcode[height=1.9cm]{https://play.google.com/store/apps/details?id=com.luvvix.onbuch&hl=fr}}%
      \hspace{8pt}\begin{minipage}[c]{0.5\linewidth}
        {\popdisplay\fontsize{11}{12.5}\selectfont\color{white}\MakeUppercase{Télécharge OnBuch}}\par\vspace{4pt}
        {\raggedright\popsemi\fontsize{8.4}{11}\selectfont\color{white}Gratuit \textbullet\ Google Play\\Tuteur IA, annales, concours, résultats\par}
      \end{minipage}
    \end{tcolorbox}
  \end{minipage}\hfill
  \begin{minipage}[t]{0.487\linewidth}
    \begin{tcolorbox}[enhanced,colback=poppurple,colframe=popink,boxrule=2.2pt,arc=10pt,
        drop shadow=popink,left=11pt,right=11pt,top=10pt,bottom=10pt,height=3.1cm,valign=center]
      \tikz[baseline=-0.7ex]\node[circle,fill=white,draw=popink,line width=1.3pt,minimum size=1.6cm,inner sep=0]{\popdisplay\fontsize{24}{24}\selectfont\color{poppurple}+};%
      \hspace{8pt}\begin{minipage}[c]{0.6\linewidth}
        {\popdisplay\fontsize{11}{12.5}\selectfont\color{white}\MakeUppercase{Passe à OnBuch+}}\par\vspace{4pt}
        {\raggedright\popsemi\fontsize{8.4}{11}\selectfont\color{white}Tous les cours signés, Tuteur IA illimité, fiches PDF — onglet OnBuch+ de l'app\par}
      \end{minipage}
    \end{tcolorbox}
  \end{minipage}
  \vspace{8pt}
  \popcontenu
  \vfill
  \signatureonbuch
  \restoregeometry
  \newpage
}

% Rangée « Ce cours contient » (page de garde)
\newcommand{\poptile}[3]{% #1 couleur #2 gros texte #3 légende
  \begin{tcolorbox}[enhanced,colback=#1,colframe=popink,boxrule=2pt,arc=9pt,shadow={2.5pt}{-2.5pt}{0pt}{popink},
      left=6pt,right=6pt,top=9pt,bottom=9pt,halign=center,valign=center,height=1.75cm,width=\linewidth,nobeforeafter]
    {\popdisplay\fontsize{12.5}{13}\selectfont\color{white}\MakeUppercase{#2}}\par\vspace{3pt}
    {\centering\popsemi\fontsize{7.8}{9.5}\selectfont\color{white}#3\par}
  \end{tcolorbox}}
\newcommand{\popcontenu}{%
  \par\noindent{\popxboldls\fontsize{7.5}{7.5}\selectfont\color{popmuted}CE COURS SIGNÉ CONTIENT}\par\vspace{7pt}
  \noindent\begin{minipage}[t]{0.235\linewidth}\poptile{popblue}{Cours}{complet, illustré, pas à pas}\end{minipage}\hfill
  \begin{minipage}[t]{0.235\linewidth}\poptile{poporange}{Méthodes}{et exercices résolus}\end{minipage}\hfill
  \begin{minipage}[t]{0.235\linewidth}\poptile{poppink}{Exercices}{type examen + corrigés}\end{minipage}\hfill
  \begin{minipage}[t]{0.235\linewidth}\poptile{popgreen}{Fiche bilan}{l'essentiel à retenir}\end{minipage}\par}

% Sommaire
\newtcolorbox{sommaire}{enhanced,colback=white,colframe=popink,boxrule=2.2pt,arc=10pt,
  drop shadow=popink,left=18pt,right=18pt,top=13pt,bottom=13pt,before skip=6pt,after skip=20pt,
  before upper={\poppill{Sommaire}{poppurple}{white}\par\vspace{9pt}\popsemi\fontsize{10.6}{14.5}\selectfont\color{popink}}}

% Signature de fin de cours — poussée tout en bas de la dernière page
\newcommand{\findecours}{%
  \par\vfill\nopagebreak
  \begin{center}\popsquiggle{poporange}\end{center}\vspace{4pt}
  \signatureonbuch
}
\AtBeginDocument{\def\llbracket{\mathopen{[\mkern-3.5mu[}}\def\rrbracket{\mathclose{]\mkern-3.5mu]}}} % intervalles d'entiers (glyphe ⟦ absent de la police)
% Glyphes absents de Fira Math : redéfinis à partir de glyphes disponibles
\AtBeginDocument{%
  \def\setminus{\mathbin{\backslash}}\let\smallsetminus\setminus
  \def\vdots{\mathord{\vbox{\baselineskip3.2pt\lineskiplimit0pt\kern4pt\hbox{.}\hbox{.}\hbox{.}}}}%
  \def\ddots{\mathinner{\mkern1mu\raise6.5pt\hbox{.}\mkern2mu\raise3.6pt\hbox{.}\mkern2mu\raise0.7pt\hbox{.}\mkern1mu}}%
  \def\triangle{\mathord{\scalebox{0.9}{$\symup{\Delta}$}}}%
  \def\nabla{\mathord{\raisebox{\depth}{\scalebox{1}[-1]{$\symup{\Delta}$}}}}%
  \def\checkmark{\popcheck{popdark}}%
  \def\star{\mathbin{\ast}}\def\bigstar{\mathord{\ast}}%
  \def\diamond{\mathbin{\scalebox{0.6}{$\Diamond$}}}%
  \def\bigcup{\mathop{\mathchoice{\raisebox{-0.3em}{\scalebox{1.7}{$\cup$}}}{\scalebox{1.25}{$\cup$}}{\cup}{\cup}}\displaylimits}%
  \def\bigcap{\mathop{\mathchoice{\raisebox{-0.3em}{\scalebox{1.7}{$\cap$}}}{\scalebox{1.25}{$\cap$}}{\cap}{\cap}}\displaylimits}%
  \def\lesseqqgtr{\mathrel{\lessgtr}}\def\bigoplus{\mathop{\oplus}\displaylimits}%
}
\providecommand{\ssi}{si et seulement si} % abréviation fréquente
\AtBeginDocument{\providecommand{\wideparen}{\overparen}} % arc de cercle

%% FIGFIX-BEGIN v5 — correctifs globaux des figures (docs/onbuch-generation/tools/figfix)
\usepackage{adjustbox}\usepackage{etoolbox}\tcbuselibrary{raster}
% toute figure TikZ/pgfplots plus large que la ligne est réduite à la largeur disponible
\BeforeBeginEnvironment{tikzpicture}{\begin{adjustbox}{max width=\linewidth}}
\AfterEndEnvironment{tikzpicture}{\end{adjustbox}}
% légendes pgfplots lisibles par-dessus les courbes
\pgfplotsset{every axis/.append style={legend style={fill=white,fill opacity=0.92,text opacity=1,draw=black!15,inner sep=3pt}}}
% étiquettes d'arêtes (syntaxe edge["..."]) avec fond blanc
\tikzset{every edge quotes/.append style={fill=white,inner sep=1.2pt}}
%% FIGFIX-END
\begin{document}

\pagegarde{V : produire un rapport}{Français}{Tle Tech}{Français – classe de Terminale technique}{N17}{9 séances (fiche de progression harmonisée)}{Cette leçon apprend à lire, analyser et rédiger un rapport, texte fonctionnel essentiel dans la vie professionnelle et administrative. Elle couvre la structure externe et interne du rapport, les outils de liaison dans la phrase et la rédaction progressive de l'introduction, du corps et de la conclusion, jusqu'à la production complète et améliorée d'un rapport authentique.
\vspace{4pt}
\begin{multicols}{2}\small
\begin{itemize}
\item À la fin de cette leçon, je sais identifier les caractéristiques d'un texte fonctionnel et d'un texte iconique
\item À la fin de cette leçon, je sais décrire la structure externe (en-tête, objet, signature) et interne (introduction, corps, conclusion) d'un rapport
\item À la fin de cette leçon, je sais employer correctement les adverbes de liaison, les locutions adverbiales et les pronoms relatifs pour articuler mes phrases
\item À la fin de cette leçon, je sais rédiger l'introduction d'un rapport (contexte, objet, annonce du plan)
\item À la fin de cette leçon, je sais rédiger le corps d'un rapport en organisant les faits de manière logique et objective
\item À la fin de cette leçon, je sais rédiger la conclusion d'un rapport (bilan, suggestions, recommandations)
\end{itemize}\end{multicols}}

\begin{sommaire}
\begin{multicols}{2}
\begin{enumerate}
\item Le rapport : un texte fonctionnel
\item Structure externe et interne du rapport
\item Les outils de liaison dans la phrase
\item Rédiger l'introduction et la conclusion
\item Rédiger le corps du rapport et exploiter les documents iconiques
\item Produire et améliorer un rapport complet
\item Activité d'intégration
\item Méthodes et exercices résolus
\item Exercices
\item Corrigés des exercices
\item Fiche bilan
\end{enumerate}
\end{multicols}
\end{sommaire}

\begin{objectifs}
\begin{itemize}
\item À la fin de cette leçon, je sais identifier les caractéristiques d'un texte fonctionnel et d'un texte iconique
\item À la fin de cette leçon, je sais décrire la structure externe (en-tête, objet, signature) et interne (introduction, corps, conclusion) d'un rapport
\item À la fin de cette leçon, je sais employer correctement les adverbes de liaison, les locutions adverbiales et les pronoms relatifs pour articuler mes phrases
\item À la fin de cette leçon, je sais rédiger l'introduction d'un rapport (contexte, objet, annonce du plan)
\item À la fin de cette leçon, je sais rédiger le corps d'un rapport en organisant les faits de manière logique et objective
\item À la fin de cette leçon, je sais rédiger la conclusion d'un rapport (bilan, suggestions, recommandations)
\item À la fin de cette leçon, je sais rédiger un rapport complet à partir d'une situation concrète de la vie courante
\item À la fin de cette leçon, je sais relire, corriger et améliorer un rapport (cohérence, orthographe, présentation)
\item À la fin de cette leçon, je sais rédiger et justifier une démarche, une réponse ou une conclusion
\end{itemize}
\end{objectifs}

\begin{prerequis}
\begin{itemize}
\item Les types de textes (narratif, descriptif, explicatif, argumentatif) vus en classe de Première
\item La phrase simple et la phrase complexe (propositions subordonnées)
\item Les connecteurs logiques de base (d'abord, ensuite, enfin, donc, mais)
\item Le compte rendu et la lettre administrative vus dans les unités précédentes
\item La relecture et la correction orthographique (accords, ponctuation)
\end{itemize}
\end{prerequis}

\newpage

\coursec{Le rapport : un texte fonctionnel}

\textbf{Situation d'accroche.} À la suite des inondations qui ont touché le quartier Ndogpassi à Douala en septembre, le maire demande au préfet un rapport circonstancié pour débloquer une aide d'urgence. Un élève de Terminale technique, stagiaire à la mairie, est chargé de rédiger ce document. Mais comment transformer des observations de terrain, des témoignages et des chiffres en un texte clair, structuré et convaincant ? Quelles parties doit contenir un rapport ? Quels mots de liaison utiliser pour enchaîner les idées ? Cette leçon donne les clés pour répondre à ces questions.

\textbf{Problématique.} Qu'est-ce qu'un rapport ? En quoi se distingue-t-il des autres textes fonctionnels comme le compte rendu ou le procès-verbal ? Quelles qualités doit-il présenter pour être efficace ?

\courssub{Qu'est-ce qu'un texte fonctionnel ?}

\textbf{L'intuition d'abord.} Dans la vie quotidienne, nous n'écrivons pas seulement pour raconter des histoires ou exprimer nos sentiments. Nous écrivons aussi pour \emph{agir} : demander un emploi, signaler un problème, attester d'un fait, réclamer un paiement. Tous ces écrits ont un point commun : ils répondent à un besoin pratique, précis, et ils sont destinés à quelqu'un qui doit \emph{faire quelque chose} après les avoir lus. Ce sont des textes fonctionnels.

\begin{definition}[Le texte fonctionnel]
Un \cle{texte fonctionnel} (ou texte utilitaire, texte pragmatique) est un texte lié à une situation de communication pratique de la vie administrative, professionnelle ou scolaire. Il vise un résultat concret : informer, demander, attester, constater, proposer. Sa réussite se mesure à son \cle{efficacité} : le destinataire doit comprendre immédiatement de quoi il s'agit et ce qu'on attend de lui.
\end{definition}

\textbf{Les grandes familles de textes fonctionnels.} On peut classer les textes fonctionnels selon leur objet :

\begin{itemize}
\item les textes qui \emph{demandent} : la lettre de demande, la requête, la demande d'emploi ;
\item les textes qui \emph{constatent} : le procès-verbal, la facture, l'attestation, le certificat ;
\item les textes qui \emph{rendent compte} : le compte rendu, le rapport ;
\item les textes qui \emph{organisent} : la convocation, l'ordre du jour, la note de service.
\end{itemize}

\begin{popfigure}
\begin{center}
\begin{tikzpicture}[
grow=right,
level 1/.style={sibling distance=2.6cm, level distance=3.6cm},
level 2/.style={sibling distance=1.15cm, level distance=3.4cm},
every node/.style={font=\small},
edge from parent/.style={draw=popmuted, thick}
]
\node[draw=popdark, fill=popblueL, rounded corners, font=\bfseries] {Textes fonctionnels}
  child { node[draw=popblue, fill=popblue!15, rounded corners] {Constater}
    child { node[fill=popcream, rounded corners] {Procès-verbal} }
    child { node[fill=popcream, rounded corners] {Facture, attestation} }
  }
  child { node[draw=poporange, fill=poporangeL, rounded corners, font=\bfseries] {Rendre compte}
    child { node[draw=poporange, fill=poporange!25, rounded corners, font=\bfseries] {LE RAPPORT}
      child { node[fill=poporange!10, rounded corners, font=\footnotesize] {Rapport de stage} }
      child { node[fill=poporange!10, rounded corners, font=\footnotesize] {Rapport d'incident} }
      child { node[fill=poporange!10, rounded corners, font=\footnotesize] {Rapport d'inspection} }
      child { node[fill=poporange!10, rounded corners, font=\footnotesize] {Rapport de mission} }
    }
    child { node[fill=popcream, rounded corners] {Compte rendu} }
  }
  child { node[draw=popgreen, fill=popgreenL, rounded corners] {Demander}
    child { node[fill=popcream, rounded corners] {Lettre administrative} }
    child { node[fill=popcream, rounded corners] {Demande d'emploi} }
  };
\end{tikzpicture}
\end{center}
\legende{Classification simplifiée des textes fonctionnels : le rapport appartient à la famille des textes qui rendent compte, mais il s'en distingue par son analyse et ses propositions.}
\end{popfigure}

\courssub{Définition du rapport}

\textbf{L'intuition d'abord.} Imaginez un médecin qui visite un malade : il observe, il note les symptômes, il pose un diagnostic, puis il prescrit un traitement. Le rapport fonctionne exactement de la même manière : celui qui rédige a \emph{vu}, \emph{constaté}, \emph{vérifié} ; il expose ensuite ce qu'il a observé et, la plupart du temps, il \emph{propose} ce qu'il convient de faire. Le rapport est donc un texte qui regarde les faits en face et qui prépare une décision.

\begin{definition}[Le rapport]
Un \cle{rapport} est un texte fonctionnel écrit qui rend compte objectivement de faits constatés, d'observations ou de l'exécution d'une mission, et qui propose souvent des solutions ou des recommandations. Il est rédigé par une personne chargée d'une enquête, d'une visite, d'un contrôle ou d'une mission, et il est destiné à un destinataire précis : un supérieur hiérarchique, une administration, un jury, un chef d'établissement.
\end{definition}

\begin{aretenir}[Les trois traits essentiels du rapport]
Un texte est un rapport s'il réunit trois caractéristiques :
\begin{enumerate}
\item il repose sur des \cle{faits réels} observés ou vérifiés par son auteur (ou par une équipe) ;
\item il est adressé à un \cle{destinataire déterminé} qui doit décider ou agir ;
\item il ne se contente pas de raconter : il \cle{analyse} et, le plus souvent, \cle{propose} des solutions.
\end{enumerate}
\end{aretenir}

\textbf{Les différents types de rapports.} Selon la mission confiée, on distingue plusieurs types de rapports, tous au programme du Baccalauréat technique :

\begin{center}
\begin{tabularx}{\linewidth}{|l|X|X|}\hline
\rowcolor{popblueL} \textbf{Type de rapport} & \textbf{Circonstance} & \textbf{Exemple concret} \\ \hline
Rapport de stage & Fin d'une période de formation en entreprise & Un élève de Tle F3 rend compte de son stage à la SODECOTON \\ \hline
Rapport d'incident & Un événement anormal s'est produit & Un court-circuit a provoqué un incendie dans l'atelier \\ \hline
Rapport d'inspection & Visite de contrôle officielle & Un inspecteur visite un atelier de mécanique à Bafoussam \\ \hline
Rapport de mission & Accomplissement d'une tâche confiée & Un chef de centre rend compte du déroulement d'un examen \\ \hline
\end{tabularx}
\end{center}

\courssub{Le rapport parmi les textes fonctionnels : rapport ou compte rendu ?}

\textbf{L'intuition d'abord.} Deux élèves assistent à la même réunion du conseil de classe. Le premier note fidèlement tout ce qui a été dit, dans l'ordre où cela a été dit : il fait un \emph{compte rendu}. Le second, chargé par le proviseur d'étudier le problème de l'absentéisme, expose les faits, les explique et suggère des sanctions et des mesures d'encadrement : il rédige un \emph{rapport}. Même réalité, deux textes différents.

\begin{attention}[Erreur fréquente : confondre rapport et compte rendu]
Au Baccalauréat technique, beaucoup de candidats perdent des points parce qu'ils rédigent un compte rendu à la place d'un rapport. Retenez la différence fondamentale : le \cle{compte rendu} \emph{relate fidèlement} des faits ou des paroles, sans les commenter ; le \cle{rapport} \emph{analyse} les faits, en dégage les causes et les conséquences, et se termine généralement par des \emph{suggestions} ou des \emph{recommandations}. Si votre texte ne contient ni analyse ni proposition, ce n'est pas un rapport.
\end{attention}

\begin{center}
\begin{tabularx}{\linewidth}{|l|X|X|}\hline
\rowcolor{popblueL} \textbf{Critère} & \textbf{Compte rendu} & \textbf{Rapport} \\ \hline
Destinataire & Toute personne absente qui veut être informée & Un destinataire précis : supérieur, administration, jury \\ \hline
Ton & Neutre, fidèle, descriptif & Objectif mais analytique, argumenté \\ \hline
Contenu & Les faits et les paroles, dans l'ordre où ils se sont produits & Les faits constatés, leur analyse (causes, conséquences), des suggestions \\ \hline
Finalité & Informer & Informer, éclairer une décision, proposer des solutions \\ \hline
Exemple & Compte rendu d'une réunion de parents d'élèves & Rapport sur les inondations de Ndogpassi \\ \hline
\end{tabularx}
\end{center}

\courssub{Les qualités d'un bon rapport}

Un rapport efficace obéit à quatre exigences, que le correcteur du Baccalauréat technique vérifie systématiquement.

\begin{propriete}[Les quatre qualités d'un bon rapport]
\begin{enumerate}
\item \textbf{L'exactitude des faits} : dates, chiffres, noms et lieux doivent être précis et vérifiables. On écrit « 247 maisons inondées le 14 septembre », et non « beaucoup de maisons, il y a quelque temps ».
\item \textbf{L'objectivité} : le rapporteur expose les faits sans jugement personnel, sans exagération, sans partialité. Il écrit « le pont est impraticable », et non « ce pont honteusement mal construit ».
\item \textbf{L'ordre logique} : les faits sont présentés selon un plan clair (ordre chronologique, ordre des causes aux conséquences, du général au particulier).
\item \textbf{La clarté de la langue} : phrases courtes, vocabulaire précis, mots de liaison explicites (\emph{d'abord}, \emph{ensuite}, \emph{en effet}, \emph{par conséquent}, \emph{enfin}).
\end{enumerate}
\end{propriete}

\begin{methode}[Vérifier qu'un texte est bien un rapport]
Face à un document, posez-vous cinq questions :
\begin{enumerate}
\item Qui est l'auteur et quelle mission lui a été confiée ?
\item À qui le document est-il destiné ?
\item Les faits rapportés sont-ils exacts, datés, chiffrés ?
\item Le texte analyse-t-il les faits (causes, conséquences) ?
\item Se termine-t-il par des suggestions ou des recommandations ?
\end{enumerate}
Si les cinq réponses sont positives, vous êtes en présence d'un rapport.
\end{methode}

\courssub{Lecture guidée d'un rapport authentique}

Lisons ensemble un extrait de rapport rédigé par un chef de centre d'examen, à l'issue du Baccalauréat technique.

\begin{exemplebox}[Extrait commenté : rapport du chef de centre d'examen de Garoua]
\emph{À Monsieur le Délégué régional des Enseignements secondaires de l'Adamaoua.}

\emph{\textbf{Objet :} rapport sur le déroulement des épreuves écrites du Baccalauréat technique, session de juin, centre de Garoua Lycée technique.}

\emph{J'ai l'honneur de rendre compte de ce qui suit. Les épreuves écrites se sont déroulées du lundi 9 au vendredi 13 juin, dans les quatorze salles du centre. Sur 812 candidats inscrits, 796 se sont présentés, soit un taux de présence de 98 \%. Dès le premier jour, nous avons constaté l'insuffisance des ventilateurs dans quatre salles, ce qui a gêné les candidats pendant les épreuves de l'après-midi. En outre, deux cas de fraude ont été relevés et consignés dans des procès-verbaux distincts. Malgré ces difficultés, l'ensemble des épreuves s'est déroulé dans le calme.}

\emph{En conséquence, je suggère : 1° l'installation de dix ventilateurs supplémentaires avant la prochaine session ; 2° le renforcement de la surveillance dans les salles 7 et 12.}

\emph{Fait à Garoua, le 16 juin. Le chef de centre.}
\end{exemplebox}

\textbf{Repérage dans le document.} Appliquons la méthode des cinq questions :

\begin{itemize}
\item \textbf{Auteur et mission} : le chef de centre, chargé d'organiser et de superviser les épreuves.
\item \textbf{Destinataire} : le Délégué régional des Enseignements secondaires, c'est-à-dire le supérieur hiérarchique.
\item \textbf{Objet} : annoncé dès l'en-tête, en une phrase précise.
\item \textbf{Faits rapportés} : dates (du 9 au 13 juin), chiffres (812 inscrits, 796 présents, 98 \%), incidents (ventilateurs insuffisants, deux cas de fraude). Tout est exact, daté, chiffré.
\item \textbf{Analyse et suggestions} : la dernière partie propose deux mesures concrètes, numérotées. C'est ce qui fait de ce texte un \emph{rapport} et non un simple compte rendu.
\end{itemize}

Remarquez aussi les \cle{outils de liaison} qui organisent le texte : \emph{dès le premier jour} (repère chronologique), \emph{en outre} (ajout), \emph{malgré ces difficultés} (opposition), \emph{en conséquence} (conclusion logique). Nous les étudierons en détail dans la troisième partie de la leçon.

\begin{exoresolu}[Identifier les éléments d'un rapport]
\textbf{Énoncé.} Voici un passage : « Le 14 septembre, les eaux ont submergé 247 maisons du quartier Ndogpassi. Cette inondation est due à l'obstruction des caniveaux par les ordures. Je propose par conséquent le curage immédiat des canaux et une campagne de sensibilisation des riverains. » Indiquez : a) les faits rapportés ; b) l'analyse ; c) les suggestions ; d) deux mots de liaison.

\tcblower
\textbf{Solution détaillée.}
\begin{enumerate}
\item \textbf{Faits rapportés} : la date (le 14 septembre), le lieu (quartier Ndogpassi), le chiffre exact (247 maisons submergées). Ce sont des données objectives et vérifiables.
\item \textbf{Analyse} : la cause de l'inondation est identifiée — « l'obstruction des caniveaux par les ordures ». Le rapporteur ne se contente pas de décrire : il explique.
\item \textbf{Suggestions} : deux mesures sont proposées — le curage immédiat des canaux et une campagne de sensibilisation. C'est la marque propre du rapport.
\item \textbf{Mots de liaison} : « est due à » (expression de la cause) et « par conséquent » (expression de la conséquence, qui introduit les propositions).
\end{enumerate}
Ce passage réunit donc les trois traits essentiels du rapport : faits exacts, analyse, propositions.
\end{exoresolu}

\begin{savaistu}[Le rapport de stage au Baccalauréat technique]
Au Cameroun, tout élève de l'enseignement technique qui effectue un stage en entreprise doit remettre un \cle{rapport de stage}. Dans plusieurs séries du Baccalauréat technique, ce rapport fait l'objet d'une évaluation : il est relu, noté et parfois défendu devant un jury. Savoir rédiger un rapport n'est donc pas seulement une compétence d'examen : c'est une compétence professionnelle que vous utiliserez dans toute votre carrière, que vous soyez technicien, comptable, infirmier ou ingénieur des travaux.
\end{savaistu}

\begin{exercice}[À vous de jouer]
Le proviseur de votre lycée vous charge de rédiger un rapport sur l'état de la salle informatique après une visite effectuée le 3 mars. Vous avez noté : 25 ordinateurs sur 40 en panne ; climatiseur hors service ; connexion Internet interrompue depuis janvier.
\begin{enumerate}
\item Rédigez l'en-tête du rapport (destinataire, auteur, objet).
\item Classez ces trois constats dans un ordre logique et justifiez votre choix.
\item Formulez deux suggestions adaptées à ces constats.
\item Dites pourquoi votre texte est un rapport et non un compte rendu.
\end{enumerate}
\end{exercice}

\begin{corrige}[Exercice]
\begin{enumerate}
\item En-tête : « À Monsieur le Proviseur du Lycée technique de \ldots{} » ; « De : [nom], chargé(e) de mission » ; « Objet : rapport sur l'état de la salle informatique, visite du 3 mars ».
\item Ordre logique possible, du plus grave au moins grave : d'abord la connexion Internet interrompue depuis janvier (elle paralyse toute la salle), ensuite les 25 ordinateurs sur 40 en panne (chiffre précis), enfin le climatiseur hors service. On peut aussi choisir l'ordre matériel : machines, puis réseau, puis confort. L'essentiel est de justifier l'ordre adopté.
\item Suggestions : « Je suggère : 1° la réparation ou le remplacement des 25 ordinateurs défectueux ; 2° le rétablissement urgent de la connexion Internet et la remise en service du climatiseur. »
\item Ce texte est un rapport parce qu'il s'adresse à un destinataire précis (le proviseur), rapporte des faits exacts et datés (visite du 3 mars, 25 ordinateurs sur 40), et se termine par des suggestions destinées à éclairer une décision. Un compte rendu se serait contenté de décrire la visite sans rien proposer.
\end{enumerate}
\end{corrige}

\begin{aretenir}[L'essentiel de la section]
Le \cle{rapport} est un texte fonctionnel qui rend compte objectivement de faits constatés au cours d'une mission et qui propose des solutions à un destinataire précis. Il se distingue du \cle{compte rendu}, qui relate fidèlement sans analyser. Un bon rapport exige l'exactitude des faits, l'objectivité, l'ordre logique et la clarté de la langue. Les types de rapports à connaître pour le Baccalauréat technique sont le rapport de stage, le rapport d'incident, le rapport d'inspection et le rapport de mission.
\end{aretenir}

\coursec{Structure externe et interne du rapport}

Dans la section précédente, nous avons identifié le rapport comme un \cle{texte fonctionnel} : un écrit professionnel qui rend compte de faits réels, observés ou vécus, à un destinataire précis (souvent un supérieur hiérarchique ou un commanditaire) en vue d'une décision ou d'un archivage. Savoir \emph{pourquoi} l'on écrit (la finalité) ne suffit pas ; il faut maîtriser \emph{comment} on le construit. Le rapport est un texte \textbf{codifié}, dont la forme répond à des normes strictes que tout technicien, administrateur ou étudiant en stage doit respecter. Cette section détaille l'architecture du rapport en distinguant la \cle{structure externe} (présentation matérielle et administrative) et la \cle{structure interne} (organisation logique du contenu).

\courssub{La structure externe : la présentation matérielle et administrative}

\begin{definition}[Structure externe]
La \cle{structure externe} d'un rapport regroupe l'ensemble des éléments de présentation qui encadrent le texte et lui confèrent sa validité administrative : l'\textbf{en-tête} (institution, lieu, date), la \textbf{référence} (le cas échéant), la mention de l'\textbf{objet}, la \textbf{formule d'appel} au destinataire, le \textbf{corps du texte} (structure interne), la \textbf{signature} de l'auteur et sa \textbf{qualité}. Pour les rapports longs, on y ajoute la \textbf{table des matières} et les \textbf{annexes}.
\end{definition}

Imaginons une situation professionnelle. Tu es chef d'équipe maintenance à la SONARA. Le directeur t'a demandé un rapport sur l'incident survenu à l'unité de distillation. Tu produis une analyse technique parfaite... mais tu omets la référence de la demande, l'objet précis et ta qualité de signataire. Le document ne peut être ni enregistré, ni transmis, ni opposable. La structure externe n'est pas un ornement : c'est la \textbf{condition de recevabilité} du rapport.

\textbf{L'en-tête institutionnel.} Il occupe le haut de la première page (et parfois les suivantes via un en-tête courant). Il contient obligatoirement :
\begin{itemize}
\item le \textbf{nom complet de l'institution} ou de la structure émettrice (ex. : \emph{Ministère de l'Éducation de Base}, \emph{Lycée Technique Industriel de Bafoussam}, \emph{SODECOTON -- Direction Régionale Nord}) ;
\item le \textbf{lieu} de rédaction et la \textbf{date} complète (jour, mois, année), alignés à droite : \emph{Garoua, le 12 mars 2027}.
\end{itemize}

\textbf{La référence administrative.} Elle permet de rattacher le rapport à un dossier ou à une correspondance antérieure. Elle se place sous l'en-tête ou en regard de l'objet. Exemple : \emph{Réf. : Note de service n° 014/DRH/SG du 05/03/2027} ou \emph{Notre réf. : 45/ST/2027}. Elle est indispensable en milieu administratif et en entreprise.

\textbf{La mention de l'objet.} C'est l'élément \textbf{capital} pour le traitement du courrier. L'\cle{objet} résume en une seule ligne nominale le sujet exact du rapport. Il doit respecter trois critères :
\begin{itemize}
\item \textbf{Concision} : une phrase nominale, pas de verbe conjugué, pas de paragraphe ;
\item \textbf{Précision} : il nomme le fait, le lieu, la période ou la mission concernée ;
\item \textbf{Visibilité} : précédé de la mention \emph{Objet :} (souvent en majuscules, souligné ou en gras), centré ou aligné à gauche.
\end{itemize}

\begin{exemplebox}[Formulation correcte de l'objet]
\begin{itemize}
\item \textbf{Objet :} Rapport d'incident sur la pompe n°3 de la station de pompage de Njock-Nkong.
\item \textbf{Objet :} Compte rendu de la mission de réception des travaux du bâtiment administratif.
\item \textbf{Objet :} Rapport de stage de fin de cycle à la CAMWATER -- Agence de Douala-Bonabéri (juin--juillet 2026).
\end{itemize}
\end{exemplebox}

\begin{attention}[Pièges à éviter sur l'objet]
\begin{itemize}
\item Ne pas confondre \emph{Objet} et \emph{Titre}. Le titre peut être plus littéraire (ex. \emph{« L'eau, source de vie : rapport sur la station de Njock »}) ; l'objet est strictement administratif.
\item Ne pas écrire : \emph{Objet : Pour vous informer que la pompe est cassée.} (Phrase verbale, trop long, familier).
\item Ne pas oublier l'objet : le destinataire ne sait pas de quoi il retourne avant de lire le texte.
\end{itemize}
\end{attention}

\textbf{La formule d'appel (salutation).} Le rapport s'adressant à une autorité hiérarchique ou fonctionnelle, on interpelle le destinataire par sa \textbf{fonction} ou son \textbf{titre officiel}, jamais par son nom propre.
\begin{itemize}
\item \emph{Monsieur le Ministre,} / \emph{Madame la Ministre,}
\item \emph{Monsieur le Directeur Général,} / \emph{Madame la Directrice Régionale,}
\item \emph{Monsieur le Proviseur,} / \emph{Madame la Cheffe de Service,}
\item \emph{Monsieur le Président du Conseil d'Administration,}
\end{itemize}
La formule se termine par une virgule et précède directement l'introduction (souvent après un saut de ligne).

\textbf{La signature et la qualité de l'auteur.} En bas à droite du document (après la conclusion), l'auteur appose sa signature manuscrite (ou numérique) et indique \textbf{sous son nom} sa \cle{qualité} : la fonction ou le titre au nom duquel il écrit. C'est ce qui engage la responsabilité.
\begin{itemize}
\item \emph{Fouda André}\par \emph{Ingénieur de maintenance, Chef de quart}
\item \emph{Mvondo Grâce}\par \emph{Élève de Terminale F2, Déléguée de classe}
\item \emph{Aboubakar Hamadou}\par \emph{Stagiaire ingénieur, Département Génie Civil}
\end{itemize}

\begin{aretenir}[Check-list de la structure externe -- Validité administrative]
\begin{itemize}
\item[$\square$] En-tête : Institution + Lieu, Date
\item[$\square$] Référence (si contexte administratif/entreprise)
\item[$\square$] \textbf{Objet :} court, précis, mis en évidence
\item[$\square$] Formule d'appel : par la fonction du destinataire
\item[$\square$] Corps du rapport (structure interne)
\item[$\square$] \textbf{Signature + Qualité} de l'auteur (obligatoire)
\item[$\square$] Pièces jointes / Annexes (mention \emph{P.J.} ou \emph{Annexes :} listées)
\end{itemize}
\end{aretenir}

\courssub{La structure interne : l'organisation logique du contenu}

\begin{definition}[Structure interne]
La \cle{structure interne} organise le discours en trois moments obligatoires et séquencés : l'\textbf{introduction} (présentation de la mission), le \textbf{corps} (exposition ordonnée des faits, constats, analyses), la \textbf{conclusion} (bilan synthétique et recommandations). Cette architecture en « entonnoir » va du général (la mission) au particulier (les faits) pour revenir au général (la décision).
\end{definition}

\textbf{1. L'introduction : situer la mission.} Elle répond, dans cet ordre logique, à quatre questions fondamentales. Elle ne doit pas contenir de résultats.
\begin{enumerate}
\item \textbf{Les circonstances (le déclencheur) :} Qui a demandé ? Quand ? Dans quel cadre ? (\emph{« Sur instruction n° 42/MINEDUB/SG du 10/09/2026, Monsieur l'Inspecteur Pédagogique Régional m'a chargé de... »})
\item \textbf{L'objectif (la finalité) :} Que devais-je constater, vérifier, évaluer, proposer ? (\emph{« L'objectif était d'évaluer l'état de fonctionnement des ateliers de mécanique auto. »})
\item \textbf{La méthode (la démarche) :} Comment ai-je procédé ? (Visite de terrain, entretiens, consultation de registres, essais, observation participante, questionnaire). (\emph{« J'ai visité les trois ateliers les 15 et 16 octobre, interrogé les cinq formateurs et consulté les registres de maintenance. »})
\item \textbf{L'annonce du plan :} Comment le rapport est-il structuré ? (\emph{« Nous présenterons successivement : I. L'état des équipements, II. La gestion des consommables, III. Les besoins en formation. »})
\end{enumerate}

\textbf{2. Le corps du rapport : exposer les faits.} C'est la partie la plus développée. Les faits sont présentés de manière \textbf{objective}, \textbf{ordonnée} et \textbf{chiffrée} autant que possible. Le plan suit l'annonce de l'introduction. Trois ordres principaux guident le découpage en parties (numérotées I, II, III ou 1, 2, 3) :
\begin{itemize}
\item \textbf{Ordre chronologique} : déroulement des événements (idéal : rapport d'accident, journal de stage, compte rendu de mission). Ex. : \emph{Jour 1 : ... ; Jour 2 : ...}
\item \textbf{Ordre thématique (ou analytique)} : regroupement par domaines, secteurs, thèmes (idéal : rapport d'inspection, d'audit, d'état des lieux). Ex. : \emph{I. Bâtiments, II. Équipements, III. Personnel, IV. Sécurité.}
\item \textbf{Ordre d'importance (ou gravité)} : du plus critique au moins critique, ou inversément (idéal : rapport d'incident, de dysfonctionnement). Ex. : \emph{I. Risques majeurs (sécurité), II. Pannes bloquantes (production), III. Améliorations confort.}
\end{itemize}
\begin{attention}[Rédiger le corps : exigences techniques]
\begin{itemize}
\item \textbf{Un fait = une phrase (ou un paragraphe) précise.} Éviter le vague : \emph{« Plusieurs machines sont vieilles »} $\rightarrow$ \emph{« Sur les 12 tours parallèles, 8 datent de 1998 et présentent un jeu excessif sur la broche. »}
\item \textbf{Chiffrer, dater, localiser.} \emph{« La toiture du bloc B fuit »} $\rightarrow$ \emph{« Trois fuites actives constatées le 12/10 sur la toiture tôle du bloc B (salle B12, B14, couloir central). »}
\item \textbf{Numérotation continue et titres clairs.} I, II, III (chiffres romains) pour les grandes parties ; 1, 2, 3 (arabes) pour les sous-parties. Chaque partie porte un titre informatif.
\item \textbf{Utiliser les outils de liaison} (adverbes, locutions, pronoms relatifs) pour assurer la cohésion entre paragraphes et parties (voir section 3 de la leçon).
\end{itemize}
\end{attention}

\textbf{3. La conclusion : décider et proposer.} Elle ne introduit \textbf{aucun fait nouveau}. Elle comporte deux temps distincts :
\begin{enumerate}
\item \textbf{Le bilan synthétique (ou appréciation globale) :} Résumé en 2--3 phrases des constats majeurs. On ne répète pas les détails. (\emph{« L'inspection révèle une vétusté critique des machines-outils du secteur tournage, un stock de consommables sous le seuil d'alerte et un besoin urgent de formation continue pour 3 formateurs sur 5. »})
\item \textbf{Les suggestions / recommandations / propositions :} Adressées au destinataire, concrètes, hiérarchisées (urgent / court terme / moyen terme), chiffrées si possible. C'est souvent la \emph{raison d'être} du rapport. (\emph{« Nous recommandons : 1) Le remplacement immédiat des 3 tours les plus dégradés (budget estimé 18 millions FCFA) ; 2) Le lancement d'une commande groupée de plaquettes carbure ; 3) L'inscription d'un module CNC au plan de formation 2027. »})
\end{enumerate}

\begin{popfigure}
\begin{center}
\begin{tikzpicture}[scale=0.9, every node/.style={font=\small}]
% Cadre page
\draw[thick, popdark] (0,0) rectangle (10,13.5);
% En-tête
\draw[fill=popblueL, draw=popblue] (0.5,12.1) rectangle (5.5,13.1);
\node[align=center] at (3,12.6) {INSTITUTION\\ \emph{(ex. Lycée Technique de Douala)}};
\draw[fill=popblueL, draw=popblue] (6,12.1) rectangle (9.5,13.1);
\node[align=center] at (7.75,12.6) {Lieu, Date\\ \emph{Douala, le 15/10/2026}};
% Référence
\draw[fill=popblue!10, draw=popblue] (0.5,11.4) rectangle (9.5,11.9);
\node[anchor=west] at (0.7,11.65) {\textbf{Réf.} : Note de service n° 042/SG du 03/10/2026};
% Objet
\draw[fill=popgoldL, draw=popgold, thick] (0.5,10.4) rectangle (9.5,11.2);
\node[anchor=west] at (0.7,10.8) {\textbf{OBJET :} RAPPORT SUR L'ÉTAT DES ATELIERS DE MÉCANIQUE AUTO};
% Appel
\node[anchor=west, font=\small] at (0.5,10.0) {Monsieur le Proviseur,};
% Introduction
\draw[fill=popgreenL, draw=popgreen] (0.5,8.3) rectangle (9.5,9.7);
\node at (5,9.0) {\textbf{INTRODUCTION} : Circonstances -- Objectif -- Méthode -- Plan};
% Corps
\draw[fill=poporangeL, draw=poporange] (0.5,4.5) rectangle (9.5,8.0);
\node at (5,7.6) {\textbf{CORPS DU RAPPORT}};
\node[anchor=west] at (0.8,7.0) {\textbf{I. État des machines-outils (tours, fraiseuses)}};
\node[anchor=west] at (0.8,6.5) {\textbf{II. Gestion des consommables et outillage}};
\node[anchor=west] at (0.8,6.0) {\textbf{III. Compétences du personnel et sécurité}};
\node[anchor=west, font=\itshape] at (0.8,5.4) {Faits précis, datés, chiffrés, objectifs -- Liaisons logiques};
% Conclusion
\draw[fill=poppinkL, draw=poppink] (0.5,2.8) rectangle (9.5,4.2);
\node at (5,3.5) {\textbf{CONCLUSION} : Bilan synthétique + Recommandations hiérarchisées};
% Signature
\draw[fill=popblueL, draw=popblue] (6,1.0) rectangle (9.5,2.3);
\node[align=center] at (7.75,1.65) {Signature\\ \emph{Nom Prénom}\\ \emph{Qualité (ex. Chef de travaux)}};
% Annexes
\node[anchor=west, font=\small] at (0.5,0.5) {\textbf{Annexes :} 1. Photos machines -- 2. Extraits registres -- 3. Devis fournisseurs};
% Légendes fléchées
\draw[-{Stealth}, thick, popblue] (10.1,12.6) -- (12,12.6) node[right, align=left]{En-tête institutionnel};
\draw[-{Stealth}, thick, popblue] (10.1,11.65) -- (12,11.65) node[right]{Référence administrative};
\draw[-{Stealth}, thick, popgold] (10.1,10.8) -- (12,10.8) node[right, align=left]{Objet : court, précis,\\ mis en évidence};
\draw[-{Stealth}, thick, popgreen] (10.1,9.0) -- (12,9.0) node[right]{Introduction : 4 étapes};
\draw[-{Stealth}, thick, poporange] (10.1,6.2) -- (12,6.2) node[right, align=left]{Corps : parties I, II, III\\ ordonnées, factuelles};
\draw[-{Stealth}, thick, poppink] (10.1,3.5) -- (12,3.5) node[right]{Conclusion : Bilan + Propositions};
\draw[-{Stealth}, thick, popblue] (10.1,1.65) -- (12,1.65) node[right, align=left]{Signature + Qualité\\ (engagement)};
\draw[-{Stealth}, thick, popmuted] (10.1,0.5) -- (12,0.5) node[right]{Annexes / Pièces jointes};
\end{tikzpicture}
\end{center}
\legende{Archétype d'une page de rapport technique : la structure externe (zones bleues/jaunes) encadre et valide la structure interne (zones verte/orange/rose).}
\end{popfigure}

\begin{popfigure}
\begin{center}
\begin{tikzpicture}[scale=1.1]
% Entonnoir 3 niveaux
\draw[fill=popgreenL, draw=popgreen, thick] (0,6.5) -- (9,6.5) -- (7.5,4.8) -- (1.5,4.8) -- cycle;
\node[align=center, font=\small] at (4.5,5.65) {\textbf{INTRODUCTION}\\[2pt] \emph{Le Général : Mission, Objectif, Méthode, Plan}\\[2pt] \textit{(Pas de résultats)}};
\draw[fill=poporangeL, draw=poporange, thick] (1.5,4.5) -- (7.5,4.5) -- (6.2,2.5) -- (2.8,2.5) -- cycle;
\node[align=center, font=\small] at (4.5,3.5) {\textbf{CORPS}\\[2pt] \emph{Le Détaillé : Faits, Constats, Analyses}\\[2pt] \textit{Parties I, II, III -- Ordre choisi -- Chiffré}};
\draw[fill=poppinkL, draw=poppink, thick] (2.8,2.2) -- (6.2,2.2) -- (5.2,0.5) -- (3.8,0.5) -- cycle;
\node[align=center, font=\small] at (4.5,1.35) {\textbf{CONCLUSION}\\[2pt] \emph{L'Essentiel : Bilan + Recommandations}\\[2pt] \textit{(Adressées au destinataire)}};
% Flèches
\draw[-{Stealth}, very thick, popdark] (4.5,4.75) -- (4.5,4.55);
\draw[-{Stealth}, very thick, popdark] (4.5,2.45) -- (4.5,2.25);
% Latéral : flèche "Précision croissante"
\draw[-{Stealth[length=3mm]}, thick, popdark] (9.5,5.6) -- (9.5,1.3) node[right, midway, align=left, font=\small] {Niveau de détail\\ \textbf{croissant} $\downarrow$};
\draw[-{Stealth[length=3mm]}, thick, popdark] (9.5,1.3) -- (9.5,5.6) node[left, midway, align=left, font=\small] {Portée décisionnelle\\ \textbf{croissante} $\uparrow$};
\end{tikzpicture}
\end{center}
\legende{Le rapport en entonnoir : l'introduction cadre la mission (général), le corps détaille les faits (précis), la conclusion resserre vers la décision (essentiel).}
\end{popfigure}

\courssub{Vérifier la structure d'un rapport : méthode pas à pas et exemple commenté}

\begin{methode}[Contrôler la conformité structurelle d'un rapport en 4 étapes]
\begin{enumerate}
\item \textbf{Valider l'externe (1 min) :} Repérer \textbf{Objet} (court, précis, visible) et \textbf{Signature+Qualité} (bas à droite). Sans eux $\rightarrow$ \textbf{irrecevable}.
\item \textbf{Analyser l'introduction (2 min) :} Cocher les 4 items : \emph{Circonstances} -- \emph{Objectif} -- \emph{Méthode} -- \emph{Plan}. Le plan annoncé correspond-il aux titres des parties du corps ?
\item \textbf{Parcourir le corps (3 min) :} Souligner les numéros/titres I, II, III. Vérifier : ordre logique (chrono/thématique/importance) respecté ? Faits objectifs, datés, chiffrés ? Liaisons fluides entre parties ?
\item \textbf{Évaluer la conclusion (1 min) :} Bilan synthétique (pas de copier-coller) + Recommandations concrètes, adressées au destinataire, hiérarchisées.
\end{enumerate}
\end{methode}

\begin{exemplebox}[Plan type corrigé -- Rapport de stage (Structure interne)]
\textbf{Objet :} Rapport de stage effectué à l'entreprise \textsc{CAMBOIS} (Douala) du 1\textsuperscript{er} au 30 juin 2026.\par
\textbf{Introduction} : Stage de fin de cycle (Terminale F4) -- Objectif : confronter acquis scolaires et réalité industrielle -- Méthode : observation participative (ateliers sciage, séchage, finition) + entretiens tuteurs -- Plan annoncé en 3 parties.\par
\textbf{I. Présentation de l'entreprise} : Historique, organigramme, capacité de production, certifications.\par
\textbf{II. Activités techniques réalisées} : Tâches confiées (réglage scie à ruban, contrôle humidité bois, maintenance préventive niveau 1), difficultés (vibrations scie n°2), solutions apportées.\par
\textbf{III. Appréciation critique et acquis} : Compétences renforcées (lecture plans, sécurité), écarts programme/réalité, suggestions pour l'entreprise (maintenance prédictive) et pour l'école (ajout module CNC).\par
\textbf{Conclusion} : Bilan global positif (autonomie acquise) -- Recommandations : 1) Convention école-entreprise pour stages longs ; 2) Dotation lycée en scie à commande numérique.\par
\textbf{Signature + Qualité} : \emph{Mbarga Éric, Élève Terminale F4, Stagiaire.}
\end{exemplebox}

\begin{exoresolu}[Reconstituer et justifier la structure d'un rapport mélangé]
\textbf{Énoncé.} Un proviseur a rédigé un rapport sur l'état des salles de classe après la saison des pluies. Les cinq blocs ci-dessous sont mélangés. Remets-les dans l'ordre chronologique et logique, puis justifie la place de chaque bloc en citant les critères de la structure externe et interne.\par
\textbf{A.} En conclusion, trois blocs sur cinq nécessitent une réfection urgente avant la rentrée ; nous suggérons de saisir la mairie pour le financement des tôles et de missionner l'entreprise BATICAM pour les maçonneries.\par
\textbf{B.} \textbf{I. État des toitures :} Les blocs A et C ont perdu 12 tôles chacun (vent violent du 12/09). \textbf{II. État des murs :} Fissures traversantes sur le pignon ouest du bloc B (recul de 3 cm). \textbf{III. État du mobilier :} 40 tables-bancs hors d'usage (pieds sectionnés, plateaux pourris), 120 récupérables après ponçage.\par
\textbf{C.} \textsc{Lycée Technique de Maroua} -- Maroua, le 20 octobre 2026. \textbf{Objet :} Rapport sur l'état des salles de classe après la saison des pluies 2026. \emph{Monsieur le Délégué Régional de l'Éducation Secondaire,}\par
\textbf{D.} \emph{Signé : Etoundi Paul, Proviseur.}\par
\textbf{E.} À la suite des pluies diluviennes de septembre, vous m'avez chargé par note n° 11/DRSE du 05/10 d'évaluer les dégâts. J'ai visité les cinq blocs du 12 au 15 octobre avec le chef de travaux. Nous examinerons successivement les toitures, les murs, puis le mobilier.
\tcblower
\textbf{Solution détaillée.}\par
L'ordre unique et correct est : \textbf{C $\rightarrow$ E $\rightarrow$ B $\rightarrow$ A $\rightarrow$ D}.\par
\noindent\textbf{Bloc C (Structure externe -- Ouverture) :} Contient l'en-tête (Institution, Lieu, Date), la Référence implicite (note n° 11), l'\textbf{Objet} (court, précis, mis en évidence) et la \textbf{Formule d'appel} (par la fonction : \emph{Monsieur le Délégué Régional...}). C'est la « couverture » administrative.\par
\noindent\textbf{Bloc E (Introduction -- Structure interne, temps 1) :} Les 4 items y sont : \emph{Circonstances} (\emph{pluies, note n° 11}), \emph{Objectif} (\emph{évaluer les dégâts}), \emph{Méthode} (\emph{visite 12--15 oct, avec chef travaux}), \emph{Annonce du plan} (\emph{toitures, murs, mobilier}). Aucun résultat chiffré.\par
\noindent\textbf{Bloc B (Corps -- Structure interne, temps 2) :} Trois parties \textbf{I, II, III} qui suivent \textbf{exactement} l'ordre annoncé en E (thématique). Faits \textbf{précis, datés, chiffrés} (12 tôles, 3 cm, 40 tables-bancs). Langage objectif, technique.\par
\noindent\textbf{Bloc A (Conclusion -- Structure interne, temps 3) :} \textbf{Bilan synthétique} (\emph{trois blocs sur cinq, réfection urgente}) + \textbf{Recommandations adressées au destinataire} (\emph{saisir la mairie, missionner BATICAM}). Pas de nouveau fait.\par
\noindent\textbf{Bloc D (Structure externe -- Fermeture) :} \textbf{Signature} + \textbf{Qualité} (\emph{Proviseur}). Engage la responsabilité de l'auteur et valide le document.
\end{exoresolu}

\begin{exercice}[Application guidée : Rédiger la structure complète d'un rapport]
Le délégué de ta classe de Terminale F3 doit rendre compte au proviseur des \textbf{coupures d'électricité récurrentes} qui perturbent les TP de l'atelier d'électrotechnique depuis la rentrée.
\begin{enumerate}
\item Rédige l'\textbf{en-tête complet} (institution fictive, lieu, date d'aujourd'hui), la \textbf{référence} (inventée : note du proviseur), l'\textbf{objet} (normé) et la \textbf{formule d'appel}.
\item Propose un \textbf{plan du corps en 3 parties numérotées et titrées}, en justifiant le choix de l'ordre (chronologique, thématique ou d'importance).
\item Rédige une \textbf{conclusion de 4 phrases} : 2 phrases pour le bilan synthétique, 2 phrases pour des recommandations concrètes et hiérarchisées (une urgence, une moyen terme).
\item Termine par la \textbf{signature} et la \textbf{qualité} de l'auteur.
\end{enumerate}
\end{exercice}

\begin{corrige}[Exercice -- Éléments attendus et pistes de rédaction]
\textbf{1. Structure externe (exemple) :}\par
\textsc{Lycée Technique Industriel de Yaoundé} \hfill Yaoundé, le 15 octobre 2026\par
\textbf{Réf.} : Note de service n° 007/PROV/2026 du 10/10/2026\par
\textbf{OBJET :} RAPPORT SUR LES COUPURES D'ÉLECTRICITÉ À L'ATELIER D'ÉLECTROTECHNIQUE\par
\emph{Monsieur le Proviseur,}\par
\vspace{2mm}
\textbf{2. Plan du corps (Ordre thématique justifié) :}\par
Les coupures touchent des domaines distincts (matériel, pédagogie, causes) $\rightarrow$ \textbf{Ordre thématique}.\par
\textbf{I. Fréquence, durée et horaires des coupures} (constats bruts : dates, heures, durée).\par
\textbf{II. Conséquences sur les activités pédagogiques et le matériel} (TP annulés, machines endommagées, sécurité).\par
\textbf{III. Causes identifiées et responsabilités} (réseau ENEO, installation interne vétuste, surcharge ateliers).\par
\vspace{2mm}
\textbf{3. Conclusion (4 phrases) :}\par
\textit{Bilan :} « Les coupures quotidiennes de 2 à 4 heures (8h--12h) depuis le 15 septembre ont entraîné l'annulation de 12 séances de TP sur 20 prévues et la détérioration de 3 variateurs de vitesse par surtension. »\par
\textit{Recommandations :} « En urgence, nous demandons l'installation d'un onduleur triphasé 20 kVA pour protéger les équipements sensibles. À moyen terme, nous suggérons la rénovation complète du tableau général de l'atelier et la signature d'un contrat de délestage programmé avec ENEO. »\par
\vspace{2mm}
\textbf{4. Signature :}\par
\emph{Mendouga Armand}\par
\emph{Délégué de la classe de Terminale F3}
\end{corrige}

\begin{aretenir}[Synthèse examinable -- Structure du rapport]
\begin{itemize}
\item \textbf{Externe (Validité) :} En-tête (Institution, Lieu, Date) -- Référence -- \textbf{Objet} (court, précis, visible) -- Appel (Fonction) -- \textbf{Signature + Qualité} -- Annexes (P.J.).
\item \textbf{Interne (Lisibilité/Efficacité) :} \textbf{Introduction} (4 items : Circonstances, Objectif, Méthode, Plan) -- \textbf{Corps} (Parties I, II, III ordonnées : chrono / thématique / importance ; faits objectifs, datés, chiffrés ; liaisons logiques) -- \textbf{Conclusion} (Bilan synthétique + Recommandations concrètes, hiérarchisées, adressées).
\item \textbf{Sanctions :} Absence Objet ou Signature = \textbf{0/5} en structure (examen) / Document irrecevable (vie pro). Plan non annoncé ou non respecté = incohérence majeure.
\end{itemize}
\end{aretenir}

\coursec{Les outils de liaison dans la phrase}

La section précédente a permis d'identifier la \textbf{structure externe} (page de garde, sommaire, introduction, corps, conclusion, annexes) et la \textbf{structure interne} (organisation thématique, paragraphes types) du rapport. Mais une architecture solide ne suffit pas : il faut encore assurer la \textbf{fluidité} et la \textbf{cohérence} du discours. C'est la fonction des \emph{outils de liaison} (connecteurs logiques, pronoms relatifs) qui transforment une suite de phrases disjointes en un texte uni, progressif et argumenté. Maîtriser ces outils, c'est garantir que le correcteur — ou le décideur qui lira votre rapport technique — suit sans effort votre raisonnement.

\courssub{Les adverbes de liaison : définition, rôle et classification}

\begin{definition}[Adverbe de liaison]
Un \cle{adverbe de liaison} (ou connecteur logique) est un \textbf{mot invariable} qui relie deux phrases ou deux propositions en exprimant une \textbf{relation logique} précise (addition, opposition, cause, conséquence, chronologie, etc.). Contrairement aux conjonctions de coordination (\emph{mais, ou, et, donc, or, ni, car}), l'adverbe de liaison conserve une autonomie syntaxique : il peut souvent se déplacer dans la phrase ou être supprimé sans briser la correction grammaticale, mais en perdant la nuance logique.
\end{definition}

\begin{aretenir}[Rôle dans le rapport]
Dans un rapport technique, les adverbes de liaison :
\begin{itemize}
    \item structurent la \textbf{démonstration} (cause $\rightarrow$ conséquence) ;
    \item organisent le \textbf{récit} chronologique des faits (visite, expérience, enquête) ;
    \item nuancent l'\textbf{analyse} (opposition, concession, restriction) ;
    \item guident le lecteur pas à pas vers les \textbf{recommandations}.
\end{itemize}
\end{aretenir}

Le choix du connecteur ne doit jamais être aléatoire. Il découle de la \textbf{relation logique} que l'auteur veut établir entre l'idée A (antérieure) et l'idée B (postérieure). Le tableau ci-dessous répertorie les principaux outils classés par fonction logique, avec des exemples contextualisés dans l'écriture d'un rapport.

\begin{popfigure}
\centering
\begin{tikzpicture}[
    cell/.style={rectangle, draw=popline, minimum height=1.1cm, text width=3.2cm, align=center, font=\small},
    header/.style={rectangle, fill=popblue, text=white, font=\bfseries\small, minimum height=0.9cm, text width=3.2cm, align=center},
    func/.style={rectangle, fill=popblueL, font=\bfseries\small, minimum height=1.1cm, text width=2.8cm, align=center, draw=popline},
    headerfunc/.style={rectangle, fill=popdark, text=white, font=\bfseries\small, minimum height=0.9cm, text width=2.8cm, align=center}
]
% En-têtes
\node[headerfunc] (hf) at (0,0) {Fonction logique};
\node[header] (h1) at (3.0,0) {Adverbes simples};
\node[header] (h2) at (6.2,0) {Locutions adverbiales};
\node[header] (h3) at (9.4,0) {Exemple en contexte (Rapport)};

% Ligne 1 : Addition
\node[func, align=center] (f1) at (0,-1.0){Addition \\ \textit{(ajouter une idée)}};
\node[cell] (c11) at (3.0,-1.0) {de plus, aussi, encore, ensuite};
\node[cell] (c12) at (6.2,-1.0) {en outre, par ailleurs, d'une part... d'autre part, de surcroît};
\node[cell] (c13) at (9.4,-1.0) {L'usine traite les eaux usées. \textbf{De plus}, elle recycle les boues.};

% Ligne 2 : Opposition
\node[func, align=center] (f2) at (0,-2.1){Opposition / Concession \\ \textit{(nuancer, contraster)}};
\node[cell] (c21) at (3.0,-2.1) {cependant, néanmoins, toutefois, pourtant};
\node[cell] (c22) at (6.2,-2.1) {en revanche, par contre, malgré tout, cela dit};
\node[cell] (c23) at (9.4,-2.1) {Le budget est validé. \textbf{Néanmoins}, les délais restent serrés.};

% Ligne 3 : Cause
\node[func, align=center] (f3) at (0,-3.2){Cause / Explication \\ \textit{(justifier)}};
\node[cell] (c31) at (3.0,-3.2) {en effet, car (conj.), c'est que};
\node[cell] (c32) at (6.2,-3.2) {du fait de, en raison de, vu que, attendu que};
\node[cell] (c33) at (9.4,-3.2) {La production a chuté. \textbf{En effet}, la machine principale est en panne.};

% Ligne 4 : Conséquence
\node[func, align=center] (f4) at (0,-4.3){Conséquence / Conclusion \\ \textit{(déduire, décider)}};
\node[cell] (c41) at (3.0,-4.3) {donc, ainsi, alors, par conséquent, c'est pourquoi};
\node[cell] (c42) at (6.2,-4.3) {par suite, en conséquence, de ce fait, au final};
\node[cell] (c43) at (9.4,-4.3) {Le stock est épuisé. \textbf{Par conséquent}, la livraison est reportée.};

% Ligne 5 : Chronologie
\node[func, align=center] (f5) at (0,-5.4){Chronologie / Énumération \\ \textit{(ordonner dans le temps)}};
\node[cell] (c51) at (3.0,-5.4) {d'abord, puis, ensuite, enfin, finalement};
\node[cell] (c52) at (6.2,-5.4) {en premier lieu, en second lieu, pour commencer, en définitive};
\node[cell] (c53) at (9.4,-5.4) {\textbf{D'abord}, visite du site. \textbf{Ensuite}, réunion technique. \textbf{Enfin}, rédaction.};

% Lignes de séparation
\draw[popline] (-1.4,0.45) -- (11.0,0.45);
\draw[popline] (-1.4,-5.95) -- (11.0,-5.95);
\foreach \y in {0,-1.0,-2.1,-3.2,-4.3,-5.4} {
    \draw[popline] (-1.4,\y-0.55) -- (11.0,\y-0.55);
}
\draw[popline] (-1.4,0.45) -- (-1.4,-5.95);
\draw[popline] (1.6,0.45) -- (1.6,-5.95);
\draw[popline] (4.6,0.45) -- (4.6,-5.95);
\draw[popline] (7.8,0.45) -- (7.8,-5.95);
\draw[popline] (11.0,0.45) -- (11.0,-5.95);
\end{tikzpicture}
\legende{Tableau de correspondance : fonctions logiques et outils de liaison (adverbes et locutions) pour le rapport technique.}
\end{popfigure}

\begin{exemplebox}[Variété des locutions adverbiales]
Les \cle{locutions adverbiales} (groupes de mots figés jouant le rôle d'adverbe) offrent une palette plus fine et souvent plus formelle, adaptée au style administratif et technique :
\begin{itemize}
    \item \textbf{Énumération structurée :} \emph{En premier lieu}, l'équipe a inspecté les cuves ; \emph{en second lieu}, elle a prélevé des échantillons ; \emph{en définitive}, elle a dressé le procès-verbal.
    \item \textbf{Articulation argumentative :} \emph{D'une part}, le coût est maîtrisé ; \emph{d'autre part}, la qualité est certifiée.
    \item \textbf{Transition thématique :} \emph{Par ailleurs}, notons la vétusté du réseau électrique (changement de sujet sans rupture brutale).
    \item \textbf{Synthèse :} \emph{Dans l'ensemble}, la mission a atteint ses objectifs.
\end{itemize}
\end{exemplebox}

\begin{attention}[Piège : Adverbe vs Conjonction]
Ne confondez pas \textbf{adverbe de liaison} et \textbf{conjonction de coordination} (\emph{mais, ou, et, donc, or, ni, car}).
\begin{itemize}
    \item \emph{Donc} (conjonction) lie deux propositions \emph{à l'intérieur d'une même phrase} : \og La panne est grave, \textbf{donc} l'arrêt est immédiat.\fg
    \item \emph{Donc} (adverbe) lie deux \emph{phrases distinctes} : \og La panne est grave. \textbf{Donc}, l'arrêt est immédiat.\fg
    \item \emph{Car} est \textbf{toujours} une conjonction (jamais en début de phrase). Pour introduire une phrase explicative, utilisez \emph{En effet}, \emph{En raison de cela}.
\end{itemize}
Au Probatoire, la ponctuation correcte autour de ces mots (point, point-virgule, virgule) est notée.
\end{attention}

\courssub{Les pronoms relatifs : reprendre l'antécédent pour éviter la répétition}

\begin{definition}[Pronom relatif]
Un \cle{pronom relatif} est un mot qui \textbf{remplace un nom ou un pronom} (appelé \emph{antécédent}) et qui \textbf{introduit une proposition subordonnée relative} apportant une information sur cet antécédent. Il assure la \textbf{cohésion référentielle} : il évite de répéter le nom tout en liant deux énoncés en un seul énoncé complexe.
\end{definition}

\begin{propriete}[Les pronoms relatifs simples et leurs fonctions]
Le pronom relatif prend la fonction grammaticale \emph{à l'intérieur de la relative} (sujet, COD, COI, complément du nom, complément circonstanciel de lieu/temps).
\begin{center}
\begin{tabularx}{\linewidth}{|l|X|X|}\hline
\textbf{Pronom} & \textbf{Fonction dans la relative} & \textbf{Antécédent typique / Exemple} \\ \hline
\textbf{qui} & Sujet & \emph{L'ingénieur \textbf{qui} a signé le rapport...} (l'ingénieur \emph{est} le sujet de \emph{a signé}) \\ \hline
\textbf{que} (qu') & COD & \emph{Le rapport \textbf{que} l'ingénieur a signé...} (le rapport \emph{est} l'objet de \emph{a signé}) \\ \hline
\textbf{quoi} & COD (après préposition, chose) & \emph{C'est à \textbf{quoi} nous devons réfléchir.} \\ \hline
\textbf{dont} & Complément du nom (objet de la préposition \emph{de}) & \emph{Les dégâts \textbf{dont} nous avons parlé...} (nous avons parlé \emph{de} les dégâts) \\ \hline
\textbf{où} & CCL (lieu) ou CCT (temps) & \emph{Le quartier \textbf{où} les eaux sont montées...} / \emph{Le jour \textbf{où} la mission a commencé...} \\ \hline
\end{tabularx}
\end{center}
\end{propriete}

\begin{propriete}[Les pronoms relatifs composés (lequel, laquelle, lesquels, lesquelles)]
Ils s'emploient \textbf{après une préposition} (à, de, pour, avec, sans, sous, sur, dans, par...) lorsque l'antécédent désigne une \textbf{chose} (ou un animal) et que l'on veut éviter l'ambiguïté ou la lourdeur. Ils s'accordent en genre et en nombre avec l'antécédent.
\begin{center}
\begin{tabularx}{\linewidth}{|l|X|}\hline
\textbf{Préposition + Lequel} & \textbf{Exemple} \\ \hline
\textbf{auquel} (à + lequel) & \emph{La norme \textbf{auquel} ce matériel répond...} \\ \hline
\textbf{duquel} (de + lequel) & \emph{L'usine \textbf{duquel} dépend l'atelier...} \\ \hline
\textbf{sur laquelle}, \emph{dans lesquels}, etc. & \emph{Le site \textbf{sur lequel} nous intervenons...} \\ \hline
\end{tabularx}
\end{center}
\end{propriete}

\begin{experience}[Fusion de deux phrases par un pronom relatif : observation guidée]
\emph{Observons le mécanisme de transformation (fusion) qui crée la subordonnée relative.}

\textbf{Phrase 1 (Antécédent) :} \og L'équipe a inspecté \textbf{les cuves de stockage}.\fg \\
\textbf{Phrase 2 (Information à ajouter) :} \og \textbf{Les cuves de stockage} présentent des fissures.\fg

\textbf{Étape 1 :} Identifier l'antécédent commun : \emph{les cuves de stockage}.
\textbf{Étape 2 :} Identifier la fonction de l'antécédent dans la Phrase 2 : \emph{Sujet} de \emph{présentent}.
\textbf{Étape 3 :} Choisir le pronom relatif sujet $\rightarrow$ \textbf{qui}.
\textbf{Étape 4 :} Remplacer l'antécédent dans la Phrase 2 par \emph{qui} et placer la relative après l'antécédent dans la Phrase 1.

$\rightarrow$ \og L'équipe a inspecté les cuves de stockage \textbf{qui} présentent des fissures.\fg
\end{experience}

\begin{popfigure}
\centering
\begin{tikzpicture}[
    node distance=1.2cm and 1.5cm,
    box/.style={rectangle, draw=popblue, fill=popblueL, rounded corners, minimum height=1.0cm, text width=3.5cm, align=center, font=\small},
    arrow/.style={-Stealth, thick, poporange},
    label/.style={font=\footnotesize, text=popdark, above, sloped}
]
% Phrase 1
\node[box] (p1a) {L'équipe a inspecté};
\node[box, fill=popgoldL, draw=popgold] (ant1) [right=of p1a] {les cuves de stockage};
\node[box] (p1b) [right=of ant1] {.};

% Phrase 2
\node[box] (p2a) [below=2.5cm of p1a] {Les cuves de stockage};
\node[box, fill=popgoldL, draw=popgold] (ant2) [right=of p2a] {présentent des fissures};
\node[box] (p2b) [right=of ant2] {.};

% Flèches de reprise
\draw[arrow, bend left=20] (ant1.north east) to node[label] {Antécédent commun} (ant2.north west);
\draw[arrow, bend right=20] (ant2.south west) to node[label, below] {Remplacé par \textbf{qui}} (ant1.south east);

% Résultat
\node[box, fill=popgreenL, draw=popgreen, text width=10cm] (res) [below=2.5cm of p2a, xshift=1.5cm] {
    \textbf{Résultat fusionné :} L'équipe a inspecté les cuves de stockage \textcolor{popgreen}{\textbf{qui}} présentent des fissures.
};
\draw[arrow, dashed] (ant2) -- (res);
\end{tikzpicture}
\legende{Schéma de la reprise de l'antécédent \og les cuves de stockage\fg par le pronom relatif \og qui\fg (fonction sujet) pour fusionner deux phrases simples en une phrase complexe.}
\end{popfigure}

\courssub{Zoom sur \og dont \fg et \og où \fg : deux outils techniques indispensables}

\begin{aretenir}[Emploi de \textbf{dont} = \og de + antécédent\fg]
\emph{Dont} remplace un complément introduit par la préposition \textbf{de} (verbe + \emph{de}, nom + \emph{de}, adjectif + \emph{de}). C'est l'outil principal pour exprimer la \textbf{possession}, la \textbf{partie/quantité} ou l'\textbf{objet d'un verbe prépositionnel} sans répéter le nom.
\begin{itemize}
    \item \textbf{Verbe + de :} \og Les risques \textbf{dont} nous avons alerté la direction.\fg (alerté \emph{de} les risques)
    \item \textbf{Nom + de :} \og La conséquence \textbf{dont} nous craignons l'issue.\fg (l'issue \emph{de} la conséquence)
    \item \textbf{Quantité / Partie :} \og 12 quartiers ont été visités, \textbf{dont} 5 sont critiques.\fg (5 \emph{sur} les 12)
\end{itemize}
\end{aretenir}

\begin{aretenir}[Emploi de \textbf{où} = Lieu ou Temps]
\emph{Où} est un pronom relatif \textbf{invariable} qui remplace un complément circonstanciel de \textbf{lieu} (dans, à, sur, sous...) ou de \textbf{temps} (le jour où, l'année où, le moment où).
\begin{itemize}
    \item \textbf{Lieu :} \og Le bassin \textbf{où} les eaux stagnent est infecté.\fg (dans \emph{lequel} $\rightarrow$ \emph{où})
    \item \textbf{Temps :} \og La semaine \textbf{où} la mission a débuté était pluvieuse.\fg (pendant \emph{laquelle} $\rightarrow$ \emph{où})
\end{itemize}
\end{aretenir}

\begin{attention}[Erreurs fatales au Bac / Probatoire]
\begin{enumerate}
    \item \textbf{Confusion \og où \fg (relatif) / \og ou \fg (conjonction)} \\
    \emph{Où} (avec accent) = lieu/temps (\og Le site \textbf{où} nous allons\fg). \\
    \emph{Ou} (sans accent) = alternative (\og Acier \textbf{ou} béton ?\fg). \textbf{Astuce :} Remplacez par \emph{ou bien}. Si ça marche, c'est \emph{ou} (sans accent).
    \item \textbf{Confusion \og dont \fg / \og que \fg} \\
    \emph{Dont} = objet de \emph{de} (\og Le problème \textbf{dont} je parle\fg = je parle \emph{de} le problème). \\
    \emph{Que} = COD direct (\og Le problème \textbf{que} je résous\fg = je résous \emph{le problème}).
    \item \textbf{Oubli de l'antécédent} : \og *Voici le rapport dont est long.\fg (Incorrect). $\rightarrow$ \og Voici le rapport \textbf{qui} est long.\fg (Sujet) OU \og Voici le rapport \textbf{dont} la lecture est longue.\fg (Complément du nom \emph{lecture}).
\end{enumerate}
\end{attention}

\courssub{Méthode : Choisir le bon connecteur (adverbe ou pronom)}

\begin{methode}[Démarche en 3 étapes pour relier deux idées]
\textbf{Étape 1 : Analyser la relation logique entre l'idée A et l'idée B.}
\begin{itemize}
    \item B \emph{s'ajoute} à A $\rightarrow$ \textbf{Addition} (\emph{de plus, en outre}).
    \item B \emph{s'oppose} ou \emph{nuance} A $\rightarrow$ \textbf{Opposition} (\emph{cependant, toutefois}).
    \item B \emph{explique la cause} de A $\rightarrow$ \textbf{Cause} (\emph{en effet, car}).
    \item B \emph{est la conséquence} de A $\rightarrow$ \textbf{Conséquence} (\emph{donc, par conséquent}).
    \item B \emph{suit} A dans le temps $\rightarrow$ \textbf{Chronologie} (\emph{ensuite, enfin}).
\end{itemize}

\textbf{Étape 2 : Déterminer le niveau syntaxique.}
\begin{itemize}
    \item Pour lier \textbf{deux phrases distinctes} (ou deux gros blocs) $\rightarrow$ \textbf{Adverbe de liaison / Locution} (souvent en début de phrase, suivi d'une virgule).
    \item Pour lier \textbf{deux propositions} en une seule phrase complexe (éviter la répétition d'un nom) $\rightarrow$ \textbf{Pronom relatif} (\emph{qui, que, dont, où, lequel}).
\end{itemize}

\textbf{Étape 3 : Vérifier la grammaire (accord, préposition, fonction).}
\begin{itemize}
    \item \emph{Lequel} et composés : accord avec l'antécédent + préposition requise.
    \item \emph{Dont} : vérifier la présence de \emph{de} dans la phrase source.
    \item \emph{Où} : vérifier lieu/temps (pas personne).
    \item Ponctuation : virgule après l'adverbe de liaison en tête de phrase ; pas de virgule devant \emph{que/qui/dont/où} (sauf incise).
\end{itemize}
\end{methode}

\courssub{Applications : Exercices résolus}

\begin{exoresolu}[Exercice 1 : Fusion par pronoms relatifs]
\emph{Reprenez les paires de phrases suivantes en une seule phrase complexe en utilisant le pronom relatif approprié (qui, que, dont, où, auquel...).}

\textbf{1.} Le technicien a vérifié la vanne. La vanne était bloquée. \\
\textbf{2.} Nous avons mesuré le débit. Le débit est conforme aux normes. \\
\textbf{3.} L'ingénieur a signalé des anomalies. Nous n'avions pas prévu ces anomalies. \\
\textbf{4.} Le chantier est situé à Douala-Bassa. Les inondations y sont fréquentes. \\
\textbf{5.} Voici le rapport. La direction attend ce rapport.

\tcblower
\textbf{Solution détaillée :}

\textbf{1.} Antécédent : \emph{la vanne}. Fonction dans P2 : \textbf{Sujet} de \emph{était} $\rightarrow$ \textbf{qui}. \\
\og Le technicien a vérifié la vanne \textbf{qui} était bloquée.\fg

\textbf{2.} Antécédent : \emph{le débit}. Fonction dans P2 : \textbf{Sujet} (attribut du sujet) de \emph{est} $\rightarrow$ \textbf{qui}. \\
\og Nous avons mesuré le débit \textbf{qui} est conforme aux normes.\fg

\textbf{3.} Antécédent : \emph{des anomalies}. Fonction dans P2 : \textbf{COD} de \emph{avions prévu} (pas de préposition) $\rightarrow$ \textbf{que} (qu'). \\
\og L'ingénieur a signalé des anomalies \textbf{qu'} nous n'avions pas prévues.\fg (Accord du participe passé \emph{prévues} avec \emph{que} mis pour \emph{anomalies} placé avant).

\textbf{4.} Antécédent : \emph{Douala-Bassa} (lieu). Fonction : \textbf{CCL} (\emph{dans} Douala-Bassa) $\rightarrow$ \textbf{où}. \\
\og Le chantier est situé à Douala-Bassa, \textbf{où} les inondations sont fréquentes.\fg

\textbf{5.} Antécédent : \emph{le rapport}. Fonction dans P2 : \textbf{COD} de \emph{attend} $\rightarrow$ \textbf{que} (qu'). \\
\og Voici le rapport \textbf{que} la direction attend.\fg
\end{exoresolu}

\begin{exoresolu}[Exercice 2 : Texte à trous — Rapport de visite d'usine (Douala-Bassa)]
\emph{Complétez le rapport ci-dessous avec les adverbes de liaison ou locutions adverbiales proposés :}
\medskip
\og \textbf{\_\_\_\_\_\_\_\_}, l'équipe de maintenance s'est rendue à l'usine de traitement d'eau de Douala-Bassa. \textbf{\_\_\_\_\_\_\_\_}, elle a inspecté les bassins de décantation. \textbf{\_\_\_\_\_\_\_\_}, elle a contrôlé les pompes de refoulement. \textbf{\_\_\_\_\_\_\_\_}, les bassins présentent un niveau de boue critique. \textbf{\_\_\_\_\_\_\_\_}, les pompes fonctionnent normalement. \textbf{\_\_\_\_\_\_\_\_}, il est urgent de programmer une vidange. \textbf{\_\_\_\_\_\_\_\_}, un devis sera transmis à la direction. \textbf{\_\_\_\_\_\_\_\_}, la mission s'est terminée à 16h.\fg
\medskip
\textbf{Liste :} \emph{en premier lieu / ensuite / en second lieu / en effet / cependant / par conséquent / pour finir / enfin}

\tcblower
\textbf{Solution détaillée et justifications :}

\og \textbf{En premier lieu}, l'équipe de maintenance s'est rendue à l'usine de traitement d'eau de Douala-Bassa. \textbf{En second lieu}, elle a inspecté les bassins de décantation. \textbf{Ensuite}, elle a contrôlé les pompes de refoulement. \textbf{En effet}, les bassins présentent un niveau de boue critique. \textbf{Cependant}, les pompes fonctionnent normalement. \textbf{Par conséquent}, il est urgent de programmer une vidange. \textbf{Pour finir}, un devis sera transmis à la direction. \textbf{Enfin}, la mission s'est terminée à 16h.\fg

\begin{itemize}
    \item \textbf{En premier lieu / En second lieu / Ensuite} : \emph{Chronologie / Énumération} des actions successives (structure du récit).
    \item \textbf{En effet} : \emph{Cause / Explication} (justifie l'urgence : le niveau de boue critique explique la nécessité de vidange).
    \item \textbf{Cependant} : \emph{Opposition / Contraste} (nuance : les pompes vont bien, mais les bassins non).
    \item \textbf{Par conséquent} : \emph{Conséquence} (décision technique logique issue du constat).
    \item \textbf{Pour finir / Enfin} : \emph{Clôture} (administrative et temporelle).
\end{itemize}
\end{exoresolu}

\begin{savaistu}[Le critère « Cohérence et Cohésion » au Baccalauréat]
Au Probatoire et au Baccalauréat technique, la grille d'évaluation de la production écrite (dissertation, rapport, synthèse) comporte un item explicite : \textbf{« Cohérence du discours / Qualité des liaisons »} (souvent noté sur 2 à 4 points sur 20).
\begin{itemize}
    \item Un texte sans connecteurs (style \emph{télégrammatique}) est pénalisé : \og L'équipe arrive. L'équipe voit. L'équipe part.\fg
    \item Un texte avec connecteurs \textbf{mal choisis} (ex: \emph{Donc} pour une opposition) montre une incompréhension de la logique du propos.
    \item Un texte où les \textbf{pronoms relatifs} remplacent efficacement les répétitions lourdes (\emph{l'usine... l'usine... l'usine...}) gagne en lisibilité et en maturité syntaxique.
\end{itemize}
\textbf{Conseil :} Relisez toujours votre production en surlignant vos connecteurs. Demandez-vous : \og Ce mot exprime-t-il \emph{exactement} le lien logique que je veux créer ?\fg
\end{savaistu}

\begin{exercice}[Entraînement : Révision complète]
\textbf{A.} Reliez ces phrases par un pronom relatif (attention aux accords) :
\begin{enumerate}
    \item Le chef de mission a signé le PV. Le PV valide les travaux.
    \item Les techniciens ont nettoyé les filtres. Les filtres étaient colmatés.
    \item Nous avons visité le laboratoire. Les analyses s'y font chaque jour.
    \item Voici les échantillons. Le laboratoire a besoin de ces échantillons.
\end{enumerate}

\textbf{B.} Choisissez le connecteur logique adapté (adverbe ou locution) :
\begin{enumerate}
    \item Le budget est serré, \_\_\_\_\_\_\_\_ (donc / cependant / en outre) nous devons réduire les frais de déplacement.
    \item La machine est neuve. \_\_\_\_\_\_\_\_ (Cependant / En effet / Par ailleurs), elle tombe souvent en panne.
    \item \_\_\_\_\_\_\_\_ (En premier lieu / Par conséquent / D'autre part), nous définissons le cahier des charges. \_\_\_\_\_\_\_\_ (Ensuite / En définitive / En effet), nous lançons l'appel d'offres.
\end{enumerate}
\end{exercice}

\begin{corrige}[Exercice Entraînement]
\textbf{A.}
\begin{enumerate}
    \item Le chef de mission a signé le PV \textbf{qui} valide les travaux. (Sujet)
    \item Les techniciens ont nettoyé les filtres \textbf{qui} étaient colmatés. (Sujet de \emph{were} / \emph{étaient})
    \item Nous avons visité le laboratoire \textbf{où} les analyses se font chaque jour. (Lieu : \emph{dans} le laboratoire)
    \item Voici les échantillons \textbf{dont} le laboratoire a besoin. (Verbe \emph{avoir besoin de})
\end{enumerate}

\textbf{B.}
\begin{enumerate}
    \item \textbf{donc} (Conséquence : budget serré $\rightarrow$ réduction frais).
    \item \textbf{Cependant} (Opposition/Concession : neuf $\neq$ tombe en panne). \emph{En effet} aurait signifié que la panne est \emph{due} à la nouveauté (absurde). \emph{Par ailleurs} changerait de sujet.
    \item \textbf{En premier lieu} (Chronologie 1). \textbf{Ensuite} (Chronologie 2). \emph{En définitive} conviendrait pour une conclusion finale, pas pour une étape intermédiaire.
\end{enumerate}
\end{corrige}

\coursec{Rédiger l'introduction et la conclusion}

Dans la section précédente, nous avons étudié les \cle{outils de liaison} dans la phrase : adverbes de liaison, locutions adverbiales et pronoms relatifs. Ces outils ne sont pas de simples ornements grammaticaux : ce sont eux qui permettent d'enchaîner harmonieusement les idées dans un texte professionnel. Nous allons maintenant les mettre immédiatement en pratique dans les deux moments les plus stratégiques d'un rapport : son \cle{introduction} et sa \cle{conclusion}. En effet, dans la vie administrative et professionnelle, un destinataire pressé lit souvent d'abord l'introduction pour savoir de quoi il retourne, puis la conclusion pour savoir ce qu'on lui propose. Si ces deux passages sont ratés, tout le rapport est discrédité.

\courssub{Fonction et éléments obligatoires de l'introduction}

\begin{definition}[L'introduction du rapport]
L'\cle{introduction} d'un rapport est la partie liminaire qui \textbf{situe le contexte} de la mission, \textbf{rappelle la mission reçue}, \textbf{précise la méthode} employée et \textbf{annonce le plan} du développement. Elle répond aux questions : \emph{qui ? quand ? où ? pourquoi ? comment ?}
\end{definition}

Pour bien comprendre, imaginons une situation concrète. Le maire de Bafoussam reçoit sur son bureau un document de trois pages. Sa première question est : « D'où vient ce document ? Qui l'a écrit ? Sur quel ordre ? À quel sujet ? » Si l'introduction ne répond pas immédiatement à ces questions, le maire est perdu. L'introduction est donc la \emph{carte d'identité} de la mission.

\begin{propriete}[Les cinq éléments obligatoires de l'introduction]
Toute introduction de rapport doit contenir :
\begin{enumerate}
\item \textbf{Qui} : l'auteur du rapport (nom, fonction) et son \cle{mandant} (la personne ou l'autorité qui a ordonné la mission) ;
\item \textbf{Quand} : la date ou la période de la mission ;
\item \textbf{Où} : le lieu de la mission ;
\item \textbf{Pourquoi} : l'objet de la mission (le problème à étudier) ;
\item \textbf{Comment} : la méthode utilisée (enquête, observation, entretiens, consultation de documents).
\end{enumerate}
\end{propriete}

\begin{methode}[Rédiger une introduction de rapport en quatre étapes]
\begin{enumerate}
\item \textbf{Contexte} : rappeler brièvement la situation qui a motivé la mission (un problème, une plainte, une décision administrative).
\item \textbf{Mission} : indiquer qui a donné l'ordre, à qui, et pour faire quoi, en citant la référence de la lettre de mission si elle existe.
\item \textbf{Méthode} : préciser en une phrase comment les informations ont été recueillies (visites, entretiens, questionnaires, observation).
\item \textbf{Annonce du plan} : présenter les grandes parties du rapport avec des connecteurs d'énumération : \emph{d'abord, ensuite, enfin}.
\end{enumerate}
\end{methode}

\begin{exemplebox}[Formules types d'introduction]
\begin{itemize}
\item Ouverture : « Conformément à votre lettre n\textsuperscript{o} 045/CAB/PR du 12 mars 2026, j'ai effectué une mission à Ndogpassi afin d'évaluer l'ampleur des dégâts causés par les inondations. »
\item Méthode : « Pour mener à bien cette mission, nous avons procédé par observation directe des lieux, puis par des entretiens avec les chefs de quartier et les sinistrés. »
\item Annonce du plan : « Ce rapport présente d'abord l'état des lieux, ensuite les causes identifiées, enfin les mesures envisageables. »
\end{itemize}
\end{exemplebox}

\courssub{Fonction et éléments obligatoires de la conclusion}

\begin{definition}[La conclusion du rapport]
La \cle{conclusion} d'un rapport est la partie finale qui \textbf{résume les constats essentiels} de la mission et \textbf{formule des suggestions ou recommandations} à l'attention du destinataire. Elle ne contient \emph{aucun fait nouveau}.
\end{definition}

Reprenons notre maire de Bafoussam. Après avoir parcouru le corps du rapport, il se demande : « Au final, que faut-il retenir ? Et surtout, que me propose-t-on de faire ? » La conclusion répond à ces deux questions. Elle a donc deux temps : le \cle{bilan} (ce qui ressort de la mission) et les \cle{suggestions} (ce qu'il convient de faire).

\begin{methode}[Rédiger une conclusion de rapport]
\begin{enumerate}
\item Rédiger un \textbf{bilan synthétique} en deux ou trois phrases, qui reprend les constats majeurs sans répéter les détails du développement.
\item Formuler \textbf{deux à quatre suggestions concrètes}, hiérarchisées de la plus urgente à la moins urgente.
\item Employer des formulations claires et mesurées : infinitif (« procéder à... »), futur (« nous installerons... ») ou conditionnel (« il conviendrait de... », « nous recommandons que... » suivi du subjonctif).
\end{enumerate}
\end{methode}

\begin{exemplebox}[Formules types de conclusion]
\begin{itemize}
\item Bilan : « En définitive, il ressort de cette mission que la fréquentation de la bibliothèque municipale a diminué de moitié en trois ans. »
\item Recommandation : « C'est pourquoi nous suggérons de réhabiliter la salle de lecture et d'organiser des animations mensuelles. »
\item Autre tournure : « Il conviendrait également que la mairie augmente le budget d'acquisition des ouvrages. »
\end{itemize}
\end{exemplebox}

\begin{attention}[Le piège de la conclusion]
L'erreur la plus fréquente consiste à \textbf{introduire des faits nouveaux dans la conclusion} : un chiffre qui n'apparaît nulle part ailleurs, une cause découverte « au dernier moment », un témoignage inédit. C'est une faute grave de méthode : la conclusion ne fait que \emph{synthétiser} ce qui a été démontré dans le corps du rapport et \emph{proposer} des solutions. Tout fait nouveau doit être replacé dans le développement. Deuxième piège : des suggestions vagues (« il faut faire quelque chose ») au lieu de propositions précises et réalisables.
\end{attention}

\begin{popfigure}
\begin{center}
\begin{tikzpicture}[
  bloc/.style={rectangle, rounded corners=4pt, draw=popblue, fill=popblueL, text width=4.6cm, align=center, minimum height=1.1cm, font=\small},
  titre/.style={rectangle, rounded corners=4pt, draw=popdark, fill=popdark, text=white, align=center, font=\bfseries\small, minimum height=0.9cm},
  egal/.style={font=\Large\bfseries, text=poporange}
]
\node[titre, text width=6.4cm] (ti) at (0,0) {INTRODUCTION};
\node[bloc, align=center] (b1) at (-2.4,-1.6){\textbf{Contexte}\\ situation, problème};
\node[bloc, align=center] (b2) at (2.4,-1.6){\textbf{Mission}\\ qui, quand, où, pourquoi};
\node[bloc, align=center] (b3) at (-2.4,-3.1){\textbf{Méthode}\\ enquête, observation, entretiens};
\node[bloc, align=center] (b4) at (2.4,-3.1){\textbf{Plan}\\ d'abord, ensuite, enfin};
\node[egal] at (0,-2.35) {=};
\draw[popblue, thick] (ti.south) -- (0,-0.9) -| (b1.north);
\draw[popblue, thick] (0,-0.9) -| (b2.north);
\draw[popblue, thick] (0,-0.9) |- (-2.4,-2.55) -- (b3.north);
\draw[popblue, thick] (0,-0.9) |- (2.4,-2.55) -- (b4.north);

\node[titre, text width=6.4cm, fill=popgreen, draw=popgreen] (tc) at (0,-4.9) {CONCLUSION};
\node[bloc, draw=popgreen, fill=popgreenL, align=center] (c1) at (-2.4,-6.5){\textbf{Bilan}\\ constats essentiels (2-3 phrases)};
\node[bloc, draw=popgreen, fill=popgreenL, align=center] (c2) at (2.4,-6.5){\textbf{Suggestions}\\ 2 à 4 recommandations concrètes};
\node[egal] at (0,-6.5) {=};
\draw[popgreen, thick] (tc.south) -- (0,-5.8) -| (c1.north);
\draw[popgreen, thick] (0,-5.8) -| (c2.north);
\end{tikzpicture}
\end{center}
\legende{Schéma fonctionnel des deux extrémités du rapport : l'introduction combine quatre éléments (contexte, mission, méthode, plan) ; la conclusion en combine deux (bilan, suggestions).}
\end{popfigure}

\courssub{Exemple commenté : l'introduction d'un rapport sur les inondations de Ndogpassi}

\begin{exemplebox}[Introduction rédigée et commentée phrase par phrase]
\textbf{Texte de l'introduction :}

« Les pluies diluviennes du mois d'octobre 2025 ont provoqué de graves inondations dans le quartier Ndogpassi, à Douala, faisant de nombreux sans-abri. \textbf{(1)} Conformément à votre lettre n\textsuperscript{o} 112/G/DLA du 25 octobre 2025, nous avons effectué, du 27 au 30 octobre 2025, une mission d'évaluation dans ce quartier. \textbf{(2)} Pour recueillir les informations, nous avons visité les zones sinistrées, interrogé une centaine de ménages et rencontré le chef du quartier. \textbf{(3)} Ce rapport présente d'abord l'étendue des dégâts, ensuite les causes de ces inondations, enfin les solutions que nous préconisons. \textbf{(4)} »

\textbf{Commentaire :}
\begin{itemize}
\item \textbf{(1)} Contexte : la catastrophe est située dans le temps, dans l'espace, et ses conséquences humaines sont annoncées.
\item \textbf{(2)} Mission : le mandant est désigné par la référence exacte de sa lettre ; les dates de la mission sont précises.
\item \textbf{(3)} Méthode : trois procédés concrets (visite, questionnaire, entretien) garantissent la fiabilité des constats.
\item \textbf{(4)} Annonce du plan : les connecteurs \emph{d'abord, ensuite, enfin} — outils de liaison étudiés dans la section précédente — ordonnent les trois parties du développement.
\end{itemize}
\end{exemplebox}

\begin{savaistu}[Le rapport administratif au Cameroun]
Dans l'administration camerounaise, un rapport sans suggestions est souvent jugé \emph{incomplet} par le destinataire. Un chef de service, un préfet ou un directeur n'attend pas seulement un constat : il attend qu'on l'aide à \textbf{décider}. C'est pourquoi la rubrique « suggestions » ou « recommandations » est devenue une convention quasi obligatoire des rapports de mission, des procès-verbaux d'inspection et des comptes rendus ministériels. Savoir la rédiger avec soin est donc un atout professionnel décisif pour tout technicien.
\end{savaistu}

\courssub{Atelier guidé : rapport sur la fréquentation de la bibliothèque municipale de Bafoussam}

\begin{exoresolu}[Rédaction guidée d'une introduction et d'une conclusion]
\textbf{Énoncé.} Le délégué communal à la Culture de Bafoussam vous a chargé, par lettre n\textsuperscript{o} 018/DCC/BFM du 5 janvier 2026, d'étudier la baisse de fréquentation de la bibliothèque municipale. Vous avez mené votre enquête du 10 au 20 janvier 2026 (observation des lieux, questionnaire auprès de 80 usagers, entretien avec la bibliothécaire). Vos constats : locaux vétustes, fonds documentaire vieilli, absence d'animations. Rédigez l'introduction et la conclusion du rapport.
\tcblower
\textbf{Solution détaillée.}

\emph{Introduction} (contexte $\rightarrow$ mission $\rightarrow$ méthode $\rightarrow$ plan) :

« La bibliothèque municipale de Bafoussam, autrefois très fréquentée par les élèves et les étudiants, enregistre depuis trois ans une baisse constante de ses usagers. Conformément à votre lettre n\textsuperscript{o} 018/DCC/BFM du 5 janvier 2026, nous avons effectué, du 10 au 20 janvier 2026, une mission d'étude dans cet établissement afin d'identifier les causes de cette désaffection. À cet effet, nous avons observé l'état des locaux, soumis un questionnaire à quatre-vingts usagers et recueilli le témoignage de la bibliothécaire en chef. Ce rapport présente d'abord l'état des lieux, ensuite les causes de la baisse de fréquentation, enfin les mesures susceptibles d'y remédier. »

\emph{Conclusion} (bilan $\rightarrow$ suggestions) :

« En définitive, il ressort de cette mission que la baisse de fréquentation de la bibliothèque municipale s'explique principalement par la vétusté des locaux, le vieillissement du fonds documentaire et l'absence totale d'animations culturelles. C'est pourquoi nous suggérons, premièrement, de rénover la salle de lecture et d'y installer un éclairage adapté ; deuxièmement, d'acquérir des ouvrages récents conformes aux programmes scolaires ; troisièmement, d'organiser chaque mois des clubs de lecture et des rencontres avec des écrivains. Il conviendrait enfin que la délégation communale dote l'établissement d'une connexion Internet. Nous restons convaincus que ces mesures ramèneront les usagers vers ce cadre de savoir. »

\emph{Vérification} : l'introduction contient bien les cinq éléments (qui, quand, où, pourquoi, comment) ; la conclusion ne contient aucun fait nouveau et formule quatre suggestions concrètes à l'infinitif et au conditionnel.
\end{exoresolu}

\begin{exercice}[À vous de rédiger]
Le proviseur de votre lycée vous a chargé, par note de service n\textsuperscript{o} 023/NS/LT du 2 février 2026, d'enquêter sur le retard croissant des élèves au premier cours. Vous avez observé les entrées pendant une semaine et interrogé cinquante élèves et trois surveillants. Rédigez l'introduction (quatre étapes obligatoires) et la conclusion (bilan de deux phrases, puis trois suggestions) du rapport que vous lui remettrez.
\end{exercice}

\begin{corrige}[Exercice — éléments de réponse]
\emph{Introduction attendue} : contexte (recrudescence des retards constatée au lycée) ; mission (référence de la note de service, auteur, dates, lieu) ; méthode (observation des entrées, questionnaire auprès de cinquante élèves, entretiens avec trois surveillants) ; annonce du plan avec \emph{d'abord, ensuite, enfin} (ampleur du phénomène, causes, solutions).

\emph{Conclusion attendue} : bilan synthétique (les retards proviennent essentiellement des difficultés de transport, de l'éloignement des domiciles et du coucher tardif) ; trois suggestions concrètes, par exemple : sensibiliser les élèves et les parents à l'heure du coucher ; négocier avec les transporteurs urbains un passage plus tôt le matin ; instaurer un contrôle systématique à la barrière dès 7 h 15. Aucun fait nouveau ne doit apparaître ; les formulations doivent employer l'infinitif ou le conditionnel (« il conviendrait de... »).
\end{corrige}

\begin{aretenir}[L'essentiel]
\begin{itemize}
\item L'introduction répond à cinq questions : \textbf{qui, quand, où, pourquoi, comment}, puis annonce le plan.
\item La conclusion fait le \textbf{bilan} des constats et formule \textbf{deux à quatre suggestions} concrètes, sans jamais introduire de fait nouveau.
\item Les formules types (« Conformément à votre lettre n\textsuperscript{o}... », « En définitive, il ressort que... ») donnent au rapport son ton administratif.
\item Les outils de liaison (\emph{d'abord, ensuite, enfin, en définitive, c'est pourquoi}) assurent la cohérence de ces deux passages.
\end{itemize}
\end{aretenir}

\coursec{Rédiger le corps du rapport et exploiter les documents iconiques}

Dans la section précédente, nous avons appris à rédiger l'introduction (qui annonce le contexte, l'objectif et le plan) et la conclusion (qui propose des recommandations). Entre ces deux moments se trouve la partie la plus longue et la plus importante du rapport : le \cle{corps}. C'est là que le rapporteur présente les faits, les chiffres et les observations. Or, dans la vie professionnelle, ces informations ne viennent pas seulement de textes écrits : elles proviennent aussi de \cle{documents iconiques} — photographies, cartes, tableaux statistiques, diagrammes. Cette section vous apprend donc deux compétences indissociables : organiser le corps du rapport et transformer une image ou un tableau en paragraphes rédigés.

\courssub{Organiser le corps du rapport}

\textbf{Classer les faits.} Un corps de rapport n'est pas un amas d'informations jetées en désordre. Les faits doivent être \emph{classés} selon un ordre que le lecteur peut suivre. Trois classements sont possibles :

\begin{itemize}
\item le classement \cle{chronologique} : on présente les faits dans l'ordre où ils se sont produits (utile pour un rapport d'accident ou d'événement) ;
\item le classement \cle{thématique} : on regroupe les faits par thèmes (constats, causes, conséquences, solutions) ;
\item le classement \cle{par ordre d'importance} : on va du fait le plus important au moins important (utile quand le destinataire est pressé).
\end{itemize}

\textbf{Numéroter les parties.} Le corps se divise en parties numérotées (\textbf{1., 2., 3.}), elles-mêmes subdivisées en sous-parties (\textbf{1.1, 1.2, 2.1}). Cette numérotation permet au lecteur de se repérer et de citer le rapport : « voir point 2.3 ».

\textbf{Un paragraphe = une idée.} Règle d'or de la rédaction : chaque paragraphe développe \emph{une seule idée principale}, précisée par des exemples, des chiffres ou des témoignages. Si vous changez d'idée, changez de paragraphe.

\begin{methode}[Rédiger un paragraphe de rapport]
\begin{enumerate}
\item Formuler l'\textbf{idée principale} en une phrase simple (le constat).
\item Énoncer le \textbf{fait précis} qui l'illustre (date, lieu, événement).
\item Ajouter une \textbf{donnée chiffrée} ou un \textbf{témoignage} cité entre guillemets.
\item Conclure par une \textbf{courte analyse} neutre (conséquence, comparaison), sans opinion personnelle.
\end{enumerate}
\end{methode}

\begin{attention}[Constat ou opinion ?]
L'erreur la plus fréquente est de mélanger constat et opinion. Le rapport rapporte des \cle{faits vérifiables}, pas des sentiments. On bannit donc « je pense que », « je crois que », « à mon avis ». Au lieu d'écrire « \emph{Je pense que le drainage est mauvais} », on écrit : « \emph{Les caniveaux du quartier sont obstrués par des déchets plastiques, comme le montre la photographie n° 2.} » Le fait remplace l'opinion ; la preuve remplace l'impression.
\end{attention}

\courssub{Lire les textes iconiques}

\begin{definition}[Texte iconique]
Un \cle{texte iconique} est un document visuel — photographie, carte, schéma, tableau statistique, diagramme — porteur d'informations que le lecteur doit \emph{décrire} puis \emph{interpréter}. Dans un rapport, il sert de \textbf{source} et de \textbf{preuve}.
\end{definition}

\textbf{Décrire une photographie.} La description suit toujours le même mouvement : \emph{du général au particulier}. On commence par l'ensemble de la scène (où ? quoi ?), puis on détaille les éléments du premier plan, puis l'arrière-plan. On décrit ce qu'on \emph{voit}, sans interprétation hâtive : on écrit « l'eau arrive à mi-jambes des vendeuses » et non « les vendeuses sont désespérées » (le désespoir ne se voit pas, il se devine).

\textbf{Exploiter un tableau de données.} Face à un tableau statistique, la démarche comporte trois gestes : \emph{relever} les chiffres clés (maximum, minimum), \emph{comparer} (entre les jours, les lieux, les années), puis \emph{dégager une tendance} (hausse, baisse, pic). Ces chiffres doivent ensuite être intégrés dans des phrases rédigées : un rapport ne colle jamais un tableau brut sans commentaire.

\begin{exemplebox}[Intégrer les chiffres dans une phrase]
À partir du tableau des pluies, on ne se contente pas d'écrire « 12 septembre : 85 mm ». On rédige : « \emph{Le 12 septembre, il est tombé 85 mm de pluie en 6 heures, soit environ le triple de la moyenne mensuelle (28 mm).} » De même : « \emph{Le niveau des eaux a atteint 1,20 m, ce qui a contraint 350 ménages à se déplacer.} » La donnée chiffrée est insérée dans une phrase complète, souvent grâce à un outil de liaison (\emph{soit}, \emph{ce qui}) qui en tire la conséquence.
\end{exemplebox}

\begin{popfigure}
\begin{center}
\begin{tikzpicture}
\begin{axis}[
ybar, bar width=0.55cm,
width=0.85\linewidth, height=6cm,
xlabel={Jour (septembre)}, ylabel={Hauteur de pluie (mm)},
symbolic x coords={10,11,12,13,14},
xtick=data,
ymin=0, ymax=100,
nodes near coords,
ymajorgrids=true, grid style={popline!40},
every axis plot/.append style={fill=popblue, draw=popdark}
]
\addplot coordinates {(10,20) (11,45) (12,85) (13,60) (14,15)};
\end{axis}
\end{tikzpicture}
\end{center}
\legende{Hauteurs de pluie quotidiennes (en mm) relevées à Douala du 10 au 14 septembre. Le pic du 12 septembre (85 mm) constitue le chiffre clé à exploiter dans le corps du rapport.}
\end{popfigure}

\courssub{De l'image au paragraphe : la démarche}

Passer d'une photographie à un paragraphe de rapport n'est pas un don : c'est une procédure en quatre étapes que le schéma suivant résume.

\begin{popfigure}
\begin{center}
\begin{tikzpicture}[
node distance=0.5cm,
etape/.style={rectangle, rounded corners, draw=popdark, fill=popblueL, text width=2.6cm, minimum height=1.6cm, align=center, font=\small},
fleche/.style={-{Stealth[length=3mm]}, thick, poporange}
]
\node[etape] (a) {\textbf{1. Observer}\par Repérer lieu, personnes, objets};
\node[etape, right=of a] (b) {\textbf{2. Décrire}\par Du général au particulier};
\node[etape, right=of b] (c) {\textbf{3. Chiffrer}\par Ajouter mesures et données};
\node[etape, right=of c] (d) {\textbf{4. Analyser}\par Dégager la conséquence};
\draw[fleche] (a) -- (b);
\draw[fleche] (b) -- (c);
\draw[fleche] (c) -- (d);
\end{tikzpicture}
\end{center}
\legende{Procédure en quatre étapes pour transformer une photographie en paragraphe descriptif de rapport : observer, décrire, chiffrer, analyser.}
\end{popfigure}

\begin{savaistu}[Les bulletins météorologiques au Cameroun]
Les services météorologiques du Cameroun publient régulièrement des bulletins de prévision et des relevés de pluies. Ces documents servent de sources officielles aux rapports administratifs rédigés après les catastrophes naturelles (inondations de Douala, glissements de terrain dans la région de l'Ouest). Un bon rapporteur cite toujours ses sources : c'est ce qui rend ses chiffres vérifiables.
\end{savaistu}

\courssub{Activité guidée : rédiger le corps d'un rapport d'inondation}

\begin{experience}[TP — Du document iconique au corps de rapport]
\textbf{Matériel (documents fournis)} : une photographie d'un marché inondé à Douala (eau à mi-jambes, étals renversés, vendeuses déplaçant leurs marchandises) ; le graphique en barres des pluies du 10 au 14 septembre ci-dessus ; le témoignage d'une vendeuse : « L'eau est montée en une heure, nous avons tout perdu. »

\textbf{Protocole} :
\begin{enumerate}
\item Observer la photographie et noter cinq éléments visibles, du général au particulier.
\item Relever dans le graphique le jour du pic de pluie et sa hauteur.
\item Comparer ce pic à la moyenne mensuelle (28 mm) et calculer le rapport : $85 \div 28 \approx 3$.
\item Rédiger le corps du rapport en deux parties numérotées : \textbf{1. Les constats} ; \textbf{2. Les causes}.
\end{enumerate}

\textbf{Observations attendues} : le pic (85 mm le 12) correspond à la date de la photographie ; la montée rapide des eaux s'explique par l'intensité exceptionnelle de la pluie et par l'obstruction des caniveaux visible sur le cliché.

\textbf{Interprétation} : croiser la photographie (preuve visuelle) et le graphique (preuve chiffrée) permet d'établir un lien de cause à effet que chaque document, pris isolément, ne suffirait pas à démontrer.
\end{experience}

\begin{exoresolu}[Rédaction du corps du rapport]
\textbf{Énoncé.} À partir des documents du TP, rédigez le corps du rapport en deux parties (constats, causes), en intégrant au moins deux données chiffrées et un témoignage.
\tcblower
\textbf{Solution détaillée.}

\textbf{1. Les constats.}

1.1. Le mardi 12 septembre, le marché a été envahi par les eaux. La photographie prise à 14 heures montre l'eau arrivant à mi-jambes des vendeuses, plusieurs étals renversés et des marchandises flottant entre les allées.

1.2. Le niveau des eaux a atteint 1,20 m en début de soirée, ce qui a contraint 350 ménages du quartier à se déplacer. Selon une vendeuse interrogée sur place, « l'eau est montée en une heure, nous avons tout perdu ».

\textbf{2. Les causes.}

2.1. Une pluie exceptionnelle. Le 12 septembre, il est tombé 85 mm de pluie en 6 heures, soit environ le triple de la moyenne mensuelle (28 mm). Le graphique des relevés confirme que ce jour constitue un pic nettement supérieur aux journées voisines (45 mm le 11, 60 mm le 13).

2.2. Un drainage insuffisant. La photographie montre en outre des caniveaux obstrués par des déchets plastiques, ce qui a empêché l'évacuation normale des eaux de ruissellement.

\emph{Remarque méthodologique} : chaque paragraphe suit le schéma « idée principale → fait précis → donnée chiffrée ou témoignage → courte analyse » ; aucune opinion personnelle n'apparaît.
\end{exoresolu}

\begin{exercice}[À vous de rédiger]
Votre chef de service vous envoie constater l'état d'une route dégradée après les pluies. Vous disposez d'une photographie (nid-de-poule de 2 m de large, véhicules immobilisés) et d'un tableau indiquant 120 mm de pluie en trois jours. Rédigez le corps du rapport en deux parties numérotées (1. constats ; 2. causes), en appliquant la procédure « observer → décrire → chiffrer → analyser » et en intégrant les deux données chiffrées dans des phrases complètes.
\end{exercice}

\begin{corrige}[Exercice — éléments attendus]
\textbf{1. Les constats.} 1.1. La photographie montre un nid-de-poule d'environ 2 m de large au kilomètre 4 de la route, ainsi que plusieurs véhicules immobilisés. 1.2. La circulation est interrompue depuis deux jours, ce qui isole trois villages (donnée à adapter). \textbf{2. Les causes.} 2.1. Il est tombé 120 mm de pluie en trois jours, soit un cumul exceptionnel qui a saturé le sol. 2.2. L'absence de fossés d'évacuation, visible sur le cliché, a concentré l'eau sur la chaussée. Le correctif valorise : la numérotation, l'objectivité (aucun « je pense »), l'intégration des chiffres dans des phrases et le lien entre la photographie et le tableau.
\end{corrige}

\begin{aretenir}[L'essentiel de la section]
Le corps du rapport classe les faits (ordre chronologique, thématique ou par importance) et les numérote (1., 1.1, 1.2). Chaque paragraphe développe une idée principale appuyée par un fait, un chiffre ou un témoignage. Les documents iconiques (photographies, tableaux, diagrammes) sont des sources et des preuves : on les décrit du général au particulier, on relève les chiffres clés, on compare, on dégage une tendance, puis on intègre ces données dans des phrases rédigées. Le rapporteur reste objectif : il rapporte des faits vérifiables, jamais ses opinions.
\end{aretenir}

\coursec{Produire et améliorer un rapport complet}

Après avoir maîtrisé la rédaction du corps du rapport et l'exploitation des documents iconiques (section précédente), nous entrons dans la phase de synthèse : produire un rapport complet, de la commande à la version finale, et développer les réflexes d'auto-évaluation indispensables à l'examen. C'est ici que toutes les compétences de l'unité — structure, objectivité, liaison, présentation — se conjuguent pour répondre à une situation de communication professionnelle réelle.

\courssub{Les cinq étapes de la production d'un rapport}

Tout rapport efficace naît d'un processus rigoureux en cinq temps. Négliger l'une de ces étapes expose à l'incohérence, à l'oubli d'informations capitales ou à des fautes de présentation qui décrédibilisent le rédacteur.

\begin{popfigure}
\begin{tikzpicture}[
    node distance=2.2cm,
    stepbox/.style={circle, draw=popblue, thick, fill=popblueL, minimum width=2.8cm, align=center, font=\small\bfseries},
    arrow/.style={-Stealth, thick, popdark},
    labelnode/.style={font=\footnotesize, text width=2.5cm, align=center, popmuted}
]
% Nœuds principaux (disposition en cercle)
\node[stepbox, align=center] (a){1. Analyser\\la situation};
\node[stepbox, right=of a, align=center] (b){2. Collecter\\les informations};
\node[stepbox, below right=of b, align=center] (c){3. Planifier\\la structure};
\node[stepbox, below left=of c, align=center] (d){4. Rédiger\\le texte};
\node[stepbox, left=of d, align=center] (e){5. Améliorer\\et relire};
% Flèches
\draw[arrow] (a) -- (b);
\draw[arrow] (b) -- (c);
\draw[arrow] (c) -- (d);
\draw[arrow] (d) -- (e);
\draw[arrow] (e) to[out=135, in=-135, looseness=1.5] (a);
% Étiquettes explicatives
\node[labelnode, above=0.2cm of a, align=center]{\textit{Qui ? Quoi ? Pour qui ?}\\\textit{Contexte, destinataire, objet}};
\node[labelnode, above=0.2cm of b, align=center]{\textit{Observation, enquête,}\\\textit{documents, témoignages}};
\node[labelnode, align=center, right=0.2cm of c] {\textit{Plan détaillé :\\intro, parties, conclusion}};
\node[labelnode, align=center, below=0.2cm of d] {\textit{Respect du plan,\\objectivité, connecteurs}};
\node[labelnode, align=center, left=0.2cm of e] {\textit{3 relectures :\\structure, langue, présentation}};
\end{tikzpicture}
\legende{Les cinq étapes itératives de la production d'un rapport. La flèche de retour symbolise la révision : on peut revenir au plan ou à la collecte si la relecture révèle des lacunes.}
\end{popfigure}

\begin{definition}[Analyse de la situation de communication]
Avant d'écrire, on identifie les \cle{cinq paramètres} de la situation :
\begin{itemize}
    \item \textbf{L'émetteur} (qui écrit ? quelle est sa fonction, sa légitimité ?).
    \item \textbf{Le destinataire} (à qui s'adresse le rapport ? quelle est sa hiérarchie, son niveau d'information ?).
    \item \textbf{L'objet} (sujet précis, circonscrit, formulé en groupe nominal).
    \item \textbf{La finalité} (informer, alerter, proposer, justifier ?).
    \item \textbf{Le contexte} (circonstances, urgence, cadre institutionnel).
\end{itemize}
\end{definition}

\begin{methode}[Analyser la commande en 3 questions]
\begin{enumerate}
    \item \textbf{Qui suis-je ?} (Ex. : Délégué de la classe de Tle G2).
    \item \textbf{À qui m'adresse-je ?} (Ex. : Monsieur le Proviseur du Lycée Technique de...).
    \item \textbf{Quel est l'objet exact et la finalité ?} (Ex. : Rapport sur les conditions de passage du Bac blanc / Alerter sur les dysfonctionnements et proposer des améliorations).
\end{enumerate}
\end{methode}

\courssub{Rédaction guidée : rapport du délégué de classe au proviseur}

Appliquons la méthode à un sujet d'entraînement type Bac : \emph{« En votre qualité de délégué de la classe de Terminale G2, vous rendez compte au Proviseur des conditions de passage du Baccalauréat blanc. Vous signalez les difficultés rencontrées et formulez des propositions. »}

\begin{exoresolu}[Sujet : Rapport sur les conditions du Bac blanc]
\textbf{Étape 1 : Analyse de la situation}
\begin{itemize}
    \item \textbf{Émetteur :} Délégué de la classe de Tle G2 (représentant légitime des élèves).
    \item \textbf{Destinataire :} Monsieur le Proviseur (autorité décisionnaire, ton respectueux, formel).
    \item \textbf{Objet :} Conditions de déroulement de l'examen blanc (du 12 au 16 mars 2026).
    \item \textbf{Finalité :} Informer objectivement, alerter sur des anomalies, proposer des solutions concrètes.
\end{itemize}

\textbf{Étape 2 : Collecte des informations (simulée)}
\begin{itemize}
    \item Horaires : Début des épreuves à 8h respecté, mais fin tardive (18h30) sans pause méridienne suffisante.
    \item Locaux : Salle 12 surchauffée (35°C), bruit de la cour de récréation mitoyenne.
    \item Matériel : 3 calculettes défaillantes en Mathématiques, sujets d'Anglais mal photocopiés (illisibles).
    \item Encadrement : Surveillants insuffisants (1 pour 40), entrées/sorties non contrôlées.
    \item Points positifs : Mise à disposition d'eau minérale, infirmerie opérationnelle.
\end{itemize}

\textbf{Étape 3 : Planification (Structure interne)}
\begin{itemize}
    \item \textbf{Introduction} (4 éléments : accroche, contexte, objet, annonce du plan).
    \item \textbf{Corps} :
        \begin{itemize}
            \item Partie 1 : Aspects positifs (organisation générale, moyens mis en œuvre).
            \item Partie 2 : Dysfonctionnements matériels et environnementaux (locaux, matériel, bruit).
            \item Partie 3 : Dysfonctionnements organisationnels (horaires, encadrement, pauses).
        \end{itemize}
    \item \textbf{Conclusion} : Bilan synthétique, 3 propositions concrètes, formule de politesse.
\end{itemize}

\textbf{Étape 4 : Rédaction (Extrait — Introduction et Début du corps)}

\begin{quote}
\textbf{RAPPORT N° 04/2026}

\textbf{Objet :} Conditions de déroulement du Baccalauréat blanc — Session de mars 2026 — Classe de Terminale G2.

\textbf{De :} \textit{Nom Prénom}, Délégué de la classe de Terminale G2.

\textbf{À :} Monsieur le Proviseur du Lycée Technique de [Nom de l'établissement].

\textbf{Date :} 18 mars 2026.

\textbf{Introduction}
Monsieur le Proviseur,
La session du Baccalauréat blanc, organisée du 12 au 16 mars 2026, constitue une étape décisive dans la préparation de nos examens officiels. En ma qualité de délégué de la classe de Terminale G2, mandaté par mes camarades pour relayer leurs observations, j'ai l'honneur de vous présenter ce rapport sur les conditions matérielles et organisationnelles de ces épreuves. Ce document fait d'abord état des dispositions satisfaisantes, avant de signaler les difficultés rencontrées et de formuler des propositions d'amélioration pour les sessions ultérieures.

\textbf{1. Des dispositions globales satisfaisantes}
Premièrement, l'organisation générale a fait preuve de rigueur : le calendrier a été communiqué à l'avance et les horaires de début d'épreuves (8h00) ont été scrupuleusement respectés. Par ailleurs, la mise à disposition quotidienne de bouteilles d'eau dans les salles et la présence effective de l'infirmière scolaire ont été très appréciées par l'ensemble des candidats. Ces mesures témoignent de la sollicitude de l'administration pour le bien-être des élèves.
\end{quote}

\textbf{Étape 5 : Amélioration (voir grille et méthode ci-après).}
\end{exoresolu}

\courssub{La grille d'auto-évaluation : un outil pour le Bac}

Cette grille doit être utilisée \textbf{systématiquement} lors de la relecture (étape 5). Elle couvre les 4 critères de notation du Baccalauréat technique.

\begin{popfigure}
\begin{center}
\begin{tabularx}{\linewidth}{|>{\bfseries}l|X|c|c|c|c|}\hline
\rowcolor{popblueL}
\textbf{Critère d'évaluation} & \textbf{Indicateurs de réussite (Me suis-je assuré que... ?)} & \textbf{Oui} & \textbf{Non} & \textbf{À corriger} & \textbf{Pts/5} \\ \hline
\textbf{1. Structure externe complète} & En-tête (Expéditeur, Destinataire, Date, Objet, Référence) présent et exact ? & $\square$ & $\square$ & & /5 \\ \hline
\textbf{2. Introduction en 4 éléments} & Accroche polie ? Contexte/Commande ? Objet précis ? Annonce du plan ? & $\square$ & $\square$ & & /5 \\ \hline
\textbf{3. Corps organisé et objectif} & Plan en 2 ou 3 parties logiques ? Titres de parties clairs ? \textbf{Zéro « je » / « nous » d'opinion} ? Faits chiffrés, datés, sourcés ? Connecteurs variés (ajout, opposition, conséquence) ? & $\square$ & $\square$ & & /10 \\ \hline
\textbf{4. Conclusion opérationnelle} & Bilan synthétique ? Propositions concrètes (numérotées, réalistes) ? Formule de politesse adaptée ? Signature ? & $\square$ & $\square$ & & /5 \\ \hline
\textbf{5. Liaisons et cohérence} & Pronoms relatifs (dont, où, lequel) utilisés ? Adverbes de liaison (toutefois, par conséquent, en outre) ? Absence de répétitions lourdes ? & $\square$ & $\square$ & & /5 \\ \hline
\textbf{6. Langue et présentation} & Orthographe/grammaire vérifiées (accords, temps) ? Paragraphes aérés ? Alignement, marges, police lisible ? Nom destinataire/date sans faute ? & $\square$ & $\square$ & & /5 \\ \hline
\rowcolor{popgoldL}
\textbf{TOTAL} & & & & & \textbf{/35} \\ \hline
\end{tabularx}
\end{center}
\legende{Grille d'auto-évaluation du rapport (35 points). Objectif Bac : viser 30/35 minimum. Chaque case « Non » entraîne une révision ciblée.}
\end{popfigure}

\begin{propriete}[Les 4 piliers d'un rapport réussi]
Un rapport technique de niveau Bac est validé s'il réunit simultanément :
\begin{enumerate}
    \item \textbf{Exhaustivité structurelle :} Aucune zone de l'en-tête n'est oubliée ; introduction (4 éléments), corps (parties équilibrées), conclusion (bilan + propositions) sont présents.
    \item \textbf{Objectivité factuelle :} Seuls des faits observables, mesurables, datés sont rapportés. Les adjectifs subjectifs (\emph{« catastrophique », « nul », « super »}) sont proscrits. On préfère \emph{« température relevée à 35°C »} à \emph{« chaleur insupportable »}.
    \item \textbf{Cohérence discursive :} La progression logique est assurée par des \cle{connecteurs variés} (adverbes, locutions, pronoms relatifs) et non par de simples « et », « mais », « aussi ».
    \item \textbf{Soin formel :} Zéro faute sur le nom du destinataire, la date, l'objet. Présentation aérée, paragraphes identifiables, écriture lisible.
\end{enumerate}
\end{propriete}

\courssub{L'amélioration du rapport : la méthode des trois relectures}

La différence entre un rapport « passable » et un rapport « excellent » se joue intégralement à l'étape 5. La relecture unique est un piège : l'œil ne peut pas tout contrôler à la fois.

\begin{methode}[La méthode des 3 relectures ciblées (10 minutes au Bac)]
\begin{enumerate}
    \item \textbf{Relecture 1 — Structure et Contenu (3 min) :} \textit{« Le squelette tient-il ? »} Vérifier : En-tête complet ? Intro 4/4 ? Plan respecté (titres = annonces) ? Corps = faits only ? Conclusion = bilan + 3 propositions ? Signature ?
    \item \textbf{Relecture 2 — Langue et Liaisons (4 min) :} \textit{« La chair est-elle saine ? »} Chasser : répétitions (remplacer par pronoms relatifs \emph{dont, où, lequel} ou synonymes), connecteurs pauvres (\emph{aussi, et, mais} $\to$ \emph{par conséquent, en outre, toutefois}), accords (sujet/verbe, participe passé), temps (imparfait pour le contexte, passé composé pour les faits ponctuels).
    \item \textbf{Relecture 3 — Présentation et Détails fatals (3 min) :} \textit{« L'habit fait le moine. »} Vérifier : Nom du Proviseur (orthographe exacte) ? Date du jour ? Objet clair ? Marges, saut de ligne entre paragraphes, lisibilité ? Numérotation des pages si > 1 page ?
\end{enumerate}
\end{methode}

\begin{attention}[Le piège fatal : la relecture négligée]
\textbf{Erreur fréquente :} « J'ai relu vite fait, ça va. »
\textbf{Conséquence :} Une faute sur le nom du destinataire (\emph{« Monsieur le Prviseur »}), une date erronée (\emph{« 32 mars »}), un objet flou (\emph{« Rapport sur le bac »}) ou un paragraphe unique sans saut de ligne font chuter la note de \textbf{5 à 10 points} sur la présentation et la cohérence, quel que soit le fond.
\textbf{Parade :} Appliquer la méthode des 3 relectures \textbf{systématiquement}, même en devoir surveillé. C'est un automatisme à acquérir dès maintenant.
\end{attention}

\courssub{Correction collective d'un rapport faible : analyse et réécriture}

L'analyse d'un contre-exemple est le meilleur entraînement à l'auto-correction. Voici un rapport « tel que rédigé par un élève en début d'unité » et sa version corrigée après application de la grille.

\begin{exemplebox}[Rapport faible — Version initiale (Avant)]
\textbf{Objet :} Le bac blanc

Monsieur le Proviseur,
Je vous écris pour vous dire comment s'est passé le bac blanc. C'était la semaine dernière. Globalement c'était pas terrible. D'abord il faisait trop chaud dans les salles, on a eu 35 degrés dans la salle 12, c'était insupportable. En plus il y avait du bruit de la cour, on entendait les cris des petits. Les sujets d'anglais étaient mal photocopiés, on voyait rien. Et les surveillants y en avait pas assez, un pour quarante, du coup y en a qui trichaient. Par contre l'eau c'était bien y en avait. Et l'infirmière elle était là. Pour finir je propose : mettre la clim, changer les salles, plus de surveillants. Voilà.
Cordialement,
Le délégué
\end{exemplebox}

\begin{exemplebox}[Analyse des défauts (Grille appliquée)]
\begin{itemize}
    \item \textbf{Structure externe :} Pas de date, pas de référence, objet trop vague (\emph{« Le bac blanc »}), pas de qualité de l'émetteur.
    \item \textbf{Introduction :} Absente (pas d'accroche, pas de contexte, pas d'annonce de plan).
    \item \textbf{Corps :} Un seul bloc (pas de paragraphes), subjectivité massive (\emph{« pas terrible », « insupportable », « on voit rien »}), familier (\emph{« y en a », « voilà »}), pas de connecteurs, mélange aspects positifs/négatifs.
    \item \textbf{Conclusion :} Propositions vagues (\emph{« mettre la clim »} — budget ?), pas de bilan, formule de politesse minimale.
    \item \textbf{Liaisons :} Quasi inexistantes (\emph{et, en plus, par contre}).
\end{itemize}
\end{exemplebox}

\begin{exemplebox}[Rapport amélioré — Version finale (Après application méthode 3 relectures)]
\begin{quote}
\textbf{RAPPORT N° 04/2026}

\textbf{Objet :} Conditions de déroulement du Baccalauréat blanc — Session de mars 2026 — Classe de Terminale G2.

\textbf{De :} \textit{Kamga Joseph}, Délégué de la classe de Terminale G2.

\textbf{À :} Monsieur le Proviseur \textit{Mbarga Étienne}, Lycée Technique de Douala-Bassa.

\textbf{Date :} 18 mars 2026.

\textbf{Introduction}
Monsieur le Proviseur,
La session du Baccalauréat blanc, organisée du 12 au 16 mars 2026, constitue une étape décisive dans la préparation de nos examens officiels. En ma qualité de délégué de la classe de Terminale G2, mandaté par mes camarades pour relayer leurs observations, j'ai l'honneur de vous présenter ce rapport sur les conditions matérielles et organisationnelles de ces épreuves. Ce document fait d'abord état des dispositions satisfaisantes, avant de signaler les difficultés rencontrées et de formuler des propositions d'amélioration pour les sessions ultérieures.

\textbf{1. Des dispositions globales satisfaisantes}
Premièrement, l'organisation générale a fait preuve de rigueur : le calendrier a été communiqué à l'avance et les horaires de début d'épreuves (8h00) ont été scrupuleusement respectés. Par ailleurs, la mise à disposition quotidienne de bouteilles d'eau dans les salles et la présence effective de l'infirmière scolaire ont été très appréciées par l'ensemble des candidats. Ces mesures témoignent de la sollicitude de l'administration pour le bien-être des élèves.

\textbf{2. Des conditions matérielles et environnementales perfectibles}
Toutefois, plusieurs dysfonctionnements ont perturbé la concentration des candidats. En premier lieu, la température dans la salle 12 a été relevée à 35°C à 14h00 les 13 et 14 mars, l'absence de ventilation rendant l'atmosphère difficilement supportable pour une épreuve de trois heures. En second lieu, la proximité immédiate de la cour de récréation a généré une pollution sonore continue (cris, ballons) lors des épreuves de l'après-midi. Enfin, la qualité de reproduction des sujets d'Anglais (session du 14 mars) était défectueuse : sur 42 copies, 15 présentaient des zones illisibles, obligeant les surveillants à dicter les consignes, ce qui a fait perdre 10 minutes aux candidats concernés.

\textbf{3. Des insuffisances dans l'encadrement et la gestion du temps}
Par ailleurs, le ratio d'encadrement (1 surveillant pour 40 candidats) s'est révélé insuffisant pour garantir la sérénité des épreuves. Des allées et venues non autorisées ont été observées. De surcroît, l'absence de pause méridienne structurée (les épreuves s'enchaînant de 8h00 à 18h30 sans interruption supérieure à 20 minutes) a engendré une fatigue marquée chez les élèves, préjudiciable aux performances de l'après-midi.

\textbf{Conclusion}
Monsieur le Proviseur, au bilan, si la volonté d'organiser cet examen blanc est louable, les conditions matérielles et d'encadrement ont nui à l'équité entre candidats. Pour les sessions futures, je me permets de formuler trois propositions concrètes :
\begin{enumerate}
    \item Affecter les épreuves de l'après-midi dans des salles éloignées de la cour de récréation ou disposant de ventilateurs.
    \item Veiller à un contrôle qualité systématique des photocopies (lisibilité, complétude) avant distribution.
    \item Renforcer l'équipe de surveillance (ratio 1/25) et inscrire une pause méridienne d'une heure au planning officiel.
\end{enumerate}
Je vous prie d'agréer, Monsieur le Proviseur, l'expression de mon respectueux dévouement.

\textit{Signature}
\textbf{Kamga Joseph}
\end{quote}
\end{exemplebox}

\begin{aretenir}[Principales corrections opérées]
\begin{itemize}
    \item \textbf{Structure :} En-tête complet (Référence, Qualité, Destinataire exact, Date). Introduction en 4 temps. Corps en 3 parties titrées (Positif / Matériel / Organisation). Conclusion Bilan + 3 propositions numérotées + Politesse administrative + Signature.
    \item \textbf{Objectivité :} \emph{« C'était pas terrible »} $\to$ \emph{« Conditions perfectibles »} ; \emph{« 35 degrés, insupportable »} $\to$ \emph{« Température relevée à 35°C »} ; \emph{« Sujets mal photocopiés »} $\to$ \emph{« 15 sur 42 copies illisibles »}.
    \item \textbf{Liaisons enrichies :} \emph{Premièrement, Par ailleurs, En premier lieu, En second lieu, Enfin, Toutefois, De surcroît} (adverbes/locutions) ; \emph{dont, lesquelles} (pronoms relatifs implicites dans la structure).
    \item \textbf{Langue :} Registre soutenu, accords corrects, pas de familier.
\end{itemize}
\end{aretenir}

\courssub{Compte rendu, remédiation et gestion du temps au Bac}

\begin{savaistu}[Au Baccalauréat technique : les 4 critères de notation]
L'épreuve de production écrite (souvent un rapport ou une synthèse) est notée sur 20 points répartis selon 4 compétences pondérées :
\begin{itemize}
    \item \textbf{Pertinence du contenu / Respect de la consigne (5 pts) :} Objet traité ? Format respecté (rapport et non lettre) ? Destinataire identifié ? Propositions réalistes ?
    \item \textbf{Cohérence et organisation (5 pts) :} Structure externe/intérieure ? Introduction/Conclusion normées ? Articulation logique des parties ? Connecteurs efficaces ?
    \item \textbf{Correction de la langue (5 pts) :} Syntaxe, orthographe, grammaire, vocabulaire technique/approprié. \textbf{Seuil d'alerte :} Plus de 5 fautes graves = note plafonnée.
    \item \textbf{Présentation (5 pts) :} Soin graphique, aération, lisibilité, en-tête complet, absence de ratures.
\end{itemize}
\textbf{Astuce :} Un rapport bien présenté, structuré et sans faute de français, mais au contenu moyen, vaut souvent mieux qu'un rapport riche mais illisible et mal structuré.
\end{savaistu}

\begin{methode}[Gestion des 50 minutes à l'épreuve du Bac (Probatoire/Bac)]
\begin{enumerate}
    \item \textbf{Minutes 0–10 : ANALYSE \& PLAN (Brouillon obligatoire).}
        \begin{itemize}
            \item Lire le sujet 2 fois. Souligner : Émetteur, Destinataire, Objet, Finalité, Contraintes (documents fournis ?).
            \item Lister les idées/arguments sur brouillon (mots-clés, chiffres).
            \item Construire le plan détaillé (Titres des parties + 2 arguments par partie + 3 propositions pour la conclusion).
        \end{itemize}
    \item \textbf{Minutes 10–40 : RÉDACTION (Copie finale).}
        \begin{itemize}
            \item Rédiger l'en-tête et l'introduction (soigner l'accroche et l'annonce de plan).
            \item Développer le corps : 1 paragraphe = 1 idée force + fait chiffré + connecteur. Sauter une ligne entre paragraphes.
            \item Rédiger la conclusion : Bilan (1 phrase) + Propositions (liste numérotée) + Politesse + Signature.
        \end{itemize}
    \item \textbf{Minutes 40–50 : RELECTURE 3 PASSES (Méthode ci-dessus).}
        \begin{itemize}
            \item Pass 1 : Structure (En-tête, Intro 4/4, Plan respecté, Conclu Bilan+Props).
            \item Pass 2 : Langue (Accords, Temps, Connecteurs variés, Répétitions $\to$ Pronoms relatifs).
            \item Pass 3 : Présentation (Nom destinataire, Date, Objet, Aération, Lisibilité, Ratures propres).
        \end{itemize}
\end{enumerate}
\end{methode}

\begin{exercice}[Entraînement chronométré — Sujet type Bac]
\textbf{Sujet :} Vous êtes le Chef d'équipe de maintenance du Lycée Technique. Suite à une panne de courant survenue le 15 avril 2026 de 10h à 14h, vous rédigez un rapport au Chef de travaux pour expliquer les causes, décrire les conséquences sur les TP et proposer des mesures préventives.
\textbf{Consignes :} Appliquez la méthode des 50 min. Utilisez la grille d'auto-évaluation pour noter votre production sur 35 pts.
\end{exercice}

\begin{corrige}[Exercice n°1 — Éléments attendus pour l'auto-correction]
\begin{itemize}
    \item \textbf{En-tête :} Référence, De (Chef équipe maintenance), À (Chef de travaux), Date (16 avril), Objet précis (Panne électrique 15/04 — Causes, Conséquences TP, Prévention).
    \item \textbf{Intro :} Contexte (panne 10h-14h), Mandat (demande du CdT), Objet, Plan (1. Causes techniques, 2. Impact TP, 3. Prévention).
    \item \textbf{Corps :} Partie 1 : Origine (câble sectionné travaux extérieurs / disjoncteur vétuste — faits techniques). Partie 2 : Conséquences (Arrêt 4 groupes électrogènes, perte 3h TP Électrotechnique G2/G3, remise à zéro paramètres machines CN). Partie 3 : Propositions (Gaine de protection câbles, Remplacement disjoncteurs > 10 ans, Contrat maintenance onduleurs).
    \item \textbf{Conclusion :} Bilan (perturbation majeure mais gérée), 3 propositions numérotées, Politesse administrative, Signature.
    \item \textbf{Langue :} Passé composé (faits), Imparfait (contexte), Futur/Conditionnel (propositions). Connecteurs : \emph{En effet, Par conséquent, Afin de prévenir, En outre}.
\end{itemize}
\end{corrige}

\begin{aretenir}[Bilan de l'unité V : Produire un rapport]
\begin{itemize}
    \item \textbf{Savoir :} Structure externe (en-tête normalisé) et interne (Intro 4 éléments, Corps parties thématiques, Conclu Bilan+Propositions). Outils de liaison (adverbes, locutions, pronoms relatifs \emph{dont, où, lequel}).
    \item \textbf{Savoir-faire :} Analyser une commande, collecter/hiérarchiser l'info, planifier, rédiger objectif, relire en 3 passes.
    \item \textbf{Réussite au Bac :} \cle{Respect du format} + \cle{Objectivité factuelle} + \cle{Connecteurs variés} + \cle{Zéro faute sur destinataire/date} + \cle{Gestion du temps (10/30/10)}.
\end{itemize}
\end{aretenir}

\newpage

\coursec{Activité d'intégration}

\courssub{Objectifs de l'activité}

À l'issue de cette activité d'intégration, l'élève sera capable de :
\begin{itemize}
\item analyser une \cle{situation de communication} professionnelle (destinataire, objet, circonstances) afin d'adapter son écriture ;
\item exploiter un \cle{dossier documentaire} varié (lettre de mission, photographies, tableau statistique, témoignages) et en hiérarchiser les informations ;
\item construire le \cle{plan} d'un rapport (introduction, corps en deux parties, conclusion) conforme à la structure externe et interne étudiée ;
\item rédiger un rapport complet, objectif et cohérent, en employant correctement les \cle{outils de liaison} (adverbes de liaison, locutions adverbiales, pronoms relatifs) ;
\item évaluer et améliorer un texte à l'aide d'une \cle{grille d'auto-évaluation}, puis produire une version finale individuelle conforme aux exigences du Baccalauréat technique.
\end{itemize}

\courssub{Situation}

Tu es \cle{agent municipal} de la commune de Batcham, dans la région de l'Ouest. Après les fortes pluies du mois de septembre, le maire t'a chargé(e) d'inspecter l'état de la route rurale qui relie les villages de Batcham-Centre, Bangang et Bameka, sur une distance totale de \num{18} km. À ton retour, tu dois lui remettre un \cle{rapport de mission} écrit. Tu disposes du dossier suivant.

\textbf{Document 1 : extrait de la lettre de mission du maire.}

\begin{exemplebox}[Lettre de mission]
« Monsieur / Madame, vous êtes chargé(e) d'inspecter, du 25 au 27 septembre, l'état de la route Batcham-Centre -- Bangang -- Bameka, fortement dégradée par les pluies. Vous recueillerez les observations nécessaires, les données chiffrées et les témoignages des usagers, puis vous me remettrez un rapport circonstancié assorti de suggestions. Le Maire. »
\end{exemplebox}

\textbf{Document 2 : deux photographies.} La photographie 1 montre de profondes \cle{ornières} remplies d'eau sur le tronçon Batcham-Centre -- Bangang ; un camion de transport de vivres y est embourbé. La photographie 2 montre le petit pont de la rivière Mfou, entre Bangang et Bameka, dont le tablier s'est effondré ; des motocyclistes traversent à gué.

\textbf{Document 3 : tableau des données recueillies.}

\begin{center}
\begin{tabularx}{\linewidth}{|l|X|X|X|}\hline
\rowcolor{popblueL} \textbf{Tronçon} & \textbf{Longueur endommagée} & \textbf{Véhicules bloqués} & \textbf{Population touchée} \\ \hline
Batcham-Centre -- Bangang & 7 km & 14 véhicules & \num{6200} habitants \\ \hline
Bangang -- Bameka & 4 km (pont effondré) & 9 véhicules & \num{4800} habitants \\ \hline
Bameka -- marché hebdomadaire & 2 km & 5 véhicules & \num{3100} habitants \\ \hline
\textbf{Total} & \textbf{13 km sur 18 km} & \textbf{28 véhicules} & \textbf{\num{14100} habitants} \\ \hline
\end{tabularx}
\end{center}

\textbf{Document 4 : trois témoignages d'usagers.}

\begin{itemize}
\item Mme Ngo Bell, commerçante : « Mes sacs de haricots pourrissent au village parce qu'aucun camion ne passe. Je perds environ \num{150000} FCFA par semaine. »
\item M. Kamdem, chauffeur de transport en commun : « Le pont effondré nous oblige à faire un détour de 22 km ; le prix du transport a doublé. »
\item Le chef du village de Bameka : « Les malades ne peuvent plus rejoindre l'hôpital de district ; une femme enceinte a dû être portée à bras pendant six kilomètres. »
\end{itemize}

\courssub{Consignes}

Travail en groupes de quatre, puis version finale individuelle.

\begin{enumerate}
\item Analyse la situation de communication : qui écrit ? à qui ? pour quoi faire ? Quel ton et quel point de vue (subjectif ou objectif) conviennent à ce texte ?
\item Lis l'ensemble du dossier et classe les informations en deux catégories : les \cle{constats chiffrés} d'une part, les \cle{conséquences pour les populations} d'autre part.
\item Établis le plan détaillé du rapport : introduction (contexte, mission, méthode, annonce du plan), corps en deux parties, conclusion.
\item Rédige l'\cle{introduction} (8 à 12 lignes) : elle doit présenter le contexte, rappeler la mission reçue, indiquer la méthode d'enquête et annoncer le plan.
\item Rédige la \cle{première partie} du corps : présente les dégâts constatés en t'appuyant sur les données du tableau et sur les deux photographies (décris objectivement ce qu'elles montrent).
\item Rédige la \cle{deuxième partie} du corps : analyse les conséquences économiques et sociales pour les populations, en exploitant les trois témoignages (citation brève ou reformulation au discours indirect).
\item Rédige la \cle{conclusion} : fais le bilan de la mission et formule \textbf{trois suggestions} précises adressées au maire.
\item Vérifie que ton rapport contient au moins \textbf{cinq outils de liaison} différents (adverbes ou locutions : \emph{d'abord, ensuite, en outre, par ailleurs, en effet, cependant, en conséquence, enfin}...) et au moins \textbf{trois pronoms relatifs} (\emph{qui, que, où, dont}).
\item Échange ton rapport avec un autre groupe et évalue-le à l'aide de la grille d'auto-évaluation ; note deux points forts et deux points à améliorer.
\item Rédige ta version finale individuelle en tenant compte des remarques reçues, puis remets-la pour correction.
\end{enumerate}

\courssub{Ressources à mobiliser}

\begin{aretenir}[Ce qu'il faut se rappeler]
\begin{itemize}
\item \textbf{Structure externe} : en-tête (auteur, destinataire, lieu, date, objet), titre, corps du texte, signature.
\item \textbf{Structure interne} : introduction (contexte, mission, méthode, plan) ; corps en parties équilibrées ; conclusion (bilan + suggestions).
\item \textbf{Outils de liaison} : adverbes de liaison (\emph{d'abord, ensuite, enfin}), locutions adverbiales (\emph{en outre, par ailleurs, en conséquence}), pronoms relatifs (\emph{qui, que, où, dont}) qui relient les phrases et évitent les répétitions.
\item \textbf{Exploitation des textes iconiques} : décrire objectivement ce que montre une photographie, sans imagination ni commentaire personnel.
\item \textbf{Ton du rapport} : objectivité (faits, chiffres, discours indirect), phrases claires, vocabulaire précis, pas de « je » émotif.
\end{itemize}
\end{aretenir}

\begin{attention}[Erreurs fréquentes à éviter]
Ne transforme pas le rapport en récit d'aventures (« j'ai vu, j'ai marché... ») : le rapport relate des \textbf{constats}, pas des exploits. N'oublie pas l'en-tête administratif. Évite l'accumulation de phrases juxtaposées sans connecteurs. Ne copie pas les témoignages au discours direct dans tout le texte : privilégie le discours indirect (« Mme Ngo Bell a déclaré que... »). Enfin, les suggestions doivent être concrètes et adressées au destinataire, pas des vœux vagues.
\end{attention}

\courssub{Démarche et corrigé commenté}

\begin{exoresolu}[Rapport de mission : modèle rédigé et commenté]
Énoncé : à partir du dossier fourni, rédiger le rapport complet que l'agent municipal remet au maire de Batcham.
\tcblower
\textbf{Étape 1 : analyse de la situation.} L'émetteur est un agent municipal ; le destinataire est le maire (supérieur hiérarchique) ; l'objet est le compte rendu d'une inspection. Le ton doit donc être \emph{formel, objectif et factuel} ; la troisième personne et le discours indirect dominent.

\textbf{Étape 2 : classement des informations.} Constats chiffrés : 13 km endommagés sur 18 (soit environ 72 \% du tracé), pont effondré, 28 véhicules bloqués, deux photographies. Conséquences : \num{14100} habitants touchés, pertes commerciales (\num{150000} FCFA par semaine), hausse des tarifs de transport, isolement sanitaire.

\textbf{Étape 3 : rédaction du rapport (modèle).}

Commune de Batcham -- Service technique municipal\par
Rapport de mission d'inspection de la route Batcham-Centre -- Bangang -- Bameka\par
Batcham, le 28 septembre\par
À l'attention de Monsieur le Maire\par

\emph{Introduction.} À la suite des fortes pluies enregistrées en septembre, la route rurale qui relie les villages de Batcham-Centre, Bangang et Bameka est devenue quasiment impraticable. Conformément à la lettre de mission du 24 septembre, une inspection a été menée du 25 au 27 septembre sur l'ensemble du tracé, qui s'étend sur 18 km. L'enquête a porté sur l'observation directe des dégâts, la collecte de données chiffrées et le recueil de témoignages d'usagers. Le présent rapport expose d'abord les constats effectués sur le terrain, puis les conséquences de cette situation pour les populations, avant de formuler des suggestions.

\emph{Première partie : les constats.} L'inspection révèle que 13 km de route sur 18 sont endommagés, soit environ 72 \% du tracé. Sur le tronçon Batcham-Centre -- Bangang, qui compte 7 km dégradés, de profondes ornières remplies d'eau rendent la circulation impossible, comme le montre la photographie 1 où un camion de vivres est embourbé. En outre, le pont de la rivière Mfou, situé entre Bangang et Bameka, s'est effondré : la photographie 2 montre des motocyclistes qui traversent à gué, au péril de leur vie. Au total, 28 véhicules sont restés bloqués pendant la période d'observation.

\emph{Deuxième partie : les conséquences pour les populations.} Cette dégradation touche environ \num{14100} habitants. Sur le plan économique, Mme Ngo Bell, commerçante, a déclaré que ses sacs de haricots pourrissent au village, ce qui lui fait perdre environ \num{150000} FCFA par semaine. Par ailleurs, le détour de 22 km imposé par l'effondrement du pont a fait doubler le prix du transport, selon M. Kamdem, chauffeur. Sur le plan sanitaire, la situation est grave : le chef du village de Bameka, dont la localité est coupée de l'hôpital de district, a signalé qu'une femme enceinte avait dû être portée à bras sur six kilomètres.

\emph{Conclusion.} La mission a permis d'établir que la route est dans un état critique et que l'isolement des trois villages menace l'économie locale et la santé des populations. En conséquence, il est suggéré à Monsieur le Maire : 1) de lancer en urgence le reprofilage des 7 km du premier tronçon ; 2) de solliciter auprès de la délégation régionale des Travaux publics la reconstruction du pont de la Mfou ; 3) d'organiser, en attendant, un passage provisoire sécurisé à gué pour les piétons et les motos.\par
L'agent municipal, [signature]

\textbf{Étape 4 : vérification.} Le texte contient bien cinq outils de liaison (\emph{d'abord, puis, en outre, par ailleurs, en conséquence}) et trois pronoms relatifs (\emph{qui relie}, \emph{où un camion}, \emph{dont la localité}). Les photographies sont décrites objectivement, les chiffres sont repris exactement (13 km sur 18, soit environ 72 \% ; 28 véhicules ; \num{14100} habitants), et les témoignages sont rapportés au discours indirect.
\end{exoresolu}

\courssub{Bilan de l'activité}

Cette activité a permis de mobiliser l'ensemble des savoirs de la leçon dans une situation professionnelle authentique. Elle a d'abord montré qu'un rapport n'est pas un texte libre : il répond à une \cle{commande} (la lettre de mission), s'adresse à un destinataire précis et obéit à une structure rigoureuse (en-tête, introduction, corps, conclusion, signature). Elle a ensuite entraîné les élèves à \cle{exploiter un dossier documentaire} : hiérarchiser les informations, décrire des photographies sans les interpréter, reprendre exactement des données chiffrées et transformer des témoignages oraux en discours indirect. Elle a enfin mis en pratique les \cle{outils de liaison}, qui donnent au texte sa cohérence et sa fluidité, compétence directement évaluée au Baccalauréat technique. La phase d'amélioration croisée a développé l'esprit critique : relire le texte d'autrui avec une grille, c'est apprendre à relire le sien. Ces compétences sont transférables à la vie professionnelle : stage, emploi, démarches administratives, partout où il faut rendre compte par écrit de manière claire et objective.

\newpage

\coursec{Méthodes et exercices résolus}

\courssub{Méthodes et exercices résolus}

\begin{methode}[Identifier et analyser un rapport comme texte fonctionnel]
    \textbf{Objectif :} reconnaître un rapport, en déterminer la nature, la finalité et le public destinataire.
    \begin{enumerate}
        \item \textbf{Repérer les indices paratextuels :} titre explicite (\og Rapport de stage\fg, \og Rapport d'enquête\fg), en-tête (émetteur, destinataire, date, objet), signature.
        \item \textbf{Analyser la finalité (fonction) :} informer objectivement (\cle{fonction informative}), rendre compte d'une action, proposer des solutions (\cle{fonction argumentative} secondaire), archiver.
        \item \textbf{Identifier le public :} hiérarchie, client, jury, administration $\rightarrow$ adapter le ton (formel, neutre, impersonnel).
        \item \textbf{Qualifier le genre :} texte \cle{fonctionnel} (utilitaire), \cle{professionnel}, \cle{structuré}, \cle{non littéraire}.
    \end{enumerate}
\end{methode}

\begin{exoresolu}[Analyse de la situation de communication d'un rapport]
    \textbf{Énoncé :} vous êtes stagiaire au service maintenance de la société \textsc{Camwater} à Douala. Votre chef de service vous demande de rédiger un rapport sur la panne survenue le 15 mars 2026 à la station de pompage de Nkolfoulou. Présentez la situation de communication de ce futur rapport sous forme de tableau.
    \tcblower
    \textbf{Solution détaillée :}
    \begin{enumerate}
        \item \textbf{Émetteur :} moi, [Prénom Nom], stagiaire au service Maintenance, \textsc{Camwater} Douala.
        \item \textbf{Destinataire (principal) :} M. le Chef de Service Maintenance, \textsc{Camwater} Douala. \emph{(Destinataires secondaires : Direction Régionale, Service Qualité.)}
        \item \textbf{Objet / Sujet :} panne de la pompe principale n°3 à la station de Nkolfoulou (15/03/2026) : causes, conséquences, interventions, recommandations.
        \item \textbf{Finalité (but) :} informer la hiérarchie sur le déroulement de l'incident, justifier les décisions prises, proposer des mesures préventives pour éviter la récidive. Fonction \textbf{principalement informative et préconisatrice}.
        \item \textbf{Contexte / Circonstances :} stage professionnel, urgence technique, cadre administratif formel.
        \item \textbf{Support / Canal :} document écrit (numérique ou papier), format administratif standard (en-tête, corps, signature).
        \item \textbf{Ton / Registre :} formel, technique, objectif, impersonnel (usage du \og nous\fg{} institutionnel ou de la voix passive), précis.
    \end{enumerate}
    \begin{center}
        \begin{tabularx}{\linewidth}{|l|X|}\hline
        \rowcolor{popblueL} \textbf{Composante} & \textbf{Détails} \\ \hline
        Émetteur & Stagiaire Maintenance, Camwater Douala \\ \hline
        Destinataire & Chef Service Maintenance (copie : Direction) \\ \hline
        Objet & Panne pompe n°3 Station Nkolfoulou (15/03/2026) \\ \hline
        Finalité & Informer, rendre compte, proposer des solutions \\ \hline
        Registre & Formel, technique, objectif \\ \hline
    \end{tabularx}
    \end{center}
    \textbf{Conclusion :} ce tableau constitue la \cle{fiche d'identité} du rapport ; elle guide tout le processus de rédaction (choix du vocabulaire, structure, degré de détail).
\end{exoresolu}

\begin{methode}[Maîtriser la structure externe et interne du rapport]
    \textbf{Objectif :} organiser le document selon les normes professionnelles (macro et micro-structure).
    \begin{enumerate}
        \item \textbf{Structure externe (apparence) :}
        \begin{itemize}
            \item \textbf{En-tête :} logo/en-tête de la structure, référence (N°/Année), date, objet (clair, précis), De / À / Copie à.
            \item \textbf{Corps :} introduction, développement (corps), conclusion/recommandations.
            \item \textbf{Bas de page :} signature (nom, grade/fonction), visa hiérarchique, pièces jointes (PJ), diffusion.
        \end{itemize}
        \item \textbf{Structure interne (organisation des idées) :}
        \begin{itemize}
            \item \textbf{Plan thématique, chronologique ou logique :} chapitres $\rightarrow$ sections $\rightarrow$ paragraphes.
            \item \textbf{Numérotation décimale :} 1. / 1.1. / 1.1.1.
            \item \textbf{Titrage :} titres courts, nominatifs (sans verbe conjugué), parallèles entre eux.
        \end{itemize}
    \end{enumerate}
\end{methode}

\begin{exoresolu}[Structurer le plan d'un rapport technique]
    \textbf{Énoncé :} à partir du sujet du rapport sur la panne de la station de Nkolfoulou (exercice précédent), proposez un plan détaillé (structure interne) numéroté, avec des titres nominatifs adaptés au contexte technique.
    \tcblower
    \textbf{Solution détaillée :}

    \textbf{Plan détaillé du rapport :}
    \begin{enumerate}
        \item \textbf{INTRODUCTION}
        \item \textbf{1. CIRCONSTANCES ET DESCRIPTION DE L'INCIDENT}
              \begin{enumerate}
                  \item[1.1.] Contexte et signalement de l'anomalie
                  \item[1.2.] Description technique de la panne (symptômes, horaire, équipement concerné)
                  \item[1.3.] Impact immédiat sur la production et la distribution
              \end{enumerate}
        \item \textbf{2. ANALYSE DES CAUSES ET INTERVENTIONS RÉALISÉES}
              \begin{enumerate}
                  \item[2.1.] Diagnostic technique et identification de la cause racine (rupture du roulement, surchauffe)
                  \item[2.2.] Chronologie des interventions de l'équipe de maintenance
                  \item[2.3.] Moyens matériels et humains mobilisés (pièces de rechange, outillage, personnel)
              \end{enumerate}
        \item \textbf{3. BILAN ET ÉVALUATION DES PERTES}
              \begin{enumerate}
                  \item[3.1.] Durée d'arrêt de la station
                  \item[3.2.] Volume d'eau non produit et non distribué
                  \item[3.3.] Coût estimé de la réparation et de la perte d'exploitation
              \end{enumerate}
        \item \textbf{4. RECOMMANDATIONS ET MESURES PRÉVENTIVES}
              \begin{enumerate}
                  \item[4.1.] Maintenance préventive : révision du planning de graissage et de contrôle des roulements
                  \item[4.2.] Stock de sécurité : constitution d'un stock de roulements de rechange
                  \item[4.3.] Formation et sensibilisation : procédure d'alerte précoce (vibrations, température)
              \end{enumerate}
        \item \textbf{CONCLUSION GÉNÉRALE}
        \item \textbf{PIÈCES JOINTES} (fiche d'intervention, photos, devis, courbes de surveillance)
    \end{enumerate}
    \textbf{Justification :} le plan suit une \cle{logique chronologico-logique} (faits $\rightarrow$ causes et actions $\rightarrow$ bilan $\rightarrow$ avenir). Les titres sont \textbf{nominatifs} (pas de phrase verbale). La numérotation est \textbf{décimale}.
\end{exoresolu}

\begin{methode}[Mobiliser les outils de liaison pour la cohérence textuelle]
    \textbf{Objectif :} assurer la fluidité et la logique du discours grâce aux connecteurs syntaxiques.
    \begin{enumerate}
        \item \textbf{Adverbes de liaison et locutions adverbiales :} placés en tête de phrase ou en incise, ils marquent la relation logique entre phrases et paragraphes.
        \begin{itemize}
            \item \cle{Chronologie :} \og d'abord\fg, \og ensuite\fg, \og finalement\fg, \og par la suite\fg.
            \item \cle{Cause et conséquence :} \og par conséquent\fg, \og donc\fg, \og c'est pourquoi\fg, \og en effet\fg{} (explication).
            \item \cle{Opposition et concession :} \og cependant\fg, \og néanmoins\fg, \og toutefois\fg, \og malgré cela\fg.
            \item \cle{Addition et énumération :} \og de plus\fg, \og par ailleurs\fg, \og en outre\fg.
        \end{itemize}
        \item \textbf{Pronoms relatifs (qui, que, dont, où, lequel...) :} ils évitent la répétition et créent des \cle{subordonnées relatives} (déterminatives ou explicatives). \emph{Exemple :} \og La pompe \textbf{dont} le roulement a cédé...\fg
        \item \textbf{Reprises pronominales (le, la, les, en, y, ce qui, ce que) :} elles assurent le chaînage anaphorique.
        \item \textbf{Connecteurs de phrase (prépositions, conjonctions) :} \og afin de\fg{} + infinitif (but), \og bien que\fg{} + subjonctif (concession).
    \end{enumerate}
    \textbf{Règle d'or :} varier les outils ; ne pas surcharger (un connecteur par phrase en général) ; ponctuer correctement (virgule après un connecteur placé en tête de phrase).
\end{methode}

\begin{exoresolu}[Réécrire un passage en améliorant la cohésion]
    \textbf{Énoncé :} le texte suivant manque de fluidité. Réécrivez-le en intégrant des outils de liaison adaptés (adverbes, pronoms relatifs, reprises) sans en changer le sens technique.

    \og L'équipe arrive sur site. L'équipe constate la surchauffe. Le roulement est détruit. L'équipe démonte la pompe. L'équipe remplace la pièce. La pompe redémarre. La pression reste instable. L'équipe vérifie l'alignement. L'alignement est défectueux. L'équipe corrige l'alignement. La pression redevient normale.\fg
    \tcblower
    \textbf{Solution détaillée :}

    \textbf{Proposition de réécriture :}

    \og \textbf{À son arrivée sur site}, l'équipe \textbf{constate} une surchauffe majeure : le roulement \textbf{est} totalement détruit. \textbf{Par conséquent}, elle procède au démontage de la pompe \textbf{afin de remplacer} la pièce défectueuse. \textbf{Suite à cette intervention}, la pompe redémarre ; \textbf{cependant}, la pression \textbf{demeure} instable. \textbf{Dès lors}, les techniciens vérifient l'alignement, \textbf{qui s'avère} défectueux. \textbf{Après la correction de ce dernier}, la pression redevient normale.\fg
    \begin{itemize}
        \item \textbf{Chronologie :} \og à son arrivée\fg, \og par conséquent\fg, \og suite à\fg, \og dès lors\fg, \og après\fg.
        \item \textbf{Pronoms relatifs et reprises :} \og \textbf{qui} s'avère\fg{} (antécédent : l'alignement), \og \textbf{ce dernier}\fg{} (reprise nominale), \og \textbf{elle}\fg{} (reprise pronominale sujet pour l'équipe).
        \item \textbf{Variété syntaxique :} voix passive (\og est détruit\fg), infinitif de but (\og afin de remplacer\fg), subordonnées relatives.
    \end{itemize}
    \textbf{Conclusion :} le texte gagne en \cle{professionnalisme} et en \cle{lisibilité} ; la logique chronologique et la relation de cause à effet deviennent explicites.
\end{exoresolu}

\begin{methode}[Rédiger l'introduction et la conclusion d'un rapport]
    \textbf{Objectif :} produire les deux parties cadres en respectant le protocole administratif.
    \begin{enumerate}
        \item \textbf{L'introduction (du général vers le particulier) :}
        \begin{enumerate}
            \item \textbf{Objet et référence :} \og J'ai l'honneur de vous présenter le rapport relatif à...\fg{} (référence de la commande ou de l'ordre de mission).
            \item \textbf{Contexte et circonstances :} quand ? où ? pourquoi ? (numéro d'ordre de mission, date du stage, survenance de l'incident).
            \item \textbf{Sujet précis :} définition claire du thème traité.
            \item \textbf{Annonce du plan :} \og Ce rapport s'articule en X parties : 1... 2... 3...\fg
        \end{enumerate}
        \item \textbf{La conclusion (du particulier vers le général) :}
        \begin{enumerate}
            \item \textbf{Synthèse des résultats :} rappel très bref des constats majeurs (sans élément nouveau).
            \item \textbf{Réponse à l'objectif :} \og Cette étude a permis de montrer que...\fg
            \item \textbf{Recommandations et perspectives :} \og Je vous propose... / Il est recommandé de...\fg{} (ouverture vers l'action).
            \item \textbf{Formule de politesse finale :} \og En espérant que ce rapport répondra à votre attente, je vous prie d'agréer...\fg
        \end{enumerate}
    \end{enumerate}
\end{methode}

\begin{exoresolu}[Rédiger l'introduction et la conclusion du rapport Camwater]
    \textbf{Énoncé :} rédigez l'introduction et la conclusion complètes du rapport sur la panne de Nkolfoulou (voir le plan de l'exercice précédent). Respectez les conventions administratives.
    \tcblower
    \textbf{Solution détaillée :}

    \textbf{INTRODUCTION}
    \begin{quote}
    \textsc{Camwater} — Direction Régionale du Littoral — Service Maintenance\\
    \textbf{Réf. :} 045/MR/CAM/DR-L/2026 \hfill Douala, le 20 mars 2026\\
    \textbf{Objet :} rapport d'incident — panne de la pompe n°3 de la station de Nkolfoulou du 15/03/2026\\
    \textbf{De :} [Nom du stagiaire], stagiaire Maintenance\\
    \textbf{À :} M. le Chef de Service Maintenance\\
    \textbf{Copie :} M. le Directeur Régional
    \vspace{0.3cm}

    Monsieur le Chef de Service,

    Par la présente, j'ai l'honneur de vous soumettre le rapport relatif à l'incident survenu le 15 mars 2026 sur la pompe principale n°3 de la station de pompage de Nkolfoulou, conformément à votre instruction du 16 mars 2026.

    Cet incident a entraîné l'arrêt total de la station pendant six heures, perturbant la desserte en eau potable des quartiers Est de Douala. Ce document retrace les circonstances de la panne, analyse les causes techniques identifiées, évalue l'impact sur la production et formule des recommandations pour sécuriser l'exploitation future.

    Ce rapport s'articule en quatre parties :
    \begin{enumerate}
        \item les circonstances et la description de l'incident ;
        \item l'analyse des causes et les interventions réalisées ;
        \item le bilan et l'évaluation des pertes ;
        \item les recommandations et les mesures préventives.
    \end{enumerate}
    \end{quote}

    \textbf{CONCLUSION GÉNÉRALE}
    \begin{quote}
    Au terme de cette analyse, il ressort que la panne de la pompe n°3 trouve son origine dans la rupture du roulement à billes côté entraînement, due à un défaut de graissage non détecté lors des dernières visites de maintenance préventive. L'intervention curative, bien que rapide (six heures d'arrêt), a engendré un déficit de production de \num{12000} m\textsuperscript{3} et un coût total estimé à \num{1250000} FCFA.

    Pour prévenir toute récidive, je recommande :
    \begin{enumerate}
        \item la révision immédiate du planning de graissage (passage d'une fréquence mensuelle à une fréquence bimensuelle) ;
        \item l'approvisionnement d'un stock de sécurité de deux jeux de roulements de rechange ;
        \item l'installation de capteurs de vibrations et de température avec alarme déportée.
    \end{enumerate}

    En espérant que les propositions ci-dessus retiendront votre attention, je vous prie d'agréer, Monsieur le Chef de Service, l'expression de ma considération distinguée.

    \hfill \textbf{[Signature]} \\
    \hfill [Nom Prénom], stagiaire Maintenance
    \end{quote}
    \textbf{Analyse :} l'introduction cite la \cle{référence administrative}, le \cle{contexte}, le \cle{sujet} et \cle{annonce le plan}. La conclusion fait la \cle{synthèse technique}, le \cle{bilan chiffré}, les \cle{recommandations numérotées} et se termine par la \cle{formule de politesse}.
\end{exoresolu}

\begin{methode}[Rédiger le corps du rapport et exploiter des documents iconiques]
    \textbf{Objectif :} développer les parties annoncées avec rigueur ; intégrer schémas, tableaux et photos comme \cle{preuves} ou \cle{illustrations}.
    \begin{enumerate}
        \item \textbf{Un paragraphe = une idée force :} phrase d'introduction de l'idée + arguments et données + phrase de transition.
        \item \textbf{Style impersonnel et objectif :} voix passive (\og a été constaté\fg), tournures comme \og il a été relevé\fg, infinitif (\og vérifier l'alignement\fg), \og nous\fg{} institutionnel.
        \item \textbf{Intégration des documents iconiques (tableaux, graphiques, schémas, photos) :}
        \begin{enumerate}
            \item \textbf{Appel dans le texte :} \og (cf. tableau 1)\fg, \og (voir figure 2)\fg.
            \item \textbf{Légende obligatoire :} numéro + titre explicite + source + unités. \emph{Exemple :} \og Tableau 1 : évolution de la pression (en bars) en fonction du temps (en minutes). Source : relevés SCADA.\fg
            \item \textbf{Commentaire et exploitation :} ne jamais laisser une figure \og muette\fg{} : la présenter, la lire et l'interpréter \textbf{juste après son insertion}.
        \end{enumerate}
        \item \textbf{Vocabulaire technique précis :} ne pas écrire \og c'est cassé\fg{} mais \og rupture de l'arbre\fg{} ou \og défaillance du roulement\fg.
    \end{enumerate}
\end{methode}

\begin{exoresolu}[Intégrer et commenter un document iconique dans le corps]
    \textbf{Énoncé :} dans la partie \og 2.1. Diagnostic technique\fg, vous disposez du relevé de température du palier (ci-dessous). Insérez ce document dans le texte, légendez-le correctement et rédigez le commentaire analytique qui l'accompagne.

    \emph{Données :} $t = 0$ min (35 °C), $t = 10$ min (42 °C), $t = 20$ min (58 °C), $t = 30$ min (85 °C), $t = 40$ min (110 °C — alarme), $t = 50$ min (arrêt d'urgence).
    \tcblower
    \textbf{Solution détaillée :}

    \textbf{Extrait du corps du rapport (partie 2.1) :}

    \og L'analyse des paramètres de supervision (SCADA) révèle une dérive thermique progressive du palier côté entraînement avant l'arrêt brutal (voir figure 1).
    \begin{popfigure}
    \begin{center}
        \begin{tikzpicture}
        \begin{axis}[
            width=0.85\linewidth, height=6cm,
            xlabel={Temps (min)}, ylabel={Température (°C)},
            grid=major, grid style={popline},
            xmin=0, xmax=55, ymin=30, ymax=120,
            xtick={0,10,20,30,40,50},
            axis lines=left,
            ]
        \addplot[popcurve, very thick, mark=*, mark size=2pt, popblue] coordinates {
            (0,35) (10,42) (20,58) (30,85) (40,110) (50,110)
        };
        \addplot[poporange, dashed, thick, domain=0:50] {80};
        \node[poporange, anchor=south west] at (axis cs:2,80) {\small Seuil d'alarme (80 °C)};
        \addplot[poppink, dashed, thick, domain=0:50] {105};
        \node[poppink, anchor=south west] at (axis cs:2,105) {\small Seuil d'arrêt (105 °C)};
        \end{axis}
        \end{tikzpicture}
    \end{center}
    \legende{Figure 1 : évolution de la température du palier côté entraînement avant l'arrêt d'urgence. Source : système SCADA, station de Nkolfoulou, 15/03/2026.}
    \end{popfigure}
    La figure 1 met en évidence une \textbf{montée en température de plus en plus rapide} entre $t = 20$ min (58 °C) et $t = 40$ min (110 °C). Le franchissement du seuil d'alarme (80 °C), atteint vers $t = 28$ min, n'a pas déclenché d'intervention humaine immédiate. Le seuil d'arrêt d'urgence (105 °C) est dépassé à $t = 40$ min, provoquant l'arrêt automatique dix minutes plus tard. Cette cinétique thermique confirme une \textbf{défaillance de la lubrification du roulement} (frottement à sec) plutôt qu'une surcharge hydraulique, qui aurait produit une chauffe plus lente et plus régulière.\fg

    \textbf{Points clés respectés :}
    \begin{itemize}
        \item \textbf{Légende complète} (numéro, titre, source, unités sur les axes).
        \item \textbf{Appel du document} dans la phrase d'introduction (\og voir figure 1\fg).
        \item \textbf{Commentaire analytique :} lecture des seuils, interprétation physique (chauffe rapide = frottement à sec), lien avec la cause racine.
    \end{itemize}
\end{exoresolu}

\begin{methode}[Produire et améliorer un rapport complet (processus itératif)]
    \textbf{Objectif :} passer du brouillon à la version finale soumise à la signature (qualité attendue au Probatoire et au Baccalauréat technique).
    \begin{enumerate}
        \item \textbf{Phase 1 — Pré-écriture (collecte et planification) :} rassembler notes, procès-verbaux, photos, données brutes. Valider le plan détaillé (méthode 2).
        \item \textbf{Phase 2 — Rédaction du brouillon :} écrire sans chercher la perfection, en suivant le plan. Insérer des repères du type \og [figure X à insérer]\fg. Utiliser les connecteurs (méthode 3).
        \item \textbf{Phase 3 — Révision macro-structurelle (contenu et structure) :}
        \begin{itemize}
            \item Adéquation entre le plan et la rédaction ? Introduction et conclusion cohérentes ?
            \item Complétude : toutes les pièces jointes annoncées sont-elles présentes ? Le plan est-il respecté ?
            \item Pertinence : supprimer le superflu (sentiments, digressions).
        \end{itemize}
        \item \textbf{Phase 4 — Révision micro-structurelle (langue et forme) :}
        \begin{itemize}
            \item \textbf{Orthographe lexicale et grammaticale} (accords, homophones).
            \item \textbf{Ponctuation} (virgule après les connecteurs, points-virgules dans les énumérations).
            \item \textbf{Typographie :} majuscules aux entités, italique pour les titres de documents, espaces dans les nombres (\num{12000} m\textsuperscript{3}), numérotation décimale alignée.
            \item \textbf{Cohésion :} vérifier les chaînes de référence (antécédents des pronoms clairs).
        \end{itemize}
        \item \textbf{Phase 5 — Mise en page finale et validation :} en-tête et pied de page, sommaire, numérotation des pages, relecture finale, signature, archivage.
    \end{enumerate}
\end{methode}

\begin{exoresolu}[Corriger un brouillon de rapport : exercice de révision]
    \textbf{Énoncé :} voici un extrait du brouillon d'un élève pour la partie \og 3. Bilan\fg. Identifiez cinq erreurs ou faiblesses (contenu, langue, forme) et proposez la version corrigée.

    \og 3. Bilan de la panne
    La panne a duré 6 heures. C'est long. On a perdu beaucoup d'eau, genre 12000 metres cubes. Ca coute cher pour l'entreprise. Le prix de la piece etait 500 000 F mais avec la main d'oeuvre et la perte d'exploitation ca fait plus de 1 million. C'est dommage car on aurait pu l'eviter si on avait graissé avant. Le graissage etait prevu le 20 mars mais la panne est arrivee le 15.\fg
    \tcblower
    \textbf{Solution détaillée :}

    \textbf{Analyse des cinq erreurs et faiblesses :}
    \begin{enumerate}
        \item \textbf{Registre familier et oral :} \og c'est long\fg, \og genre\fg, \og c'est dommage\fg, \og ça coûte cher\fg.
        \item \textbf{Imprécision chiffrée et unités :} \og 12000 metres cubes\fg{} (pas d'espace dans le nombre, unité non symbolisée), \og plus de 1 million\fg{} (vague).
        \item \textbf{Syntaxe et ponctuation faibles :} phrases courtes hachées, absence de connecteurs logiques, accents et apostrophes manquants (\og Ca\fg, \og etait\fg, \og l'eviter\fg).
        \item \textbf{Manque de structure interne :} pas de sous-parties numérotées (3.1., 3.2.), pas de phrase d'introduction de l'idée.
        \item \textbf{Subjectivité et jugement de valeur :} \og c'est dommage\fg{} doit être remplacé par une analyse factuelle (écart entre planning prévu et réalité).
    \end{enumerate}

    \textbf{Version corrigée (version finale) :}
    \begin{quote}
    \textbf{3. BILAN ET ÉVALUATION DES PERTES}

    \textbf{3.1. Durée d'indisponibilité}

    La station est restée à l'arrêt pendant \textbf{six heures} (de 8 h 12 à 14 h 15), soit un quart de la journée d'exploitation.

    \textbf{3.2. Pertes de production}

    Le volume d'eau non produit est estimé à \textbf{\num{12000} m\textsuperscript{3}}, ce qui correspond aux besoins journaliers d'environ \num{60000} abonnés des quartiers Est.

    \textbf{3.3. Impact financier}

    Le coût global de l'incident s'élève à \textbf{\num{1250000} FCFA}, répartis comme suit :
    \begin{itemize}
        \item pièce de rechange (roulement) : \num{480000} FCFA ;
        \item main-d'œuvre et logistique d'urgence : \num{220000} FCFA ;
        \item perte d'exploitation (eau non facturée) : \num{550000} FCFA.
    \end{itemize}

    \textbf{3.4. Analyse de l'écart entre planning et réalité}

    La maintenance préventive (graissage) était programmée au \textbf{20 mars 2026} (échéance mensuelle). La survenue de la panne le \textbf{15 mars} révèle un \textbf{écart de cinq jours} entre la fréquence théorique et la durée de vie réelle du lubrifiant dans ces conditions de charge. Cet écart justifie la révision de la périodicité de graissage.
    \end{quote}
    \textbf{Gains :} registre \textbf{formel et technique}, \textbf{chiffres précis} avec unités, \textbf{structure numérotée}, \textbf{connecteurs logiques} (\og soit\fg, \og ce qui correspond à\fg, \og révèle\fg), \textbf{analyse factuelle} remplaçant le regret. Vérification : $480\,000 + 220\,000 + 550\,000 = 1\,250\,000$ FCFA, total cohérent avec la conclusion du rapport.
\end{exoresolu}

\begin{attention}[Les cinq erreurs qui coûtent des points à l'examen (rapport)]
    \begin{enumerate}
        \item \textbf{Oubli des mentions obligatoires de l'en-tête et du pied de page :} pas de référence, pas de date, pas d'objet clair, pas de destinataire ni d'émetteur identifiés, pas de signature ni de visa $\rightarrow$ \textbf{sanction :} note plafonnée (non-respect de la structure externe attendue).
        \item \textbf{Introduction sans annonce de plan (ou plan non respecté) :} l'introduction pose le sujet mais ne dit pas \emph{comment} il sera traité, ou le corps du rapport ne suit pas l'ordre annoncé $\rightarrow$ \textbf{sanction :} incohérence structurelle majeure (critère \og organisation\fg).
        \item \textbf{Registre inapproprié (familier, subjectif, littéraire) :} usage de \og je pense que\fg, \og c'est nul\fg, \og genre\fg, de métaphores ou d'un style narratif (\og nous sommes arrivés, nous avons vu...\fg) au lieu du style \textbf{procédural et impersonnel} (\og l'équipe a constaté...\fg, \og il a été procédé à...\fg).
        \item \textbf{Documents iconiques \og muets\fg{} ou mal légendés :} insérer un tableau, un graphique ou une photo sans numéro, sans titre, sans source, sans unité, et \textbf{sans commentaire analytique} dans le texte qui suit $\rightarrow$ \textbf{sanction :} perte de points sur \og exploitation des documents\fg{} et \og qualité de l'information\fg.
        \item \textbf{Recommandations floues, non chiffrées ou hors sujet :} \og il faut faire attention\fg, \og il faut améliorer la maintenance\fg. \textbf{Exigence :} des recommandations \textbf{précises, mesurables et datées} : \og remplacer le roulement par la référence X avant le 30/04\fg, \og passer la périodicité de graissage de 30 à 15 jours à compter du 01/04\fg.
    \end{enumerate}
\end{attention}

\newpage

\coursec{Exercices}

\courssub{Je vérifie mes connaissances}

\begin{exercice}[QCM : caractéristiques du rapport]
Parmi les affirmations suivantes, cochez celle qui correspond à la définition du rapport en milieu technique et professionnel.
\begin{enumerate}
    \item Un texte narratif qui raconte une histoire vécue à la première personne.
    \item Un texte fonctionnel, objectif et structuré, rendant compte d'une mission, d'une visite ou d'une enquête.
    \item Un texte argumentatif visant à convaincre un destinataire d'adopter un point de vue.
    \item Un texte purement descriptif et esthétique d'un lieu ou d'un objet.
\end{enumerate}
\end{exercice}

\begin{exercice}[Vrai ou Faux justifié : structure du rapport]
Indiquez si les affirmations suivantes sont vraies (V) ou fausses (F). Justifiez brièvement chaque réponse.
\begin{enumerate}
    \item L'introduction d'un rapport doit obligatoirement contenir le résumé détaillé de toutes les observations effectuées.
    \item La structure externe comprend la page de garde, la table des matières, le corps du texte et les annexes.
    \item Le pronom relatif \emph{dont} s'utilise uniquement lorsque l'antécédent est une personne.
    \item La conclusion d'un rapport technique se termine généralement par des recommandations ou des perspectives.
\end{enumerate}
\end{exercice}

\begin{exercice}[Définitions et outils de liaison]
\begin{enumerate}
    \item Définissez le terme \cle{rapport} en contexte professionnel.
    \item Citez trois \cle{adverbes de liaison} (ou locutions adverbiales) marquant l'addition d'idées.
    \item Citez deux \cle{pronoms relatifs} permettant d'éviter la répétition du nom d'un lieu ou d'un moment.
    \item Quelle est la différence fondamentale entre le \cle{compte rendu} et le \cle{rapport} ?
\end{enumerate}
\end{exercice}

\begin{exercice}[Application directe : choix du pronom relatif]
Complétez les phrases suivantes par le pronom relatif approprié (\emph{qui, que, dont, où, lequel, auxquels}).
\begin{enumerate}
    \item L'atelier \_\_\_\_\_\_\_ nous avons visité se trouve à Yaoundé.
    \item Le chef de mission \_\_\_\_\_\_\_ a dirigé l'enquête est un ingénieur civil.
    \item Voici le rapport \_\_\_\_\_\_\_ vous avez besoin pour la réunion.
    \item Le mois \_\_\_\_\_\_\_ a eu lieu la rentrée était chargé.
    \item Les machines \_\_\_\_\_\_\_ le moteur a été remplacé fonctionnent à nouveau.
\end{enumerate}
\end{exercice}

\courssub{Je m'entraîne}

\begin{exercice}[Rédaction de l'introduction d'un rapport]
Vous avez effectué une visite d'inspection à l'atelier de menuiserie \emph{« Bois Noble »} à Yaoundé, du 10 au 12 octobre 2024, sur instruction du Directeur de l'Enseignement Technique. Rédigez l'introduction complète de ce rapport (contexte, objet, méthode, plan annoncé) en respectant le style impersonnel et objectif. (15 à 20 lignes)
\end{exercice}

\begin{methode}[Rappel : les quatre éléments de l'introduction]
\begin{enumerate}
    \item \textbf{Contexte} : qui a ordonné la mission, quand, pourquoi.
    \item \textbf{Objet} : ce que le rapport examine.
    \item \textbf{Méthode} : comment les données ont été recueillies (observation, entretiens, documents).
    \item \textbf{Annonce du plan} : les grandes parties du corps, sans les développer.
\end{enumerate}
\end{methode}

\begin{exercice}[Exploitation de document statistique : rédaction du corps]
Le tableau ci-dessous présente la fréquentation mensuelle du Centre de Santé Intégré de Bafoussam sur le premier semestre 2024.

\begin{center}
\begin{tabularx}{\linewidth}{|c|X|}\hline
\rowcolor{popblueL}\textbf{Mois} & \textbf{Nombre de patients} \\ \hline
Janvier & 1\,240 \\ \hline
Février & 1\,180 \\ \hline
Mars & 1\,350 \\ \hline
Avril & 1\,520 \\ \hline
Mai & 1\,480 \\ \hline
Juin & 1\,610 \\ \hline
\end{tabularx}
\end{center}

Rédigez le corps du rapport en trois paragraphes distincts :
\begin{enumerate}
    \item Évolution générale de la fréquentation (tendances, chiffres clés).
    \item Analyse comparative (pics, creux, écarts).
    \item Hypothèses explicatives objectives (saisonnalité, campagnes de vaccination, etc.).
\end{enumerate}
Utilisez des outils de liaison (\emph{d'abord, ensuite, par conséquent, en outre, toutefois}) et le style impersonnel.
\end{exercice}

\begin{attention}[Pièges fréquents dans l'exploitation d'un tableau statistique]
\begin{itemize}
    \item Ne recopiez pas tous les chiffres : sélectionnez les valeurs remarquables (minimum, maximum, écart).
    \item Calculez correctement : ici, la hausse globale entre janvier et juin est de $1\,610 - 1\,240 = 370$ patients, soit environ $30\,\%$ ($370/1\,240 \approx 0{,}298$).
    \item Distinguez le creux (février : 1\,180) du pic (juin : 1\,610) ; l'écart maximal est de 430 patients.
    \item Les hypothèses explicatives doivent rester prudentes : employez le conditionnel (\emph{cela pourrait s'expliquer par...}).
\end{itemize}
\end{attention}

\begin{exercice}[Réécriture avec outils de liaison : rapport à trous]
Complétez le texte suivant en choisissant les connecteurs logiques appropriés dans la liste : \emph{d'abord / en outre / par conséquent / toutefois / en définitive / par exemple / ensuite}.

« \_\_\_\_\_\_, l'équipe s'est rendue sur le site du marché central. \_\_\_\_\_\_, elle a procédé au recensement des étals. \_\_\_\_\_\_, elle a interrogé les commerçants sur leurs conditions de travail. \_\_\_\_\_\_, l'absence de registre officiel a ralenti la collecte des données. \_\_\_\_\_\_, 40\,\% des vendeurs n'ont pas de titre d'occupation. \_\_\_\_\_\_, l'insalubrité des abords constitue un risque sanitaire majeur. \_\_\_\_\_\_, la réhabilitation du marché s'impose comme une priorité. »
\end{exercice}

\begin{exercice}[Remise en ordre et justification : rapport de rentrée]
Les paragraphes suivants, extraits d'un rapport sur la rentrée scolaire 2024 au Lycée Technique de Douala-Bassa, sont mélangés.
\begin{enumerate}
    \item Numérotez-les de 1 à 5 pour reconstituer la structure logique (Introduction, Partie 1, Partie 2, Partie 3, Conclusion).
    \item Justifiez l'ordre choisi en identifiant les indices textuels (connecteurs, référents, progression thématique).
\end{enumerate}
\begin{itemize}
    \item[$\square$] \textbf{Paragraphe A :} En définitive, malgré les efforts consentis, le manque de salles de classe et de manuels techniques demeure préoccupant. Nous recommandons la construction d'un bâtiment de 6 salles et la dotation urgente en équipements de TP.
    \item[$\square$] \textbf{Paragraphe B :} D'abord, l'effectif global a atteint 2\,450 élèves, soit une hausse de 8\,\% par rapport à 2023. Les séries industrielles (F4, F2) absorbent 60\,\% de cet effectif.
    \item[$\square$] \textbf{Paragraphe C :} Ce rapport rend compte de la rentrée scolaire 2024 au Lycée Technique de Douala-Bassa, effectuée du 2 au 6 septembre. Il analyse les effectifs, l'encadrement et les infrastructures.
    \item[$\square$] \textbf{Paragraphe D :} Ensuite, le corps enseignant compte 142 professeurs, dont 15 contractuels. Le ratio élèves/professeur s'élève à 17,2, supérieur à la norme de 15.
    \item[$\square$] \textbf{Paragraphe E :} Par ailleurs, sur les 42 salles de classe, 5 sont inutilisables (toitures endommagées). Les ateliers de mécanique et d'électricité manquent de matériel de pointe.
\end{itemize}
\end{exercice}

\courssub{Je me prépare au Bac}

\begin{exercice}[Sujet type Baccalauréat Technique : rapport de visite d'inspection]
\textbf{Situation :} Vous êtes Inspecteur Pédagogique Régional de l'Enseignement Technique pour la région du Centre. Du 15 au 17 janvier 2025, vous avez effectué une visite d'inspection à l'atelier de menuiserie-école du Lycée Technique Industriel de Mfou.

\textbf{Consignes :}
Rédigez le rapport complet de cette mission (environ 25 à 30 lignes) structuré en quatre parties obligatoires :
\begin{enumerate}
    \item \textbf{Introduction :} contexte (ordre de mission n° 012/MINESEC/SG/DRH du 10/01/2025), objet, dates, méthode (observation, entretiens, consultation des registres).
    \item \textbf{Corps (deux parties) :}
        \begin{itemize}
            \item \textbf{Partie 1 :} état des lieux des infrastructures et équipements (surface 120 m², 15 établis, 3 combinées à bois dont 1 en panne, stock de bois insuffisant, ventilation défaillante).
            \item \textbf{Partie 2 :} pédagogie et encadrement (2 professeurs pour 48 élèves de Terminale, programme respecté, sécurité : port des EPI inégal, 2 accidents mineurs en 2024).
        \end{itemize}
    \item \textbf{Conclusion :} bilan synthétique + 3 recommandations concrètes, chiffrées et datées (ex. : réparation de la combinée avant mars 2025, recrutement d'un vacataire, dotation en EPI).
\end{enumerate}
\textbf{Contraintes :} style impersonnel, vocabulaire technique (établi, combinée, EPI, vacataire), connecteurs logiques variés, objectivité totale.
\end{exercice}

\begin{exercice}[Sujet type Probatoire : rapport d'enquête sur un marché]
\textbf{Situation :} En tant que membre de la Commission Municipale d'Hygiène et de Sécurité de la Commune de Bafoussam, vous avez mené une enquête du 5 au 9 février 2025 au Marché Central (Marché A) suite à des plaintes récurrentes des usagers.

\textbf{Données brutes collectées :}
\begin{itemize}
    \item 320 étals recensés (180 alimentaires, 140 non alimentaires).
    \item 45\,\% des étals alimentaires sans certificat médical valide.
    \item 60 poubelles pour 320 étals (ratio d'environ 1 poubelle pour 5,3 étals) ; 70\,\% débordantes à 13 h.
    \item 3 points d'eau potable fonctionnels sur 8 prévus.
    \item 12 cas de troubles gastro-intestinaux signalés au CMS voisin en janvier 2025 (lien suspecté).
    \item Règlement municipal n° 04/2020 non affiché, non appliqué.
\end{itemize}

\textbf{Consignes :}
\begin{enumerate}
    \item Rédigez la page de garde (titre, auteur, commanditaire, date, destinataire).
    \item Rédigez l'introduction (contexte, objet, méthodologie : observation, entretiens avec 50 commerçants et usagers, analyse des registres du CMS).
    \item Rédigez le corps en trois parties thématiques titrées (ex. : \emph{1. Conformité sanitaire des commerçants}, \emph{2. Gestion des déchets et assainissement}, \emph{3. Ressources en eau et prévention des risques}). Intégrez tous les chiffres fournis. Utilisez des tableaux de synthèse si pertinent.
    \item Rédigez la conclusion avec 4 mesures correctives prioritaires (administratives, techniques, pédagogiques, coercitives).
\end{enumerate}
\textbf{Volume :} corps du rapport $\approx$ 30 lignes. Langue soignée, ton administratif.
\end{exercice}

\begin{exercice}[Amélioration d'un rapport d'élève : correction et réécriture]
Voici un extrait de rapport rédigé par un élève de Terminale F4 après un stage en entreprise. Ce texte comporte au moins 10 défauts (subjectivité, absence d'objet clair, fautes d'accord et de grammaire, liaisons absentes ou maladroites, conclusion sans suggestion argumentée, vocabulaire familier, structure confuse, répétitions, absence de chiffres précis, ton émotionnel).

\textbf{Texte à corriger :}
« Mon stage c'était super bien passé à l'entreprise CAMBOIS à Douala du 15 au 30 juin. J'ai trop appris des choses. Le patron il était sympa. Il y a des machines qui marchent pas bien, la scie à ruban elle fait un bruit bizarre et la dégauchisseuse elle est vieille. Les ouvriers ils sont contents mais ils disent que le salaire est bas. Moi j'ai aimé faire des meubles. À la fin je pense que l'entreprise doit acheter des machines neuves et augmenter les salaires. Voilà c'est tout. »

\textbf{Consignes :}
\begin{enumerate}
    \item Listez et numérotez les 10 défauts identifiés en les classant par catégorie : \emph{Structure}, \emph{Langue/Grammaire}, \emph{Style/Ton}, \emph{Contenu}.
    \item Proposez une réécriture professionnelle de ce même rapport (introduction + corps en 2 parties + conclusion avec suggestions), en appliquant toutes les règles du genre (style impersonnel, objectivité, vocabulaire technique, connecteurs, données chiffrées inventées mais plausibles : effectif, machines, durée, etc.).
\end{enumerate}
\end{exercice}



\newpage

\coursec{Corrigés des exercices}

\coursec{Leçon 17 : Produire un rapport}

\begin{corrige}[Exercice 1 — QCM : caractéristiques du rapport]
La bonne réponse est la \textbf{2}. Le rapport est un \cle{texte fonctionnel} : il répond à une commande (ordre de mission), rend compte de faits observés de manière \cle{objective} (style impersonnel) et suit une \cle{structure} codifiée (introduction, corps, conclusion). Les réponses 1, 3 et 4 renvoient à d'autres genres : le récit autobiographique (1), le texte argumentatif (3) et la description littéraire (4).

\textbf{Point de vigilance :} ne pas confondre rapport et compte rendu argumentatif ; le rapport n'a pas pour fonction première de convaincre mais d'informer et de proposer.
\end{corrige}

\begin{corrige}[Exercice 2 — Vrai ou Faux justifié : structure du rapport]
\begin{enumerate}
    \item \textbf{Faux.} L'introduction présente uniquement le contexte, l'objet, la méthode et le plan ; les observations détaillées appartiennent au corps du rapport.
    \item \textbf{Vrai.} La page de garde, la table des matières, le corps du texte et les annexes constituent bien les quatre éléments de la structure externe (ou « appareil ») du rapport.
    \item \textbf{Faux.} \emph{Dont} remplace un complément introduit par \emph{de}, que l'antécédent soit une personne (\emph{l'élève dont je parle}) ou une chose (\emph{la machine dont le moteur est en panne}).
    \item \textbf{Vrai.} La conclusion dresse un bilan synthétique et formule des recommandations ou des perspectives ; c'est précisément ce qui distingue le rapport du simple compte rendu.
\end{enumerate}
\end{corrige}

\begin{corrige}[Exercice 3 — Définitions et outils de liaison]
\begin{enumerate}
    \item Le \cle{rapport} est un texte fonctionnel, objectif et structuré, rédigé à la demande d'une autorité (hiérarchie, administration, commanditaire), qui rend compte d'une mission, d'une visite ou d'une enquête et qui aboutit à des conclusions assorties de recommandations.
    \item Trois marqueurs d'addition : \emph{en outre}, \emph{de plus}, \emph{par ailleurs} (sont aussi acceptés : \emph{également}, \emph{non seulement... mais aussi}).
    \item Deux pronoms relatifs de lieu ou de temps : \emph{où} (\emph{l'atelier où nous sommes allés} ; \emph{le jour où la mission a débuté}) et \emph{lequel} après préposition (\emph{la salle dans laquelle se tient la réunion}).
    \item Le \cle{compte rendu} relate fidèlement des faits ou des paroles sans les commenter ; le \cle{rapport}, lui, analyse et interprète les faits, puis propose des recommandations. Le rapport suppose donc un travail d'expertise que le compte rendu n'exige pas.
\end{enumerate}
\end{corrige}

\begin{corrige}[Exercice 4 — Application directe : choix du pronom relatif]
\begin{enumerate}
    \item L'atelier \emph{que} nous avons visité se trouve à Yaoundé. (\emph{que} = COD de « avons visité » : nous avons visité \emph{l'atelier}.)
    \item Le chef de mission \emph{qui} a dirigé l'enquête est un ingénieur civil. (\emph{qui} = sujet de « a dirigé ».)
    \item Voici le rapport \emph{dont} vous avez besoin pour la réunion. (\emph{dont} = complément introduit par \emph{de} : « avoir besoin \emph{de} ».)
    \item Le mois \emph{où} a eu lieu la rentrée était chargé. (\emph{où} = complément de temps.)
    \item Les machines \emph{dont} le moteur a été remplacé fonctionnent à nouveau. (\emph{dont} = complément du nom « le moteur » : le moteur \emph{des} machines.)
\end{enumerate}
\textbf{Point de vigilance :} \emph{où} ne s'emploie pas seulement pour les lieux ; il convient aussi pour le temps (\emph{le mois où}, \emph{l'année où}).
\end{corrige}

\begin{corrige}[Exercice 5 — Rédaction de l'introduction d'un rapport]
\textbf{Exemple de production attendue :}

Sur instruction du Directeur de l'Enseignement Technique, une mission d'inspection s'est déroulée du 10 au 12 octobre 2024 à l'atelier de menuiserie « Bois Noble », situé à Yaoundé. Le présent rapport a pour objet de rendre compte de cette visite et d'évaluer le fonctionnement de l'atelier, tant sur le plan matériel que sur le plan humain. Les informations recueillies proviennent de trois sources complémentaires : l'observation directe des locaux et des équipements, des entretiens menés avec le gérant et les ouvriers, et la consultation des documents de gestion de l'entreprise. Ce rapport examine d'abord l'état des infrastructures et du matériel, ensuite l'organisation du travail et les conditions du personnel, avant de formuler, en conclusion, un bilan et des recommandations.

\textbf{Critères de réussite (barème indicatif sur 5 points) :}
\begin{itemize}
    \item Contexte (ordre de mission, dates, lieu) : 1 point.
    \item Objet clairement formulé : 1 point.
    \item Méthode d'enquête précisée : 1 point.
    \item Annonce du plan (deux ou trois parties) : 1 point.
    \item Style impersonnel et correction de la langue : 1 point.
\end{itemize}
\textbf{Points de vigilance :} aucune observation détaillée ne doit figurer dans l'introduction ; bannir la première personne (\emph{j'ai visité} $\rightarrow$ \emph{une visite a été effectuée} ou \emph{la mission a porté sur}).
\end{corrige}

\begin{corrige}[Exercice 6 — Exploitation de document statistique : rédaction du corps]
\textbf{Vérification préalable des calculs :}
\begin{itemize}
    \item Hausse globale janvier--juin : $1\,610 - 1\,240 = 370$ patients, soit $\dfrac{370}{1\,240} \approx 0{,}298$, c'est-à-dire environ $30\,\%$ d'augmentation.
    \item Creux : février avec 1\,180 patients ; pic : juin avec 1\,610 patients ; écart maximal : $1\,610 - 1\,180 = 430$ patients.
    \item Recul janvier--février : $1\,240 - 1\,180 = 60$ patients, soit environ $4{,}8\,\%$ de baisse.
\end{itemize}

\textbf{Exemple de corps rédigé :}

\emph{Paragraphe 1 — Évolution générale.} D'abord, il convient de souligner que la fréquentation du Centre de Santé Intégré de Bafoussam connaît une progression sensible au cours du premier semestre 2024 : le nombre de patients passe de 1\,240 en janvier à 1\,610 en juin, soit une hausse de 370 patients, c'est-à-dire une augmentation d'environ 30\,\% en six mois. Cette tendance générale à la hausse traduit une sollicitation croissante des services du centre.

\emph{Paragraphe 2 — Analyse comparative.} Ensuite, l'examen comparatif des données révèle deux moments contrastés. Le mois de février enregistre le niveau le plus bas (1\,180 patients), avec un léger recul de 60 consultations par rapport à janvier. À l'inverse, le mois de juin constitue le pic du semestre (1\,610 patients). L'écart entre le creux et le pic atteint ainsi 430 patients. Toutefois, après la forte progression observée entre mars (1\,350) et avril (1\,520), on note un ralentissement en mai (1\,480), avant une reprise en juin.

\emph{Paragraphe 3 — Hypothèses explicatives.} Enfin, plusieurs hypothèses pourraient expliquer cette évolution. La baisse de février pourrait correspondre à un mois plus court et à une moindre circulation des maladies saisonnières. En outre, la hausse du second trimestre pourrait s'expliquer par le retour des pluies, favorable aux affections respiratoires et au paludisme, ou par d'éventuelles campagnes de vaccination ayant attiré de nouveaux usagers. Ces hypothèses devraient être confirmées par l'analyse des registres médicaux.

\textbf{Points de vigilance :} ne pas recopier tous les chiffres ; employer le conditionnel pour les hypothèses ; varier les connecteurs (\emph{d'abord, ensuite, toutefois, en outre, enfin}) ; maintenir le style impersonnel.
\end{corrige}

\begin{corrige}[Exercice 7 — Réécriture avec outils de liaison : rapport à trous]
« \emph{D'abord}, l'équipe s'est rendue sur le site du marché central. \emph{Ensuite}, elle a procédé au recensement des étals. \emph{En outre}, elle a interrogé les commerçants sur leurs conditions de travail. \emph{Toutefois}, l'absence de registre officiel a ralenti la collecte des données. \emph{Par exemple}, 40\,\% des vendeurs n'ont pas de titre d'occupation. \emph{Par conséquent}, l'insalubrité des abords constitue un risque sanitaire majeur. \emph{En définitive}, la réhabilitation du marché s'impose comme une priorité. »

\textbf{Justification des choix :}
\begin{itemize}
    \item \emph{D'abord / ensuite / en outre} ordonnent les étapes chronologiques de la mission (addition et succession).
    \item \emph{Toutefois} introduit une opposition : la difficulté rencontrée.
    \item \emph{Par exemple} illustre concrètement cette difficulté.
    \item \emph{Par conséquent} marque la conséquence logique des constats.
    \item \emph{En définitive} signale la conclusion du raisonnement.
\end{itemize}
\end{corrige}

\begin{corrige}[Exercice 8 — Remise en ordre et justification : rapport de rentrée]
Ordre correct : \textbf{C (1) -- B (2) -- D (3) -- E (4) -- A (5)}.

\textbf{Justification :}
\begin{itemize}
    \item \textbf{C} ouvre le rapport : formule typique d'introduction (« Ce rapport rend compte de... »), dates de la mission et annonce du plan en trois points (effectifs, encadrement, infrastructures).
    \item \textbf{B} : le connecteur \emph{d'abord} marque la première partie, consacrée aux effectifs, premier point du plan annoncé ; la hausse de 8\,\% illustre l'usage des données chiffrées.
    \item \textbf{D} : \emph{ensuite} enchaîne sur le deuxième point annoncé (l'encadrement) ; le ratio élèves/professeur ($2\,450 / 142 \approx 17{,}2$) prolonge logiquement les données d'effectif du paragraphe B.
    \item \textbf{E} : \emph{par ailleurs} introduit le troisième point du plan (les infrastructures et les équipements).
    \item \textbf{A} : \emph{en définitive} et la formulation de recommandations (construction d'un bâtiment, dotation en équipements) signalent sans ambiguïté la conclusion.
\end{itemize}
\textbf{Point de vigilance :} la cohérence entre le plan annoncé en introduction et l'ordre des parties est l'indice textuel décisif ; les connecteurs d'ouverture de paragraphe confirment cet ordre.
\end{corrige}

\begin{corrige}[Exercice 9 — Sujet type Baccalauréat Technique : rapport de visite d'inspection]
\textbf{Exemple de production attendue :}

\emph{Introduction.} En exécution de l'ordre de mission n° 012/MINESEC/SG/DRH du 10 janvier 2025, une visite d'inspection a été effectuée du 15 au 17 janvier 2025 à l'atelier de menuiserie-école du Lycée Technique Industriel de Mfou. Le présent rapport a pour objet d'évaluer l'état des infrastructures et des équipements, ainsi que les conditions d'encadrement pédagogique de l'atelier. Il s'appuie sur l'observation directe des lieux, sur des entretiens avec le personnel enseignant et sur la consultation des registres de l'établissement. Il examine d'abord l'état des lieux matériel, ensuite la pédagogie et l'encadrement, avant de formuler des recommandations.

\emph{Partie 1 — État des lieux des infrastructures et équipements.} D'abord, l'atelier occupe une surface de 120 m² et dispose de 15 établis en état d'usage. Toutefois, le parc machines présente des insuffisances : sur les trois combinées à bois, une est en panne, ce qui réduit d'un tiers la capacité d'usinage. En outre, le stock de bois est insuffisant pour assurer les travaux pratiques du trimestre. Enfin, la ventilation de l'atelier est défaillante, ce qui expose les usagers à la poussière de bois.

\emph{Partie 2 — Pédagogie et encadrement.} Ensuite, l'encadrement est assuré par deux professeurs pour 48 élèves de Terminale, soit un ratio de 24 élèves par enseignant ($48 / 2 = 24$), nettement supérieur aux normes des travaux pratiques. Il convient néanmoins de noter que le programme officiel est respecté. En revanche, le port des équipements de protection individuelle (EPI) demeure inégal : deux accidents mineurs ont été enregistrés en 2024.

\emph{Conclusion.} En définitive, l'atelier de Mfou assure un enseignement conforme au programme, mais souffre d'un équipement dégradé et d'un encadrement insuffisant. Il est recommandé : premièrement, la réparation de la combinée à bois en panne avant mars 2025 ; deuxièmement, le recrutement d'un vacataire dès le deuxième trimestre 2025 ; troisièmement, la dotation urgente de l'atelier en EPI et la remise en état du système de ventilation.

\textbf{Barème indicatif (sur 20) :}
\begin{itemize}
    \item Introduction complète (contexte, objet, méthode, plan) : 4 points.
    \item Partie 1 : toutes les données exploitées, style impersonnel : 4 points.
    \item Partie 2 : données exploitées, ratio calculé, objectivité : 4 points.
    \item Conclusion : bilan + 3 recommandations chiffrées et datées : 4 points.
    \item Langue, connecteurs, vocabulaire technique (établi, combinée, EPI, vacataire) : 4 points.
\end{itemize}
\textbf{Points de vigilance :} citer l'ordre de mission avec sa référence exacte ; ne pas oublier le calcul du ratio ; chaque recommandation doit découler d'un constat du corps et comporter une échéance.
\end{corrige}

\begin{corrige}[Exercice 10 — Sujet type Probatoire : rapport d'enquête sur un marché]
\textbf{Vérification préalable des calculs :}
\begin{itemize}
    \item Étals alimentaires sans certificat : $180 \times 0{,}45 = 81$ étals.
    \item Ratio poubelles : $320 / 60 \approx 5{,}3$, soit environ 1 poubelle pour 5,3 étals ; poubelles débordantes : $60 \times 0{,}70 = 42$.
    \item Points d'eau défaillants : $8 - 3 = 5$ sur 8, soit $62{,}5\,\%$ hors service.
\end{itemize}

\textbf{1. Page de garde (éléments attendus) :}
\begin{itemize}
    \item Titre : \emph{Rapport d'enquête d'hygiène et de sécurité au Marché Central (Marché A) de Bafoussam}.
    \item Auteur : la Commission Municipale d'Hygiène et de Sécurité.
    \item Commanditaire et destinataire : le Maire de la Commune de Bafoussam.
    \item Date : février 2025 (période d'enquête : du 5 au 9 février 2025).
\end{itemize}

\textbf{2. Introduction (modèle) :} Suite aux plaintes récurrentes des usagers, le Maire de la Commune de Bafoussam a chargé la Commission Municipale d'Hygiène et de Sécurité d'une enquête au Marché Central (Marché A), menée du 5 au 9 février 2025. Le présent rapport rend compte de cette mission. Il s'appuie sur l'observation directe des lieux, sur des entretiens avec 50 commerçants et usagers, et sur l'analyse des registres du Centre Médico-Social voisin. Il examine successivement la conformité sanitaire des commerçants, la gestion des déchets et l'assainissement, puis les ressources en eau et la prévention des risques.

\textbf{3. Corps (attendus par partie) :}
\begin{itemize}
    \item \emph{Conformité sanitaire :} 320 étals recensés (180 alimentaires, 140 non alimentaires) ; 45\,\% des étals alimentaires, soit 81 étals, sans certificat médical valide ; le règlement municipal n° 04/2020 n'est ni affiché ni appliqué.
    \item \emph{Déchets et assainissement :} 60 poubelles pour 320 étals (environ 1 pour 5,3 étals) ; 70\,\% d'entre elles, soit 42 poubelles, débordantes à 13 h.
    \item \emph{Eau et prévention :} 3 points d'eau potable fonctionnels sur 8 prévus (5 hors service, soit 62,5\,\%) ; 12 cas de troubles gastro-intestinaux signalés au CMS en janvier 2025, avec un lien suspecté mais non confirmé (conditionnel obligatoire).
\end{itemize}

\textbf{4. Conclusion (modèle) :} En définitive, l'enquête révèle une situation sanitaire préoccupante au Marché A. Quatre mesures correctives sont proposées : premièrement, sur le plan administratif, l'affichage et l'application du règlement n° 04/2020 ; deuxièmement, sur le plan technique, l'installation de 40 poubelles supplémentaires et la réparation des 5 points d'eau défaillants avant avril 2025 ; troisièmement, sur le plan pédagogique, une campagne de sensibilisation des 81 commerçants non conformes ; quatrièmement, sur le plan coercitif, des sanctions progressives contre les récalcitrants après un délai de trois mois.

\textbf{Barème indicatif (sur 20) :} page de garde complète : 2 points ; introduction (contexte, objet, méthode, plan) : 4 points ; corps en trois parties titrées intégrant tous les chiffres : 8 points ; conclusion avec 4 mesures différenciées : 4 points ; langue et ton administratif : 2 points.

\textbf{Points de vigilance :} le lien entre les 12 cas médicaux et le marché doit rester au conditionnel (\emph{lien suspecté}) ; les quatre mesures doivent être de natures distinctes (administrative, technique, pédagogique, coercitive) ; aucun chiffre fourni ne doit être omis.
\end{corrige}

\begin{corrige}[Exercice 11 — Amélioration d'un rapport d'élève : correction et réécriture]
\textbf{1. Grille des défauts (10 défauts classés) :}

\emph{Structure :}
\begin{enumerate}
    \item Absence d'introduction : ni contexte, ni objet, ni méthode.
    \item Absence de plan : les constats sont mélangés sans parties titrées.
    \item Conclusion expéditive (« Voilà c'est tout ») sans bilan ni recommandations argumentées.
\end{enumerate}
\emph{Langue/Grammaire :}
\begin{enumerate}
    \item[4.] Construction fautive : « Mon stage c'était super bien passé » (mélange de « s'est passé » et « c'était »).
    \item[5.] Négation incomplète : « des machines qui marchent pas bien » (absence de \emph{ne}).
    \item[6.] Dislocations familières répétées : « Le patron il était... », « la scie à ruban elle fait... », « Les ouvriers ils sont... ».
\end{enumerate}
\emph{Style/Ton :}
\begin{enumerate}
    \item[7.] Subjectivité et première personne : « J'ai trop appris », « Moi j'ai aimé », « je pense ».
    \item[8.] Vocabulaire familier et vague : « super », « sympa », « trop appris des choses », « un bruit bizarre », « vieille ».
\end{enumerate}
\emph{Contenu :}
\begin{enumerate}
    \item[9.] Aucune donnée chiffrée précise (effectif, nombre de machines, année des dates).
    \item[10.] Suggestions non justifiées : « acheter des machines neuves et augmenter les salaires » sans lien avec des constats mesurés ni échéances.
\end{enumerate}

\textbf{2. Exemple de réécriture professionnelle :}

\emph{Introduction.} Sur instruction du Proviseur du Lycée Technique de Douala-Bassa, un stage d'observation et de pratique s'est déroulé du 15 au 30 juin 2024 au sein de l'entreprise CAMBOIS, spécialisée dans la fabrication de meubles. Le présent rapport rend compte du fonctionnement de l'atelier et des conditions de travail. Il s'appuie sur l'observation directe des postes de travail et sur des entretiens avec le personnel. Il examine d'abord l'état du parc machines, ensuite les conditions de travail, avant de formuler des suggestions.

\emph{Partie 1 — État du parc machines.} L'atelier compte 8 machines, dont une scie à ruban et une dégauchisseuse. D'abord, la scie à ruban émet un bruit anormal qui traduit une usure probable de la lame ou des roulements. Ensuite, la dégauchisseuse, mise en service depuis plus de quinze ans, présente des signes de vétusté qui ralentissent la production.

\emph{Partie 2 — Conditions de travail.} L'entreprise emploie 12 ouvriers. Si la motivation du personnel est réelle, les entretiens font état d'une insatisfaction liée au niveau des rémunérations. En outre, le port des équipements de protection individuelle n'est pas systématique.

\emph{Conclusion.} En définitive, le stage a permis de constater un savoir-faire réel mais un parc machines vieillissant. Il est suggéré, d'une part, la révision de la scie à ruban avant septembre 2024, d'autre part, le remplacement progressif de la dégauchisseuse, enfin l'ouverture d'une négociation salariale avec le personnel.

\textbf{Points de vigilance :} chaque suggestion de la conclusion doit correspondre à un constat du corps ; les données inventées doivent rester plausibles et cohérentes (dates, effectifs) ; la première personne doit disparaître totalement de la version réécrite.
\end{corrige}

\newpage

\coursec{Fiche bilan}

\begin{popfigure}
\begin{center}\tcbox[colback=popblue,colframe=popblue,coltext=white,arc=9pt,boxsep=4pt,left=6pt,right=6pt]{\bfseries\small Le Rapport}\end{center}
\vspace{-2pt}
\begin{tcbraster}[raster columns=3,raster equal height=rows,raster column skip=6pt,raster row skip=6pt,enhanced,blankest]
\begin{tcolorbox}[enhanced,colback=popblue,colframe=popblue,coltext=white,boxrule=0.9pt,arc=3pt,left=4pt,right=4pt,top=3pt,bottom=3pt,boxsep=1pt,halign=left]\bfseries\scriptsize Nature Texte fonctionnel Informer / Décider\end{tcolorbox}
\begin{tcolorbox}[enhanced,colback=poporange,colframe=poporange,coltext=white,boxrule=0.9pt,arc=3pt,left=4pt,right=4pt,top=3pt,bottom=3pt,boxsep=1pt,halign=left]\bfseries\scriptsize Structure Externe Garde, Sommaire, Annexes\end{tcolorbox}
\begin{tcolorbox}[enhanced,colback=popgreen,colframe=popgreen,coltext=white,boxrule=0.9pt,arc=3pt,left=4pt,right=4pt,top=3pt,bottom=3pt,boxsep=1pt,halign=left]\bfseries\scriptsize Structure Interne Introduction, Corps, Conclusion\end{tcolorbox}
\begin{tcolorbox}[enhanced,colback=poppurple,colframe=poppurple,coltext=white,boxrule=0.9pt,arc=3pt,left=4pt,right=4pt,top=3pt,bottom=3pt,boxsep=1pt,halign=left]\bfseries\scriptsize Langage Outils de liaison Adverbes, Relatifs, Locutions\end{tcolorbox}
\begin{tcolorbox}[enhanced,colback=poppink,colframe=poppink,coltext=white,boxrule=0.9pt,arc=3pt,left=4pt,right=4pt,top=3pt,bottom=3pt,boxsep=1pt,halign=left]\bfseries\scriptsize Intro /\& Conclu Contexte, Objet, Plan Bilan, Perspectives\end{tcolorbox}
\begin{tcolorbox}[enhanced,colback=popgold,colframe=popgold,coltext=white,boxrule=0.9pt,arc=3pt,left=4pt,right=4pt,top=3pt,bottom=3pt,boxsep=1pt,halign=left]\bfseries\scriptsize Corps du rapport Faits, Analyses, Documents iconiques\end{tcolorbox}
\begin{tcolorbox}[enhanced,colback=popdark,colframe=popdark,coltext=white,boxrule=0.9pt,arc=3pt,left=4pt,right=4pt,top=3pt,bottom=3pt,boxsep=1pt,halign=left]\bfseries\scriptsize Processus Planifier, Rédiger, Relire, Améliorer\end{tcolorbox}
\end{tcbraster}

\legende{Carte mentale : les 7 piliers de la rédaction d'un rapport en Terminale technique.}
\end{popfigure}

\begin{bilan}

\courssub{Définitions clés}

\begin{definition}[Rapport]
Texte \cle{fonctionnel} à visée professionnelle : il rend compte d'une activité, d'une enquête ou d'un stage à un destinataire précis (hiérarchie, client, jury) pour \cle{informer} et/ou \cle{aider à la décision}. Style objectif, impersonnel, structuré.
\end{definition}

\begin{definition}[Structure externe]
Éléments matériels encadrant le texte : \textbf{page de garde} (titre, auteur, destinataire, date), \textbf{sommaire}, \textbf{liste des abréviations/sigles}, \textbf{annexes} (documents bruts, questionnaires).
\end{definition}

\begin{definition}[Structure interne]
Organisation logique du contenu en trois parties obligatoires :
\begin{enumerate}
\item \textbf{Introduction} : contexte, objet, délimitation, plan annoncé.
\item \textbf{Corps (Développement)} : parties/chapitres équilibrés (faits $\rightarrow$ analyses $\rightarrow$ interprétations).
\item \textbf{Conclusion} : bilan synthétique, réponses à l'objet, perspectives/recommandations.
\end{enumerate}
\end{definition}

\courssub{Outils de liaison (Morphosyntaxe) — À maîtriser pour la cohérence}

\begin{center}
\begin{tabularx}{\linewidth}{|l|X|X|}\hline
\rowcolor{popblueL}
\textbf{Catégorie} & \textbf{Exemples} & \textbf{Fonction principale} \\ \hline
\cle{Adverbes de liaison} & cependant, donc, ensuite, par ailleurs, en effet & Relier deux phrases/segments (relation logique) \\ \hline
\cle{Locutions adverbiales} & en outre, par conséquent, pour cette raison, d'une part... d'autre part & Marquer transition, cause, conséquence, énumération \\ \hline
\cle{Pronoms relatifs} & qui, que, dont, où, lequel, auquel & Relier deux propositions (antécédent + relative) \\ \hline
\cle{Conjonctions de sub.} & comme, puisque, afin que, bien que, si & Introduire une proposition subordonnée (cause, but, concessive) \\ \hline
\end{tabularx}
\end{center}

\courssub{Méthodes express}

\begin{methode}[Rédiger l'Introduction (entrée en 4 temps)]
\begin{enumerate}
\item \textbf{Accroche / Contexte} : Situation générale, commande, actualité (1--2 phrases).
\item \textbf{Objet du rapport} : « Le présent rapport a pour objet de... » (verbe d'action : analyser, proposer, rendre compte).
\item \textbf{Délimitation / Méthodologie} : Période, terrain, outils, population concernée.
\item \textbf{Plan annoncé} : « Nous présenterons d'abord..., puis nous analyserons..., enfin nous proposerons... »
\end{enumerate}
\end{methode}

\begin{methode}[Rédiger le Corps (règle des 3 A)]
\begin{enumerate}
\item \textbf{Annoncer} : Titre de partie clair (verbe nominalisé ou phrase nominale).
\item \textbf{Apporter les faits} : Données brutes, observations, citations, \cle{documents iconiques} (tableaux, graphiques, photos) \textbf{intégrés et commentés} (légende + lecture guidée).
\item \textbf{Analyser / Argumenter} : Interpréter les faits, croiser sources, utiliser connecteurs logiques (cause/conséquence, opposition, concession).
\end{enumerate}
\end{methode}

\begin{methode}[Exploiter un document iconique]
\begin{enumerate}
\item \textbf{Présenter} : « Le tableau 1 ci-dessous montre... » (source, date, nature).
\item \textbf{Décrire} : Tendances majeures, valeurs extrêmes, écarts (lecture globale).
\item \textbf{Interpréter} : Lier aux objectifs du rapport, expliquer \emph{pourquoi} (causes, conséquences).
\end{enumerate}
\end{methode}

\begin{methode}[Rédiger la Conclusion (sortie en 3 temps)]
\begin{enumerate}
\item \textbf{Bilan / Synthèse} : Reprendre les réponses principales à l'objet (sans ajouter d'info nouvelle).
\item \textbf{Portée / Limites} : Fiabilité, difficultés rencontrées, marge d'erreur.
\item \textbf{Perspectives / Recommandations} : Actions concrètes, suites à donner, ouverture.
\end{enumerate}
\end{methode}

\courssub{Check-list structurelle (Externe + Interne)}
\begin{itemize}
\item Page de garde complète (titre précis, commande, auteur, date).
\item Sommaire avec numérotation cohérente (1, 1.1, 1.2 / 2, 2.1...).
\item Introduction : Contexte $\rightarrow$ Objet $\rightarrow$ Démarche $\rightarrow$ Plan.
\item Corps : Parties équilibrées, titres parlants, transition entre parties.
\item Documents iconiques : Numérotés, légendés, source cités, \textbf{exploités} dans le texte.
\item Conclusion : Synthèse $\rightarrow$ Limites $\rightarrow$ Recommandations.
\item Annexes : Questionnaire, raw data, docs longs référencés dans le corps.
\item Relecture : Cohérence, impersonnel (passif, « nous » institutionnel), connecteurs, orthographe.
\end{itemize}

\end{bilan}

\begin{aretenir}[Checklist avant le Bac — 10 points de contrôle]
$\square$ 1. Identifier le \textbf{type de rapport} (stage, activité, enquête, technique) et le \textbf{destinataire}.\par
$\square$ 2. Respecter la \textbf{structure externe} : page de garde normée, sommaire automatique, annexes référencées.\par
$\square$ 3. Construire une \textbf{introduction complète} : accroche, objet clair, délimitation, plan annoncé explicite.\par
$\square$ 4. Organiser le \textbf{corps} en parties logiques (2 à 4) avec titres informatifs et transitions fluides.\par
$\square$ 5. Appliquer la règle \textbf{Faits $\rightarrow$ Analyse} : jamais de faits bruts sans interprétation.\par
$\square$ 6. Varier les \textbf{outils de liaison} (adverbes, locutions, relatifs, conjonctions) pour la cohérence textuelle.\par
$\square$ 7. Intégrer au moins un \textbf{document iconique} (tableau, graphique, schéma) \textbf{commenté} (pas juste inséré).\par
$\square$ 8. Rédiger une \textbf{conclusion en 3 temps} : synthèse répondant à l'objet, limites/honnêteté, recommandations concrètes.\par
$\square$ 9. Adopter le \textbf{style administratif} : impersonnel, précis, objectif, vocabulaire technique maîtrisé.\par
$\square$ 10. \textbf{Relire} : cohérence globale, concordance des temps, orthographe lexicale/grammaticale, mise en page soignée.
\end{aretenir}

\findecours

\end{document}
